\documentclass[11pt,a4 paper]{article}
\setlength{\textwidth}{6.1in}
\setlength{\textheight}{8.3in}
\setlength{\hoffset}{-0.4in}
\setlength{\voffset}{-0.2in}
\newtheorem{theo}{Theorem}[section]
\newtheorem{cor}{Corollary}[section]
\newtheorem{rem}{Remark}[section]
\newtheorem{defi}{Definition}[section]
\newtheorem{lemma}{Lemma}[section]
\newtheorem{for}{Formula}[section]
\newtheorem{prop}{Proposition}[section]
\newtheorem{ex}{Example}[section]
\newtheorem{conj}{Conjecture}
\newtheorem{op}{Open Problem}
\newtheorem{question}{Question}
\newtheorem{algo}{Algorithm}
\newtheorem{con}{Construction}
\newtheorem{idea}{Research Idea}
\newcommand{\proof}{\noindent{\em Proof.}\quad}


%\usepackage[utf8]{inputenc}
\usepackage{multirow}
\usepackage{hhline}
%\usepackage{fourier}
\usepackage{array}
\usepackage{makecell}
%\usepackage{enumitem}

\renewcommand\theadalign{cb}
\renewcommand\theadfont{\bfseries}
\renewcommand\theadgape{\Gape[4pt]}
\renewcommand\cellgape{\Gape[4pt]}

%\def\bproof{\textbf{Proof}: }
%\def\eproof{\hfill$\Box$}

\newtheorem{com}{Comment}
\usepackage{float}
\usepackage{tocloft}
\usepackage{amsfonts,amsmath,amssymb}
\usepackage{multirow, verbatim}
\usepackage{shadow,enumerate}
\usepackage{makeidx}  % allows index generation
\usepackage{graphicx}
\usepackage{color}
\def\Z{{\mathbb Z}}

\usepackage{mathrsfs}

\renewcommand{\theenumi}{\alph{enumi}}
\newcommand{\FB}{{\mathbb F}_{2}}
\def\whitebox{{\hbox{\hskip 1pt
        \vrule height 6pt depth 1.5pt
        \lower 1.5pt\vbox to 7.5pt{\hrule width
                  3.2pt\vfill\hrule width 3.2pt}%
        \vrule height 6pt depth 1.5pt
        \hskip 1pt } }}
\def\qed{\ifhmode\allowbreak\else\nobreak\fi\hfill\quad\nobreak\whitebox\medbreak}

\newtheorem{co}{Comment}

\begin{document}
\title{Designing plateaued Boolean functions in spectral domain  \\ and their classification}

\author{
S. Hod\v zi\'c \footnote {University of Primorska, FAMNIT, Koper,
Slovenia, e-mail: samir.hodzic@famnit.upr.si}
\and E.~Pasalic\footnote{
University of Primorska, FAMNIT \& IAM, Koper, Slovenia, e-mail: enes.pasalic6@gmail.com}
%\and S. Hod\v zi\'c \footnote {University of Primorska, FAMNIT, Koper,
%Slovenia, e-mail: samir.hodzic@famnit.upr.si}
%\and E.~Pasalic\footnote{
%University of Primorska, FAMNIT \& IAM, Koper, Slovenia, e-mail: enes.pasalic6@gmail.com}
\and Y. Wei \footnote{
%Guangxi Experiment Center of Information Science,
 Guilin University of Electronic Technology, Guilin,
P.R. China,
 e-mail: walker$_{-}$wyz@guet.edu.}
\and F. Zhang \footnote{School of Computer Science and Technology,
China University of
 Mining and Technology, Xuzhou, Jiangsu 221116,  P.R. China, e-mail:
zhfl203@cumt.edu.cn.}
}

\date{}
\maketitle


\begin{abstract}
The design of plateaued functions over $GF(2)^n$, also known as 3-valued Walsh spectra functions (taking the values from the set $\{0, \pm 2^{\lceil \frac{n+s}{2} \rceil}\}$), has been commonly approached by specifying a suitable algebraic normal form  which then induces this particular Walsh spectral characterization.
% takes values in the set $\{0, \pm 2^{\lceil \frac{n+s}{2} \rceil}\}$. Such a function $f$ defined on $GF(2)^n$ is then called an $s$-plateaued function.
In this article, we  consider the reversed design method which specifies these functions in the spectral domain by specifying a suitable  allocation of the nonzero spectral values and their signs. We analyze the properties of {\em trivial and nontrivial} plateaued functions (as affine inequivalent distinct subclasses), which are distinguished by their  Walsh support $S_f$ (the subset of $GF(2)^n$ having the nonzero spectral values) in terms of whether it is an affine subspace or not. The former class  exactly corresponds to partially bent functions and admits linear structures, whereas the latter class may contain functions without linear structures. A simple sufficient condition on $S_f$, which ensures the nonexistence of linear structures, is derived  and some generic design methods of nontrivial plateaued functions without linear structures are given.
%affine inequivalent to those plateaued functions having $S_f$ which is not an affine subspace. These nontrivial plateaued functions are much more interesting objects which can provide functions that do not admit linear structures which is confirmed though some generic design methods.
The extended affine equivalence of plateaued functions is also addressed using the concept of {\em dual} of plateaued functions.
Furthermore, we solve the problem of specifying disjoint spectra (non)trivial plateaued functions of maximal cardinality whose concatenation can be used to construct bent functions in a generic manner. This approach may lead to new classes of bent functions due to large variety of possibilities to select underlying duals that define these disjoint spectra plateaued functions.
%introduced in \cite{Secondary}, we firstly consider the concept of the EA-equivalence in terms of corresponding Walsh supports and duals. We show that two plateaued functions, regardless of whether their Walsh supports affine or non-affine subspaces, can be inequivalent if their duals do not satisfy certain conditions. Thus, a nontrivial characterization of plateaued functions is given which is the first step towards better understanding of these objects. In addition, we provide generic constructions of plateaued functions (of any amplitude) with non-affine Walsh supports (which correspond to non-partially bent functions). In this context, the construction of the maximally possible number of disjoint spectra functions with non-affine Walsh supports is addressed.
An additional method of specifying affine inequivalent plateaued functions, obtained by applying a nonlinear transform to their input domain, is also given.
%Finally, we introduce the concept of {\em extendable} set of disjoint spectra plateaued functions referring to the possibility to extend this set with additional disjoint spectra functions so that maximal cardinality is attained.\\\\
%\textbf{Samir: This blue part is not good. It is a trivial fact that affine and non-affine subspaces can not be related by an invertible matrix, and thus we can not say "we show..." (since there is nothing to prove about that). Therefore, I propose the following sentence: "By calling (non)trivial palteaued functions those with (non)affine Walsh support, we completely describe affine inequivalent plateaued functions in terms of their Walsh supports and duals." Even this sentence is not the best one, but we need to emphasize our main result related to equivalency -- that is Theorem 3.2, which describes equivalency in general!!! More than half of the abstract speaks only about constructions etc (and less about equivalency), but equivalency is the first and also quite important part of this work! (We need to say more about it -- to justify the word "classification" in the title of the paper...)}
\newline \newline
\noindent
\textbf{Keywords:} Plateaued functions, Extended affine equivalence, Disjoint spectra plateaued functions, Extendable and nontrivial plateaued functions.

\end{abstract}

\section{Introduction}

%Construction methods of cryptographically significant Boolean functions which possess some of important properties (such as high (maximal) non-linearity/degree, balanceness, correlation/algebraic immunity, resiliency), are of a great importance in many modern cryptographic systems.  A lot of research has been devoted towards efficient design methods  of highly/maximal non-linear functions.


A class of Boolean functions on $GF(2)^n$ characterized by the property that their Walsh spectra is three-valued (more precisely taking values in $\{0,\pm 2^{\frac{n+s}{2}}\}$ for a positive integer $s<n$) is called  $s$-\emph{plateaued} functions  \cite{Zheng}.
% and  later studied in \cite{CarletAPN,CarletQ,Zheng2,Wang,Jong,SihemPlate}.
The notion of plateaued functions as defined here does not include bent and linear functions (which only have  two different values in their Walsh spectrum), though sometimes in the literature these families are also included. This class of functions has a wide range of applications such as  cryptography, in the design of sequences for communications, and related combinatorics and designs.
For instance, vectorial plateaued functions were efficiently used to provide large sets of orthogonal binary sequences suitable for CDMA (Code Division Multiple Access) applications which resulted in a significant improvement of the number of users per cell in regular hexagonal networks \cite{WeiGuoCDMA}.
The design methods of plateaued functions have been addressed in several works in the past twenty years \cite{CarletAPN,CarletQ,Feng2,SihemPOP2014,THOMCusick2016,SihemM2017AMC,Fengrong2018IT,Zheng2}, but none of these construction methods is generic and furthermore a simple characterization of these functions in the Walsh domain has not been utilized as a construction rationale.  In \cite{CarletAPN}, Carlet provided a survey on the construction methods of Boolean plateaued functions along with  certain characterizations of (vectorial) plateaued functions which extends the initial attempt towards their characterization in \cite{CarletQ}. Nevertheless, these  characterizations of plateaued functions, given in terms of their second order derivatives, autocorrelation properties or power moments of the Walsh transform in \cite{CarletAPN}, are not efficient as generic construction methods.
%many primary constructions () and secondary constructions () were provided.

%The initial attempt towards some efficient  construction methods and  characterization  of plateaued functions was made by Carlet and Prouff \cite{CarletQ}.
%In \cite{CarletAPN}, Carlet provided a survey on the construction methods of Boolean plateaued functions along with  certain characterizations of plateaued (vectorial) Boolean functions. Initially, the work on characterization of plateaued functions was conducted Carlet and Prouff \cite{CarletQ}.
% Mesnager  introduced in \cite{CarletAPN}.
%2014 new characterizations of $p$-ary plateaued functions by the constant of the ratio of two
%consecutive Walsh power moments of even order \cite{}.
In addition,  some  methods  for specifying  plateaued functions in the algebraic normal form domain was considered in \cite{THOMCusick2016} and \cite{Logatchev}.   In \cite{SihemM2017AMC}, a secondary construction of plateaued functions was presented based on the use of three suitable  bent functions.
% in $n$ variables can  satisfy some conditions, a plateaued function in $n$ variables will be obtained )
The importance of plateaued functions also relates to construction of highly nonlinear resilient functions and for this purpose  the so-called disjoint spectra plateaued functions  are required
\cite{Zhang2009}. Certain methods, mainly employing the (generalized) indirect sum method of Carlet \cite{SecEnf},  for generating the sets of disjoint spectra plateaued functions without linear structures were investigated  in \cite{Feng2,Fengrong2018IT}.
%In \cite{Feng2}, a method for constructing  quadruples of disjoint spectra  plateaued functions without  linear structures was provided.
%
%Recently, by using the indirect sum,  an efficient method for constructing a large set of disjoint spectra plateaued functions without linear structures was presented \cite{Fengrong2018IT}. In addition, a much larger set of disjoint spectra plateaued functions can be constructed by using the generalized indirect sum \cite{Fengrong2018IT}.
We briefly mention that the notion of plateaued functions has been generalized over different mathematical structures, the interested reader is referred to   \cite{AycaWilffried,SihemPOP2014,Jong,SihemPlate,SihemM2017IT}.
%then $p$-ary plateaued functions on $GF(p)^n$ (where $p$ is on odd prime) were introduced in \cite{SihemPlate} and similarly to the notion of generalized Boolean (bent) functions generalized plateaued functions were considered in \cite{SihemM2017IT}.
%However, this article mainly deals with the classical  Boolean plateaued functions and the interested reader is referred

 %   Mesnager \cite{SihemPOP2014} plateaued function over the Galois Field $F_{2^n}$.

%She
% introduced  new characterizations of $p$-ary plateaued functions \cite{SihemM2014} and  provided
%characterizations of plateaued functions (including vectorial functions) in arbitrary characteristic \cite{SihemPlate}. In \cite{SihemM2017AMC}, the authors provided a secondary construction of plateaued functions ( If three known bent functions in $n$ variables can  satisfy some conditions, a plateaued function in $n$ variables will be obtained ).
% A method for producing highly plateaued functions is presented \cite{THOMCusick2016}. Since the  use of disjoint spectra functions  provided a method for constructing resilient functions \cite{Zhang2009}, the sets of disjoint spectra plateaued functions were investigated \cite{Feng2,Fengrong2018IT}. In \cite{Feng2}, a method for constructing  quadruples of disjoint spectra  plateaued functions without  linear structures was provided. Recently, by using the indirect sum,  an efficient method for constructing a large set of disjoint spectra plateaued functions without linear structures was presented \cite{Fengrong2018IT}. In addition, a much larger set of disjoint spectra plateaued functions can be constructed by using the generalized indirect sum \cite{Fengrong2018IT}.


%The properties of plateaued functions were also widely studied.  A   characterization of plateaued functions was initiated in 2003 by Carlet and Prouff \cite{CarletQ}. Further, several  results about plateaued functions, especially on their characterizations in arbitrary
%characteristic (see for instance \cite{CesmDCC,SihemM2014,SihemPlate,Jong,SihemM2017,SihemM2017IT}).
%Recently, the value distribution of the
%second-order derivatives of plateaued functions (in particular cubic plateaued functions) were further investigated \cite{SihemM2017}. In \cite{SihemM2017IT}, the notion of plateaued
%functions was extended  to generalized functions by introducing the notion of
%gplateaued functions, which included the notion of generalized
%bent functions. In addition, the gplateaued functions were also studied and  several
%secondary constructions were provided.


%together with all functions which have exactly three values in spectrum. In this work, due to technical reasons, by a plateaued function we mean a function with three-valuated spectrum only. \textbf{In some other article(s) it is mentioned that $s$-plateaued functions (in a certain way) actually represent a generalization of bent functions, since it is about the value $s$ ($0$-palateaued are bent, etc..). That is the reason why I used "The concept of plateaued" above...}
%Certain design methods of plateaued functions have been suggested \cite{CarletAPN,DisjointSpect,Logathev,Mesnager} though these are rather specified as a particular algebraic normal form which has the desired property that the Walsh spectrum of such  a function is  3-valued taking the values in the set  $\{0,\pm 2^r\}$.

%the construction methods for previously mentioned functions are known as \emph{primary} and \emph{secondary}.
% Primary construction methods (referring mainly to bent functions), also referred to as direct construction methods,   employ a collection of suitable algebraic structures on $n/2$-dimensional subspaces  rather than using known   bent functions as their building blocks.
%Two well-known classes of bent functions derived form these primary methods are Dillon's partial spread class \cite{Dillon} and
 %or if it uses other bent functions then certain non-obvious conditions are required.

The main objective of this article is to completely revise the design approach by transferring it to the Walsh domain which gives much more flexibility and generality.
Indeed, any $s$-plateaued function $f$ on $GF(2)^n$ uniquely corresponds to its Walsh support, denoted by $S_f$ (a subset of $GF(2)^n$ for which the Walsh spectral values are nonzero $\pm 2^{\frac{n+s}{2}}$), and a sequence of $(\pm 1)$-signs that precisely specifies these  non-zero Walsh coefficients.
%is characterized by the fact that its Walsh support (a subset of $GF(2)^n$ for which the Walsh spectral values are nonzero $\pm 2^{\frac{n+s}{2}}$), denoted by $S_f$,  is of cardinality $2^{n-s}$.
The cardinality  of these positive and negative values in the Walsh support  is also known and it corresponds to the cardinality of zeros and ones of a bent function defined on $GF(2)^{n-s}$. This implies that an alternative approach to design plateaued functions on $GF(2)^n$ is to provide a proper placement of the spectral values  $\{0,\pm 2^{\frac{n+s}{2}}\}$ over the vector space $GF(2)^n$ so that this placement indeed yields a Boolean function (which is then plateaued). A useful categorization of  plateaued functions in this context is their separation  into two classes that we call {\em trivial} respectively {\em nontrivial} plateaued functions. The former class corresponds to the case when the Walsh support of a plateaued function is an affine subspace. In this case a plateaued function corresponds to a partially bent function and therefore for convenience it is called trivial (see relation (3.1) in \cite{Ayca2}). Furthermore, identifying the signs of nonzero values in the Walsh spectrum of an $s$-plateaued function with any  $(\pm 1)$-bent  sequence on $GF(2)^{n-s}$ gives an efficient design method for trivial plateaued functions. In particular, if our goal is to design a set of disjoint spectra $s$-plateaued functions of maximal cardinality $2^s$ it is sufficient to define these plateaued functions on the cosets of some fixed affine subspace $S_f$ of dimension $n-s$. In difference to previous approaches \cite{Zhang2009,Feng2,Fengrong2018IT} this method is generic, provides a greater variety of the obtained disjoint spectra plateaued functions, and finally the cardinality of these disjoint spectra functions is maximal. We emphasize the fact that any affine subspace $S_f$ and any bent function on $GF(2)^{n-s}$ can be selected in the design, which then essentially exhausts the possibilities of specifying trivial plateaued functions.

To address the problem of EA-equivalence of trivial and nontrivial plateaued functions on $GF(2)^n$, whose Walsh support $S_f \subset \FB^n$ is of cardinality $2^{n-s}$
%the positive and negative values in the Walsh domain for a plateaued function whose Walsh support $S_f$ is an arbitrary subset of $GF(2)^n$ of cardinality $2^{n-s}$
we use the concept of a {\em dual} of plateaued function, denoted by $f^*$, which essentially gives a correspondence to the (lexicographically) ordered vector space $GF(2)^{n-s}$. It refers to a function associated  to the signs of non-zero Walsh coefficients and is  viewed as a function on $GF(2)^{n-s}$ after fixing a particular ordering on the corresponding Walsh support.
%The question concerning extended affine (EA) equivalence of nontrivial and trivial plateaued functions is also treated in details.
While trivial and nontrivial plateaued functions are always inequivalent, we show that the EA-equivalence between the functions inside these two classes can be completely described in terms of their Walsh supports and duals.
% where the notion of {\em dual} of a plateaued function refers to a function associated  to the signs of non-zero Walsh coefficients and is  viewed as a function on $GF(2)^{n-s}$ after fixing a particular ordering on the corresponding Walsh support. At the same time, using the spectral design idea for construction of plateaued functions introduced in \cite{Secondary}, the complete description of inequivalent plateaued functions that we provide in addition gives a technique for construction of inequivalent plateaued functions, regardless of whether they are trivial or not. In addition, we provide a sufficient condition that nontrivial plateaued functions do not possess linear structures which is later used to confirm the existence of nontrivial disjoint spectra of plateaued functions of maximal cardinality without linear structures.
%\textbf{Samir: It is hard to speak about equivalence results without mentioning the word "dual". This word is used in the next paragraph, and therefore these two paragraphs needs to be adopted. In my opinion, the blue part below must be mentioned much earlier, and then it allows us to use the word "dual" freely in all these paragraphs (instead of sequence of signs at multiple places). Not to forget that the dual is important for the construction of both trivial and nontrivial functions (and reader should not be confused that we are using the word "dual" only when treating nontrivial functions, but not trivial ones...) OR maybe to switch the positions of the previous and next paragraph??}

We first completely solved the question related to the design of trivial plateaued functions in the spectral domain and then we investigate the possibilities of designing  nontrivial ones. In this case, the Walsh support is not an affine subspace which then makes a direct correspondence to bent functions on $GF(2)^{n-s}$ more difficult.
%\textcolor[rgb]{0.00,0.00,1.00}{To address the problem of assigning the positive and negative values in the Walsh domain for a plateaued function whose Walsh support $S_f$ is an arbitrary subset of $GF(2)^n$ of cardinality $2^{n-s}$ we use the concept of a dual of plateaued function, denoted by $f^*$, which essentially gives a correspondence to the (lexicographically) ordered vector space $GF(2)^{n-s}$.}
The property of being plateaued is then characterized through the bent distance of the dual $f^*$ to so-called {\em sequence profile}  of $S_f$. In difference to trivial plateaued functions the sequence profile of nontrivial plateaued functions necessarily contains sequences of nonlinear functions. In this direction, we give efficient design methods of nontrivial plateaued function which are apparently affine inequivalent to trivial ones. In particular, we also provide an efficient method for constructing nontrivial disjoint spectra plateaued functions of maximal cardinality $2^s$. It is shown that these functions cannot have linear structures which essentially improves upon recent results in \cite{Fengrong2018IT} where large sets of disjoint spectra plateaued functions, but of cardinality less than $2^s$, were specified.




%In the last part of this article we apply a similar method, used to define EA-inequivalent bent functions introduced originally in \cite{HouLang} and later extended in \cite{Houextension2017}, to generate EA-inequivalent plateaued functions by applying a suitable nonlinear permutations of the input space.
Our analysis leads to two significant research problems that we only address partially.
The first question regards the notion of {\em extendable} set of disjoint spectra $s$-plateaued functions. More precisely, having demonstrated that disjoint spectra $s$-plateaued functions of maximal cardinality can always be specified for trivial plateaued functions and in certain cases for nontrivial ones, the question is whether a given set $\{f_1, \ldots,f_r\}$  of arbitrary nontrivial $s$-plateaued functions can always be extended to a set of disjoint spectra $s$-plateaued functions of maximal cardinality.
% in which case the function is called {\em extendable} is left as an important open problem.
% \textbf{Samir: This has to be adopted according to Section 5 - the extendability is something that can be considered/posed also for other plateaued functions, see Open problem 3!}
Another question which deserves a full attention of the research community is whether the flexibility of defining disjoint spectra $s$-plateaued functions of maximal cardinality might lead to new classes of bent functions. More precisely, using  a set of $2^s$ disjoint spectra of $s$-plateaued functions on $GF(2)^n$ one can easily construct a bent function on the space $GF(2)^{n+s}$ by concatenating these plateaued functions. Then, since the duals of these plateaued functions may be taken from different classes of bent functions it is very challenging to conjecture that the resulting bent function  does not necessarily belong to known primary classes of bent functions. This question is extremely important due to the fact that the primary classes of bent functions only constitute an insignificant
%\textbf{Maybe to say "small" instead of insignificant? Even 1 bent functions can be significant, if we know how to make new ones out of it :)}
portion of all bent functions.

In the last part of this article we apply a similar method, used to define EA-inequivalent bent functions introduced originally in \cite{HouLang} and later extended in \cite{Houextension2017}, to generate EA-inequivalent plateaued functions by applying a suitable nonlinear permutations of the input space. This approach can also be used to specify nontrivial disjoint spectra plateaued functions but the conditions on the corresponding nonlinear permutations (acting on the input variable space) become rather complicated.

The rest of this article is organized as follows. In Section \ref{sec:pre}, we introduce some basic definitions and the properties of bent and plateaued functions. The EA-equivalence of plateaued functions and their classification into trivial and nontrivial ones is addressed in Section \ref{sec:equivalence}. Furthermore, a necessary and sufficient condition that nontrivial plateaued functions admit linear structures ia also given. In Section \ref{sec:nontrivial}, we provide some generic methods of constructing (non)trivial plateaued functions and show that these methods also gives us the possibility of defining a set of disjoint spectra plateaued functions of maximal cardinality. To generate EA-inequivalent plateaued functions,  we apply in Section \ref{sec:nonlperm} a similar technique of applying suitable nonlinear permutations of the input space, as already employed in \cite{HouLang,Houextension2017}.
%We introduce the notion of \textcolor[rgb]{0.00,0.00,1.00}{extendable plateaued functions} \textbf{To adopt this part according to Section 5!} as an attempt to structurally distinguish their intrinsic properties concerning their specification in the spectral domain.


%  In this context, a useful result is Titsworth theorem which gives both necessary and sufficient condition that a Boolean function specified in the spectral domain is Boolean. Based on this result we deduce some interesting properties which allows to specify plateaued functions in the spectral domain. It will be shown that a class of {\em trivial plateaued functions}  can be easily deduced by taking an arbitrary affine subspace and specifying  the nonzero spectral values on such a subspace. The design and analysis of {\em nontrivial plateaued functions}, whose nonzero spectral values are not defined on an affine subspace, is much more difficult. We show that there is a possibility to specify such functions in certain cases. Also ...........................
%
%%In the second part of this article we address the problem of specifying 5-valued spectra functions having its Walsh spectra values in the set $\{0,\pm 2^r, \pm 2^s\}$ ...
%{\bf BLA BLA BLA }


% and the corresponding conditions imposed on initial functions in this case are in general  hard to satisfy.
% compared to the methods that extend the variable space.
%.} % (cf. Example~\ref{ex:5_3}) of new families of bent functions defined on the same variable space as the initial functions.
%In Remark \ref{gbent} we  emphasize the fact that these  conditions are related to so-called  generalized bent (gbent) functions, which where characterized recently in \cite{SHWM,tang}.
%In general, the secondary constructions which do not increase the variable space seem to induce  harder sufficient bent conditions  than those constructions that generate bent functions on larger  variable spaces. The main problem that remains to be considered, in both cases, regards the question to which classes these bent functions belong to. It should be noticed that depending on the selection of initial functions we might easily have the case that these methods provide functions that are not contained in the known primary classes of bent functions. This classification is however hard in general.






\section{Preliminaries}\label{sec:pre}

The vector space $\mathbb{F}_2^n$ is the space of all $n$-tuples $x=(x_1,\ldots,x_n)$, where $x_i \in \mathbb{F}_2$.
 %with the standard operations.
For $x=(x_1,\ldots,x_n)$ and $y=(y_1,\ldots,y_n)$  in $\mathbb{F}^n_2$, the usual scalar (or inner) product over $\mathbb{F}_2$ is defined as $x\cdot y=x_1 y_1+\cdots+ x_n y_n.$ The Hamming weight of  $x=(x_1,\ldots,x_n)\in \mathbb{F}^n_2$ is denoted and computed as  $wt(x)=\sum^n_{i=1} x_i.$ %By "$\sum$" we denote the integer sum (without modulo evaluation), whereas "$\bigoplus$" denotes the sum evaluated modulo two.

The set of all Boolean functions in $n$ variables, which is the set of mappings from $\mathbb{F}_2^n$ to $\mathbb{F}_2$, is denoted by $\mathcal{B}_n$.  Especially, the set of affine functions in $n$ variables is given by $\mathcal{A}_n=\{a\cdot x+ b\;|\;a\in\mathbb{F}_2^n,\; b\in\{0,1\}\},$ and similarly  $\mathcal{L}_n=\{a\cdot x:a\in\mathbb{F}_2^n\}\subset \mathcal{A}_n$ %(also, $\mathcal{A}^1_n=\{a\cdot x\oplus 1 :a\in\mathbb{F}_2^n\}$)
 denotes the set of linear functions. It is well-known that any $f:\mathbb{F}^n_2 \rightarrow \mathbb{F}_2$ can be uniquely represented by its associated algebraic normal form (ANF) as follows:
\begin{eqnarray}\label{ANF}
f(x_1,\ldots,x_n)={\sum_{u\in \mathbb{F}^n_2}{\lambda_u}}{(\prod_{i=1}^n{x_i}^{u_i})},
\end{eqnarray}
where $x_i, \lambda_u \in \mathbb{F}_2$ and $u=(u_1, \ldots,u_n)\in \mathbb{F}^n_2$. The support of an arbitrary function $f\in \mathcal{B}_n$ is defined as $supp(f)=\{x\in \mathbb{F}^n_2:f(x)=1\}.$ %We say that a function $f\in \mathcal{B}_n$ is \emph{balanced} if $wt(f)=2^{n-1}.$

For an arbitrary function $f\in \mathcal{B}_n$, the set of its values on $\mathbb{F}^n_2$ (\emph{the truth table}) is defined as $T_f=(f(0,\ldots,0,0),f(0,\ldots,0,1),f(0,\ldots,1,0),\ldots,f(1,\ldots,1,1))$. The corresponding $(\pm 1)$-{\em sequence of $f$} is defined as
$\chi_f=((-1)^{f(0,\ldots,0,0)},(-1)^{f(0,\ldots,0,1)},(-1)^{f(0,\ldots,1,0)}\ldots,$ $(-1)^{f(1,\ldots,1,1)})$. The {\em Hamming distance} $d_H$ between two arbitrary Boolean functions, say $f,g\in \mathcal{B}_n,$ we define by $d_H(f,g)=\{x\in \mathbb{F}^n_2:f(x)\neq g(x)\}=2^{n-1}-\frac{1}{2}\chi_f\cdot \chi_g$, where $\chi_f\cdot \chi_g=\sum_{x\in \mathbb{F}^n_2}(-1)^{f(x)+ g(x)}$.
%The concatenation of two functions, say $f$ and $g$, is denoted by $f||g$ which means that the sequence of $f||g$ is obtained by merging the sequences of  $f$ and $g,$ i.e., $\chi_{f||g}=(\chi_f,\chi_g).$ Equivalently, the truth table of  $f||g$ is obtained by merging the truth tables of $f$ and $g$, i.e., $T_{f||g}=(T_f,T_g)$.
%\textcolor[rgb]{0.98,0.00,0.00}{For any two sets $A=\{\alpha_1,\ldots,\alpha_{r}\}$ and $B=\{\beta_1,\ldots,\beta_{r}\}$, let $A\wr B=\{(\alpha_i,\beta_i): i=1,\ldots,r\}$.} %For arbitrary sized sets of binary vectors, the Kronecker product of $A$ and $B$ is $A\times B=\{(\alpha,\beta):\alpha\in A,\; \beta\in B\}$.
By $\textbf{0}_k=(0,\ldots,0)$ we denote the all-zero vector in $\mathbb{F}^k_2.$


The \emph{Walsh-Hadamard transform} (WHT) of $f\in\mathcal{B}_n$, and its inverse WHT, at any point $\omega\in\mathbb{F}^n_2$ are defined respectively by
\begin{eqnarray}\label{WHT}
W_{f}(\omega)=\sum_{x\in \mathbb{F}_2^n}(-1)^{f(x)+ \omega\cdot x},\;\;\;\;(-1)^{f(x)}=2^{-n}\sum_{\omega\in \mathbb{F}_2^n}W_f(\omega)(-1)^{\omega\cdot x}, \;\;\;\omega\in \mathbb{F}^n_2.
\end{eqnarray}
%The inverse WHT is given by
%\begin{eqnarray}\label{IWHT}
%(-1)^{f(x)}=2^{-n}\sum_{\omega\in \mathbb{F}_2^n}W_f(\omega)(-1)^{\omega\cdot x}, \;\;x\in \mathbb{F}^n_2.
%\end{eqnarray}
%A \emph{vectorial Boolean function}, say $H:\mathbb{F}^n_2\rightarrow\mathbb{F}^k_2$, can be represented uniquely as $H(x)=(h_1(x),\ldots,h_k(x))$, where $h_i:\mathbb{F}^n_2 \rightarrow \mathbb{F}_2$. For any fixed non-zero vector $\omega\in \mathbb{F}^k_2$ and any $u\in \mathbb{F}^n_2,$ the notation $W_{\omega\cdot H}(u)$ will denote the WHT  of the  {\em component} function  $\omega\cdot H=\omega_1h_1 \oplus \cdots \oplus \omega_k h_k$ of the function $H$ at point $u$.

%For any subset $U\subseteq \mathbb{F}^n_2,$  $\phi_U$ will denote a Boolean function for which $\phi_U(x)=1$ if and only if $x\in U.$ Also,
For a set $U\subseteq \mathbb{F}^n_2$, by $U^{\bot}$ we denote the set  $U^{\bot}=\{y\in \mathbb{F}^n_2:x\cdot y=0,\;\forall x\in U\}.$ The cardinality of any set $U$ is  denoted by  $\#U.$ The \emph{Sylvester-Hadamard}  matrix of size $2^k \times 2^k$, is defined recursively as:
%denoted by $H_{2^{k}},$ which is constructed using Kronecker product $H_{2^{k}}=H_2\otimes H_{2^{k-1}},$ where
\begin{eqnarray*}\label{HM}
H_1=(1);\hskip 0.4cm H_2=\left(
                           \begin{array}{cc}
                             1 & 1 \\
                             1 & -1 \\
                           \end{array}
                         \right);\hskip 0.4cm H_{2^k}=\left(
      \begin{array}{cc}
        H_{2^{k-1}} & H_{2^{k-1}} \\
        H_{2^{k-1}} & -H_{2^{k-1}} \\
      \end{array}
    \right).
\end{eqnarray*}
The $i$-th row of $H_{2^k}$ we denote by $H^{(i)}_{2^k}$ ($i\in[0,2^k-1]$). For any two sets $A=\{\alpha_1,\ldots,\alpha_{r}\}$ and $B=\{\beta_1,\ldots,\beta_{r}\}$, we define  $A\wr B=\{(\alpha_i,\beta_i): i=1,\ldots,r\}$.
%are denoted  by $h_{i,l}$, $0\leq l \leq 2^{k}-1,$ where $0\leq i \leq 2^{k}-1.$
%In general, for $f$ defined on $\mathbb{F}^n_2$ and a given set $S=\{\omega_0,\ldots,\omega_{2^m-1}\}\subseteq \mathbb{F}^n_2$ ($\#S=2^m$), we denote by $f|_{S}$ the restriction of $f$ to $S$ and the corresponding sequence as
%% $\chi^S_f$ we denote the  sequence of $f$ restricted to $S$, $f|_{S}$,  as
%$\chi^S_f=((-1)^{f(\omega_0)},(-1)^{f(\omega_1)},\ldots,(-1)^{f(\omega_{2^{m}-1})}),$ where $\omega_i\in S.$
% we define by formula (\ref{WHT}) as $W_{\omega\cdot h}(u).$


\subsection{Bent and plateaued functions and their duals}

%{\bf I have left this section in the original form as in the article on secondary constructions. We will later see what is essentially needed and whether we can efficiently employ the Titsworth theorem throughout the article ... } \\ \\
Throughout this article we use the following definitions and useful facts related to bent and plateaued functions.
%\begin{itemize}
 A function $f\in\mathcal{B}_n,$ for even $n$, is called {\em bent} if $W_f(u)=2^{\frac{n}{2}}(-1)^{f^*(u)}$
for a Boolean function $f^*\in \mathcal{B}_n$ which is also a bent function, called the {\it dual} of $f$.
%If $f^*=c\oplus f$, $c\in \{0,1\}$, then $f$ is called self-dual if $c=0$, or anti self-dual if $c=1$.
%Note that bent functions only exist for even $n$.
%Similarly, we say that a vectorial function $h=(h_1,\ldots,h_k):\mathbb{F}^n_2\rightarrow\mathbb{F}^k_2$ is bent if for any non-zero vector $\omega\in \mathbb{F}^k_2,$ the function $\omega\cdot (h_1,\ldots,h_k)$ is a bent function on $\mathbb{F}^n_2.$ Due to Nyberg \cite{Nyberg}, vectorial bent functions exist only for $n$ even and $k\leq \frac{n}{2}.$ {\bf Do we need this or this is just adding more unnecessary information ?? I could not find it later in the text, and even if I could it is more logical to comment at place it appears !!}
% \textcolor[rgb]{0.00,0.00,1.00}{In the same way, we have the definition of vectorial plateaued functions, where a plateaued Boolean function is defined below. }
Two functions $f$ and $g$ on $\FB^n$ are said to be at {\em bent distance} if $d_H(f,g)= 2^{n-1}\pm 2^{n/2-1}$. Similarly, for a subset $B\subset \mathcal{B}_n$, a function $f$ is said to be at bent distance to $B$ if for all $g\in B$ it holds that $d_H(f,g)=2^{n-1}\pm 2^{n/2-1}$.

A function $f\in \mathcal{B}_n$ is called {\em $s$-plateaued} if its Walsh spectrum only takes three values $0$ and $\pm 2^{\frac{n+s}{2}}$ ($\leq 2^n$), where $s\geq 1$ if $n$ is odd and $s\geq 2$ if $n$ is even ($s$ and $n$ always have the same parity). The Walsh distribution of $s$-plateaued functions (cf. \cite[Proposition 4]{Decom}) is given by
\begin{table}[H]
\caption{Walsh spectra distribution}
\label{tab:plateaued}
%\caption{Walsh spectra distribution}
\begin{eqnarray*}
%\scriptsize
\begin{tabular}{|c|c|}
  \hline
  % after \\: \hline or \cline{col1-col2} \cline{col3-col4} ...
  $W_f(u)$ & Number of $u\in \mathbb{F}^n_2$ \\ \hline
  $0$ & $2^n-2^{n-s}$ \\ \hline
  $2^{\frac{n+s}{2}}$ & $2^{n-s-1}+(-1)^{f(0)}2^{\frac{n-s}{2}-1}$ \\ \hline
  $-2^{\frac{n+s}{2}}$ & $2^{n-s-1}-(-1)^{f(0)}2^{\frac{n-s}{2}-1}$ \\
  \hline
\end{tabular}
\end{eqnarray*}
%\caption{Walsh spectra distribution}
\end{table}
% Clearly, the class of $0$-plateaued functions is a class of bent functions.
In particular, a class of $1$-plateaued functions for $n$ odd, or  $2$-plateaued for $n$ even, corresponds to so-called {\em semi-bent} functions.
%  and $n$-plateaued functions are linear (affine) functions. \textcolor[rgb]{0.98,0.00,0.00}{In this paper, by an $s$-plateaued function defined on $\mathbb{F}^n_2$ ($1\leq s\leq n-1$), we mean a function which is not bent and which have exactly three values in its Walsh spectrum.}\textbf{In many other articles, authors use $0\leq s\leq n$, which includes affine and bent functions. By "a plateaued function" I want to call a functions which is neither affine/bent and which have three valuated spectrum. Just want to distinguish better these three different functions...}
%\item \textcolor[rgb]{0.98,0.00,0.00}{For a function $f\in \mathcal{B}_n$ and any vector $u\in \mathbb{F}^n_2$, the autocorrelation of $f$ is defined as $\Delta_f(u)=\sum_{x\in \mathbb{F}^n_2}(-1)^{f(x)\oplus f(x\oplus u)}$. We say that $f$ is partially bent function \cite{CarletPartial} if $\# S_f\cdot \# \mathfrak{R}=2^n$, where $\mathfrak{R}=\{u\in \mathbb{F}^n_2:\; \Delta(u)\neq 0\}.$}
In general, the {\em Walsh support} of  $f\in \mathcal{B}_n$ is defined as $S_f=\{\omega\in \mathbb{F}^n_2\; :\; W_f(\omega)\neq0\}.$

 For an arbitrary  $s$-plateaued function $f\in \mathcal{B}_n$ with $W_f(u)\in \{0,\pm 2^{\frac{n+s}{2}}\}$ the value $2^{\frac{n+s}{2}}$ is called the \emph{amplitude} of $f$. We define its dual function $f^*$ on the set $S_f$ of cardinality $2^{n-s}$ by $W_f(\omega)=2^{\frac{n+s}{2}}(-1)^{f^*(\omega)},$ for $\omega\in S_f$.
%\footnote{In some other articles, the dual of plateaued functions is defined by including $W_f(\omega)=0$ on $\mathbb{F}^n_2\backslash S_f$, i.e., $f^*(\omega)=0$ if $W_f(\omega)=0$, and $f^*(\omega)=1$ if $W_f(\omega)\neq 0$ ($\omega\in \mathbb{F}^n_2$). However, this definition of dual $f^*$ does not distinguish positive and negative values of non-zero coefficients $W_f(\omega)$, neither it uniquely determines the corresponding plateaued function $f$. Our definition of the dual seems  to be more sensible due to possibility of establishing a strong connection between the dual and bent functions on $\FB^{n-s}$.}.
     %
%     {\bf From here including some parts on the next page nothing is clear !!!! In the first place, the vectors $\omega_{i_j}$ belong to $\mathbb{F}^n_2$ and writing  $S_f=\{\omega_{i_0},\omega_{i_1},\ldots,\omega_{i_{2^{n-s}-1}}\}\subseteq \mathbb{F}^{n-s}_2$ is just wrong !!! Why having double index here, not enough with $\omega_{i}$ ? If not a multi-set $S_f$ and it cannot be w.r.t. $\mathbb{F}^n_2$ then why $\subseteq$ ? The same applies to $E$, where are these elements $e_i$ in $\mathbb{F}^{n-s}_2$ ?? If yes, then $E=\mathbb{F}^{n-s}_2$, if not $E \subseteq \FB^{n-s}$ is wrong as above. This is much more imprecise now than it was !! Maybe it is best to explain in one page if necessary what are you trying to do and then I can reduce it later. What I think you are doing here is to try to identify a subset of $\FB^n$ of cardinality $2^{n-s}$ with $\FB^{n-s}$ but in a very cumbersome way !! It does not need Carlet to have several comments on the notation here and the basic idea ! If this part is truly important then it must be very clear !! Are you trying to say that we order the subset of $\FB^n$ (using integer representation e.g.) and then associate the elements one by one to elements in $\FB^{n-s}$ ??? }
     %
To specify the dual function as $f^*:\FB^{n-s} \rightarrow \FB$ we use the concept of {\em lexicographic ordering}. That is, a subset  $E=\{e_0,\ldots,e_{2^{n-s}-1}\}\subset \mathbb{F}^{n}_2$ is ordered lexicographically if $|e_i| < |e_{i+1}|$ for any $i \in [0,2^{n-s}-2]$, where $|e_i|$ denotes the integer representation of $e_i \in \mathbb{F}^n_2$. More precisely, for $e_i=(e_{i,0}, \ldots, e_{i,n-1})$ we have $|e_i|= \sum_{j=0}^{n-1}e_{i,n-1-j} 2^j$, thus having the most significant bit of $e_i$ on the left-hand side.
Since $S_f$ is not ordered in general, {\em we will always represent it} as $S_f=v + EM$,\footnote{The reason for embedding $M$ to describe $S_f$ and not use simply $S_f=v +E$ will be explained in Section \ref{sec:class}, in the context of EA-equivalence of plateaued functions.} for some lexicographically ordered set $E$ with $e_0=\textbf{0}_{n}$, $M\in GL(n,\mathbb{F}_2)$ and $v \in S_f$. Here, $GL(n,\mathbb{F}_2)$ denotes the group of all invertible $\mathbb{F}_2$-linear transformations on $\mathbb{F}^n_2$, i.e., the group of invertible binary matrices of size $n \times n$.

For instance, if $S_f=\{(0,1,0),(0,1,1),(1,1,1), (1,0,1)\}$, by fixing $v=(0,1,1)\in S_f$ and $M$ to be the identity matrix $I_{3\times 3}$, then $E=\{e_0,e_1,e_2,e_3\}=\{(0,0,0), (0,0,1),(1,0,0),(1,1,0)\}$ is ordered lexicographically and consequently $S_f=v+ E$ is "ordered" as $S_f=\{\omega_0,\omega_1,\omega_2,\omega_3\}=\{(0,1,1),(0,1,0),(1,1,1),(1,0,1)\}$. Notice that in this example $S_f$ is not an affine subspace.


 This way we can make a direct correspondence between  $\FB^{n-s}$ and $S_f$ through $E$ so that for $\FB^{n-s}=\{x_0, x_1, \ldots, x_{2^{n-s}-1}\}$, where $\FB^{n-s}$ is lexicographically ordered, we assign
\begin{eqnarray}\label{DPL}
 \overline{f}^*(x_j)=f^*(v+e_jM)=f^*(\omega_j),\;\;\; x_j \in \FB^{n-s}, \; e_j \in E,\;j\in[0,2^{n-s}-1].
 %\leftrightsquigarrow f^*(e_j)
 \end{eqnarray}
%{\bf Is this clear enough ?? We do not mention $\omega$ in Eq. (\ref{DPL})??}
Then, the set $S_f=\{\omega_0,\ldots,\omega_{2^{n-s}-1}\}$ is "ordered" so that $\omega_i=v+ e_i M$, where $E=\{e_0,\ldots,e_{2^{n-s}-1}\}$ is ordered lexicographically.  In the above example, we have for instance that $\overline{f}^*(x_1)=\overline{f}^*(0,1)= f^*(0,1,0)=f^*(\omega_1)$, $x_1\in \mathbb{F}^2_2$.
 %$f^*(x_1)=f^*(0,1)\leftrightsquigarrow f^*(0,1,0)=f^*(\omega_1)$, $x_1\in \mathbb{F}^2_2$.%, where the latter input to $f^*$ is the second  input in $S_f$.
This simply means that the truth table of $\overline{f}^*:\FB^{n-s} \rightarrow \FB$ is fully specified by assigning its  values with respect to   the signs of the Walsh spectral values in the ordered set $S_f=\{\omega_0,\omega_1,\ldots,\omega_{2^{n-s}-1}\}$.
%An additional example which illustrates the definition of the dual $f^*$ on the vector space $\FB^{n-s}$ through relation (\ref{DPL}) (for $M=I_{n\times n}$) is given in \cite[Appendix]{Secondary}.
%\begin{rem}
%Formally, one can specify  $P$ to be a mapping from $S_f=v\oplus E$ to $\mathbb{F}^{n-s}_2$ ($E=\{e_0,\ldots,e_{2^{n-s}-1}\}$ is ordered lexicographically, %$\omega_i=v\oplus e_i)$ defined by $P:\omega_i\rightarrow x_i$, so that  the identification  in (\ref{DPL}) is given by $f^*(x_i)=f^*(P(\omega_i))$. Throughout the article, we will always consider $f^*(\omega_i)$ as $f^*(x_i)$ without mentioning the mapping $P$, using the fact that $S_f=v \oplus E$ is sorted with respect to some vector $v \in S_f$ and a lexicographically ordered set $E$. It appears to be difficult to capture  intrinsic properties of the mapping $P$ which would essentially establish the same results.
%\end{rem}
\begin{rem}\label{rem:dual}
This definition of  $\overline{f}^*$ (given by (\ref{DPL})), when viewed as a function on $\FB^{n-s}$ is not unique, since the given set $S_f \subset \FB^n$ can be represented in many ways as $S_f=v+ EM$, for some $v \in \FB^n$, $M \in GL(n,\mathbb{F}_2)$ and lexicographically ordered set $E$.  This means that the same function $f$ gives rise to different duals defined on $\mathbb{F}^{n-s}_2$ through the relation (\ref{DPL}).
%%with respect to representation of its Walsh support as $S_f=v\oplus EM$ actually enhances the definition given in \cite{Secondary}, since the representation of $S_f$ used in \cite{Secondary} is a special case when $M$ is an identity matrix.
%In this context, it should be clear that the definition of $f^*$ as a function defined on $\mathbb{F}^{n-s}_2$ by relation (\ref{DPL}) strongly depends on the ordering of $S_f$, which means that different orderings of $S_f$ may impose different duals of the same function.
%Therefore, in this work by dual $f^*$ of a function $f$ we will always denote a function on $\mathbb{F}^{n-s}_2$ whose truth table is given by (\ref{DPL}) with respect to "some" ordering of $S_f$, which is  induced by some choice of $v\in \mathbb{F}^{n}_2$, $M\in GL(n,\mathbb{F}_2)$ and a lexicographically ordered set $E\subset \mathbb{F}^n_2$.
%\end{rem}
%{\bf Enes: I think that already here we need to introduce a mapping $P:S_f \rightarrow \FB^{n-s}$ and to avoid the use of $\leftrightsquigarrow$ operation !! Once we are very clear with this we can restate our results in Section 3.1 especially and see whether it simplifies the whole treatment !! It is of importance to restate the main results such as Proposition 3.2, Theorem 3.2 and Corollary 3.1 in terms of this embedded mapping rather then using this $\leftrightsquigarrow$ operation. I think that the reviewers will not like this weird association. We do not need to say anything about this $P$ mapping even though I think it can be fully specified for the case of affine subspaces, using Lemma 3.1 in our CF paper !!}\\\\
%\textbf{\textcolor[rgb]{0.98,0.00,0.00}{Samir:} Let us say that we want to use the mapping $P:S_f\rightarrow \mathbb{F}^{n-s}_2$ in our proofs. How to explain properly the equality "$z_i=c+\omega_iA^T\leftrightsquigarrow y_iB^{-1}+tB^{-1}$"? Here, $A$  and $B$ are not of the same size. Even using $P(\omega_i)=y_i$ does not help too much in my view (is $P$ linear? -- probably linear for affine Walsh supports with orderings involved, but what for non-affine Walsh supports??). The fact is that "$\leftrightsquigarrow$" means much more than mapping from $S_f$ and $\mathbb{F}^{n-s}_2$. I understand that "$\leftrightsquigarrow$" can bring some confusion in the article, but it is due to the fact that we can not describe it efficiently and fully at the same time. We use it in (\ref{DPL}), then we use it to identify $c+\omega_iA^T\leftrightsquigarrow y_iB^{-1}+tB^{-1}$, we use it to identify $\phi_b(\omega_i)=b\cdot \omega_i$ with $\phi_b(x_i)$. Hence, its meaning is understood from the context, in which  we are relating values of "some mappings that involve vectors from different spaces"....I do not have better explanation of the symbol "$\leftrightsquigarrow$"...If someone has a proposal on how we can replace "$\leftrightsquigarrow$" with $P$ in the article, then feel free to suggest "how to do that" (in my view $P(\omega_i)=y_i$ is not enough to clarify all places where we use "$\leftrightsquigarrow$"). Therefore, I am providing the remark below...}
%
%%{\bf Enes: Remove the remark below and shorten the above remark !! Or merge the two remarks but make it short !!}
%
%\begin{rem}
%Throughout the article, the symbol "$\leftrightsquigarrow$" means the equality (that is "$=$"), and it is used to emphasize the equality between the values of some function(s) which is (are) defined in vectors which belong to different spaces (for instance (\ref{DPL}) means that $f^*(\omega_i)=f^*(x_i)$, where $\omega_i\in \mathbb{F}^n_2$ and $x_i\in \mathbb{F}^{n-s}_2$).
%
An alternative way of describing the relation (\ref{DPL}) would be the use of a mapping, say $P:S_f\rightarrow \mathbb{F}^{n-s}_2$, and write (\ref{DPL}) as $f^*(\omega_j)=f^*(P(\omega_j))=\overline{f}^*(x_j)$ (where $f^*\circ P$ is viewed as $\overline{f}^*$ defined on $\mathbb{F}^{n-s}_2$), but its properties are not known due to the  translation/preservation of an ordering imposed on $S_f$ to $\mathbb{F}^{n-s}_2$ (see \cite[Remark 2.1]{Secondary}). Note that when $S_f$ is an affine subspace, then the mapping $P$ can be described in terms of (\ref{eq:Q}), but if $S_f$ is non-affine its description is not known.
\end{rem}
%Throughout this article we consider $E$ to be lexicographically ordered if not stated otherwise.
%\begin{rem}
%To avoid the confusion between binary variables $x_i \in \FB$ (used in the ANF representation) and the vectors $x_i \in \FB^k$ we always specify to which space $x_i$ belongs to, unless it is clear from the context. For convenience, a binary vector $u \in \FB^k$ whose integer representation is $d \in [0,2^k-1]$ is sometimes denoted by $u_{[d]}$.
%\end{rem}
%\textbf{Samir: Maybe to delete the remark 2.2, and just to explain the notation at the places where we use the notation $u_{[d]}$??}
%     More precisely, let a  (non-multi) subset $S\subset \mathbb{F}^{n}_2$ be equal to $v\oplus E$, for some lexicographically ordered (non-multi) subset  $E=\{e_0,\ldots,e_{2^{n-s}-1}\}\subset \mathbb{F}^{n}_2$ which contains zero vector $e_0=\textbf{0}_{n}$. The ordering of the elements $S_f=\{\omega_{i_0},\omega_{i_1},\ldots,\omega_{i_{2^{n-s}-1}}\}$, where $\omega_{i_j}=v\oplus e_j$, we define as follows. For $\omega_{i_j}, \omega_{i_t}\in S_f$, we say that $\omega_{i_j}\leq_{\sigma} \omega_{i_t}$ if and only if $e_j\leq_{lex} e_t$ (where $\omega_{i_r}=v\oplus e_r$, $r\in[0,2^{n-s}-1]$). It is not difficult to see (due to the lexicographic ordering of $E$) that this ordering of elements in $S_f$ is actually a partial ordering, and we will call it the \emph{$\sigma$-ordering}. For instance, consider in the lexicographic ordering the set $E=\{e_0,e_1,e_2,e_3\}=\{(0,0,0), (0,0,1),(1,0,1),(1,1,1)\}$ and $v=(0,1,0)$, the $\sigma$-ordering of the set $S_f=v\oplus E$ is
%     $$S_f=v\oplus E=v\oplus\left(
%             \begin{array}{ccc}
%               0 & 0 & 0 \\
%               0 & 0 & 1 \\
%               1 & 0 & 1 \\
%               1 & 1 & 1 \\
%             \end{array}
%           \right)=\left(
%                     \begin{array}{ccc}
%                       0 & 1 & 0 \\
%                       0 & 1 & 1 \\
%                       1 & 1 & 1 \\
%                       1 & 0 & 1 \\
%                     \end{array}
%                   \right)
%           =\left(
%              \begin{array}{c}
%                v\oplus e_0 \\
%                v\oplus e_1 \\
%                v\oplus e_2 \\
%                v\oplus e_3 \\
%              \end{array}
%            \right)
%           =\left(
%             \begin{array}{c}
%               \omega_{i_0} \\
%               \omega_{i_1} \\
%               \omega_{i_2} \\
%               \omega_{i_3} \\
%             \end{array}
%           \right).
%     $$
%     %
%     %Then, we say that $\omega_{i_j}<_{\sigma}\omega_{i_t}$ (where $\omega_{i_j}=v\oplus e_j$ and $\omega_{i_j}=v\oplus e_t$) if and only if $e_j<_{\ell}e_t$.
%     %
%     % For convenience, assuming that the space $S_f$ is ordered lexicographically (what we define later), say $S_f=\{\omega_0,\omega_1,\ldots,\omega_{2^{n-s}-1}\}$,
%     %
%     In connection to Lemma \ref{l1} (see Section \ref{sec:plateWS}), our identification of $f^*$ with a function in $(n-s)$ variables is achieved through indices, i.e., the value of $f^*$ at $\omega_{i_j}\in S_f\subset \mathbb{F}^n_2$ is considered to be the value of $f^*$ in $x_j\in \mathbb{F}^{n-s}_2$, i.e., we are considering the correspondence
% \begin{eqnarray}\label{DPL}
% f^*(\omega_{i_j})\leftrightsquigarrow f^*(x_j),\;\;\; j\in[0,2^{n-s}-1].
% \end{eqnarray}
%     Clearly, in this identification we actually have that $f^*(e_j)\leftrightsquigarrow f^*(x_j).$

%     It is important to note that values of $f^*$ are defined with respect to vectors $\omega\in S_f$, and thus the results deduced for $f^*$ throughout the article do not depend on ordering of the space $\mathbb{F}^n_2$. In other words, this definition of dual implies that an arbitrary  ($s$-plateaued) function $f$ can be uniquely identified through its Walsh support $S_f$ and its dual $f^*$.

%If $S_f$ is an affine subspace in $\mathbb{F}^n_2$, i.e., for some subspace $E \subseteq\mathbb{F}^n_2$ of dimension $m$ and vector $v\in \mathbb{F}^n_2,$ we say that $S_f=a\oplus E$ is also of dimension $m,$ i.e., $S_f$ is $m$-dimensional affine subspace (writing $dim(S_f)=m$).
%In our definition, it is clear that an arbitrary  $s$-plateaued function $f$ can be uniquely identified with its Walsh spectrum $S_f$ and dual $f^*$.
 %\item Two  functions $f,g\in \mathcal{B}_n$ are said to be disjoint spectra functions if $W_f(\omega)W_g(\omega)=0,$ for every $\omega\in \mathbb{F}^n_2$ \cite{CCA}.
% \item Two  functions $f, g\in \mathcal{B}_n$  are said to be affine equivalent if and only if there
%exist $A \in  GL(\mathbb{F}^n_2)$, $b \in \mathbb{F}^n_2$, $\mu\in \mathbb{F}^n_2$ and $\epsilon \in \mathbb{F}_2$ such that $g(x)=f(xA\oplus b)\oplus \mu\cdot x\oplus \epsilon,$ for all $x\in \mathbb{F}^n_2.$ Here, $GL(\mathbb{F}^n_2)$ is the group of all invertible $\mathbb{F}_2$-linear transformations on $\mathbb{F}^n_2.$
 %\item The notion of CCZ-equivalence between vectorial functions has been introduced in \cite{CarletCCZ}, but later named CCZ-equivalence in \cite{Buda}. For two vectorial functions $h,h':\mathbb{F}^n_2\rightarrow \mathbb{F}^k_2,$ are said to be CCZ-equivalent if their graphs $G_h=\{(x,y)\in \mathbb{F}^n_2\times \mathbb{F}^k_2: y=h(x)\}$ and $G_{h'}=\{(x,y)\in \mathbb{F}^n_2\times \mathbb{F}^k_2: y=h'(x)\}$ are affine equivalent, that is, if there exists an affine permutation $L=(L_1,L_2)$ of $\mathbb{F}^n_2\times \mathbb{F}^k_2$ ($L_1:\mathbb{F}^n_2\times \mathbb{F}^k_2\rightarrow \mathbb{F}^n_2$, $L_2:\mathbb{F}^n_2\times \mathbb{F}^k_2\rightarrow \mathbb{F}^k_2$) such that $L(G_h)=G_{h'}$, i.e., $y=h(x)$ if and only if $L_2(x,y)=h'(L_1(x,y)).$
%\end{itemize}



\section{Classification of plateaued functions}\label{sec:class}

In this section we discuss different ways to classify plateaued functions depending on their spectral characterization. In the first place, we observe that the standard notion of extended affine (EA) equivalence transfers to the spectral domain which enables us to distinguish EA-inequivalent classes of  plateaued functions. A plateaued function with affine Walsh support is then called trivial and its equivalence class is obtained by applying invertible affine transformation in the spectral domain. In difference to these trivial objects (that correspond to partially bent functions \cite[Relation (3.1)]{Ayca2}) there are nontrivial plateaued functions with non-affine Walsh support that are EA-inequivalent to any trivial plateaued function. By analyzing the EA-equivalence through relation (\ref{DPL}), we provide both description and generic construction of inequivalent plateaued functions, regardless of whether they are trivial or not. We also show that certain subclasses of nontrivial plateaued functions functions do not admit linear structures.

\subsection{On affine equivalence between (non)trivial plateaued functions}\label{sec:equivalence}

%In this section we analyze the EA-equivalence between two arbitrary plateaued functions in terms of its Walsh supprt and dual.
Since the design of plateaued functions described in \cite[Section 3]{Secondary} is related to the Walsh spectrum, it is convenient to analyze the EA-equivalence between two arbitrary plateaued functions in terms of their Walsh supports and duals. The  notion of EA-equivalence, which goes back to early works of Marioana-McFarland \cite{MMclass} and Dillon \cite{Dillon} can be stated as follows. Throughout this article $GL(n,\mathbb{F}_2)$ denotes the group of all invertible $\mathbb{F}_2$-linear transformations on $\mathbb{F}^n_2$.
%is is just a special case of the so called extended affine (EA) equivalence.
\begin{defi}\label{def:EA}
Two Boolean functions $h,f:\mathbb{F}^n_2\rightarrow \mathbb{F}_2$ are said to be EA-equivalent if there exists a matrix $A\in GL(n,\mathbb{F}_2)$,
%(where $GL(n,\mathbb{F}_2)$ denotes the group of all invertible $\mathbb{F}_2$-linear transformations on $\mathbb{F}^n_2$),
vectors $b,c\in \mathbb{F}^n_2$ and a constant $\varepsilon \in \{0,1\}$ such that
\begin{equation} \label{eq:EAeq}
h(x)=f(xA+b)+c\cdot x+\varepsilon.
\end{equation}
\end{defi}
This equivalence naturally applies to plateaued functions and if $f$ is an $s$-plateaued function on $\FB^n$ with the support $S_f$ then the Walsh spectra of $h(x)=f(xA+b)+c\cdot x+\varepsilon$ (in terms of relation (\ref{DPL})), at any $u\in \mathbb{F}^n_2$, is given by:
\begin{eqnarray}\label{eq:equiv}\nonumber
W_h(u)&=&\sum_{x\in \mathbb{F}^n_2}(-1)^{f(xA+b)+c\cdot x+\varepsilon+u\cdot x}=\sum_{y\in \mathbb{F}^n_2}(-1)^{f(y)+(u+c)A^{-T}\cdot y+(u+c)\cdot bA^{-1}+\varepsilon}\\
&=&\left\{\begin{array}{cc}
                                                                                                 2^{\frac{n+s}{2}}(-1)^{f^*((u+c)A^{-T})+(u+c)\cdot bA^{-1}+\varepsilon}=2^{\frac{n+s}{2}}(-1)^{h^*(u)}, & (u+c)A^{-T}\in S_f \\
                                                                                                 0, &  (u+c)A^{-T}\not\in S_f
                                                                                               \end{array}
\right..
\end{eqnarray}
From (\ref{eq:equiv}), we have that the Walsh supports of $h$ and $f$ are related by the equality $S_h=c+S_fA^T$. In order to establish a more precise connection between the duals of two EA-equivalent $s$-plateaued functions, say $f$ and  $h$ on $\mathbb{F}^n_2$, we fix the following notation.
\begin{itemize}
 \item We  always assume that  $S_f=v+E=\{\omega_0,\ldots,\omega_{2^{n-s}-1}\}$ is  ordered so that $\omega_i=v+e_i$, where $v\in S_f$ and $E=\{e_0,\ldots,e_{2^{n-s}-1}\}$ is ordered lexicographically ($e_0=\textbf{0}_n$).
\item  Since by (\ref{eq:equiv}) we have $S_h=c+S_fA^T$, the support of $h$ is  ordered as $S_h=\{z_0,\ldots,z_{2^{n-s}-1}\}$ with $z_i=c+\omega_iA^T$ $(i\in[0,2^{n-s}-1])$. With respect to  $S_f$ and $S_h$, the duals $\overline{f}^*,\overline{h}^*:\mathbb{F}^{n-s}_2\rightarrow \mathbb{F}_2$ (where e.g. $x_i \rightarrow \overline{f}^*(x_i$)) will be defined using the identification:
\begin{equation}
\label{eq:identif}
\overline{f}^*(x_i)= f^*(\omega_i);  \;\;\;\;\;\;\;\;  \overline{h}^*(x_i)= h^*(z_i); \; \; \;\; x_i \in \FB^{n-s}, \omega_i \in S_f, z_i \in S_h,
\end{equation}
where in particular (by relations (\ref{eq:equiv}) and (\ref{eq:identif})) we have that
 \begin{eqnarray}\label{eq:equiv22}
h^*(z_i)=f^*(\omega_i)+\omega_i\cdot b+\varepsilon. %;\;\;\; ~  \overline{h}^*(x_i)=\overline{f}^*(x_i)+x_i\cdot b+\varepsilon,\;
\end{eqnarray}
where $z_i=c+\omega_iA^T$  and $i\in[0,2^{n-s}-1]$.
\item
%Notice that in \cite{Secondary} the dual of $f^*$ is defined with respect to $S_f$ represented as $S_f=v+E$, where $v\in S_f$ and $E$ ordered lexicographically. As already pointed out in Remark \ref{rem:dual}, different orderings of $S_f$ give different duals in terms of relation (\ref{DPL}) and their connection is difficult to establish in general.
Using $S_f=v+E$ (with lexicographically ordered $E$) and expressing   $S_h=c+S_fA^T=\widehat{v}+EA^T$, where $\widehat{v}=c+vA^T\in \mathbb{F}^n_2$,  gives us the possibility to analyze the EA-equivalence in a sensible way by representing $S_h=\widehat{v} + EM$, where we have that $M=A^T$.
%due to the connection between duals  $f^*$ and $h^*$ imposed by (\ref{eq:equiv22}). %, where the equivalence between $f$ and $h$ (with supports $S_f= \{\omega_0,\ldots,\omega_{2^{n-s}-1}\}$ and $S_h=\{z_0,\ldots,z_{2^{n-s}-1}\}$) is partially given through $z_i=c+\omega_iA^T$. This allows us to identify the EA-equivalence through the natural connection between $M$ and $A^T$.
%Clearly, the Walsh support $S_f$ we always represent as $S_f=v+E$ (where $M=I_{n\times n}$) for simplicity, since we have the natural connection between $S_h=c+S_fA^T$ (with $M=A^T$) and $S_f$ which follows from relation (\ref{eq:equiv}).
\end{itemize}
%
%{\bf Where do we use the above representation $S_h=u + EM$ along with this "identification" ??} \\
% Assume that $S_f$ is given as $S_f=\{\omega_0,\ldots,\omega_{2^{n-s}-1}\}$ without using any particular ordering of $S_f$. Using the fact that $S_h=c+S_fA^T$, %we have that the support $S_h$ is given as $S_h=\{z_0,\ldots,z_{2^{n-s}-1}\}$ where the vectors $z_i$ and duals $h^*$, $f^*$ are related as
%\begin{eqnarray}\label{eq:equiv22}
%\left\{\begin{array}{c}
%          z_i=c+\omega_iA^T,\\
%         h^*(z_i)=f^*(\omega_i)+\omega_i\cdot b+\varepsilon,\;\;i\in[0,2^{n-s}-1]
%       \end{array}
%\right..
%\end{eqnarray}
%where $S_f=\{\omega_0,\ldots,\omega_{2^{n-s}-1}\}$, $S_h=\{z_0,\ldots,z_{2^{n-s}-1}\}$ and $z_i=c+\omega_iA^T$ ($i\in[0,2^{n-s}-1]$).
%\end{defi}
\begin{rem}\label{rem:ordering}
In relation (\ref{eq:equiv22}), if  $S_f$ is an ordered set then it  induces some  ordering of $S_h$ via the equality $z_i=c+\omega_iA^T$, and clearly the Walsh supports of any two equivalent plateaued functions $h$ and $f$ are then related using this connection. %Note that this does not mean that the EA-equivalence between $h$ and $f$ is dependent on any ordering, but in fact any two equivalent plateaued functions can be treated in terms of (\ref{eq:equiv22}).
%In context of relation (\ref{DPL}), notice that in $S_h=c+S_fA^T=(c+vA^T)+EA^T$ the matrix $M$ is equal to $A^T.$
\end{rem}
%At this point we would like to discuss the concept of dual of an $s$-plateaued function in $n$ variables given by relation (\ref{DPL}):
%\begin{itemize}
%\item Notice that in \cite{Secondary} the dual of $f^*$ is defined with respect to $S_f$ represented as $S_f=v+E$, where $v\in S_f$ and $E$ ordered lexicographically. As already pointed out in Remark \ref{rem:dual}, different orderings of $S_f$ give different duals in terms of relation (\ref{DPL}) and their connection is difficult to establish in general. \textcolor[rgb]{0.98,0.00,0.00}{The representation of $S_h$ in terms of (\ref{DPL}) given as $S_h=\widehat{v}+EA^T$ ($\widehat{v}=c+vA^T\in \mathbb{F}^n_2$, $M=A^T$) that we employ in \textbf{Notation} gives us the possibility to analyze the EA-equivalence in a sensible way due to the connection between duals  $f^*$ and $h^*$ imposed by (\ref{eq:equiv22}). %, where the equivalence between $f$ and $h$ (with supports $S_f= \{\omega_0,\ldots,\omega_{2^{n-s}-1}\}$ and $S_h=\{z_0,\ldots,z_{2^{n-s}-1}\}$) is partially given through $z_i=c+\omega_iA^T$. This allows us to identify the EA-equivalence through the natural connection between $M$ and $A^T$.
%Clearly, the Walsh support $S_f$ we always represent as $S_f=v+E$ (where $M=I_{n\times n}$) for simplicity, since we have the natural connection between $S_h=c+S_fA^T$ (with $M=A^T$) and $S_f$ which follows from relation (\ref{eq:equiv}).}
%\end{itemize}
In general, the equivalence between $h$ and $f$ is possible regardless of whether $S_h$ and $S_f$ are affine subspaces or not. %In general, the treatment of the former case ($S_f$ being an affine subspace) is much easier.
%Therefore, we provide the following definition  which also distinguishes these subclasses with respect to EA-equivalence.
%%at the same time classifies plateaued functions in terms of intrinsic difficulty of specifying subclasses of these functions.
\begin{defi} \label{def:nontr}
An $s$-plateaued function $f \in \mathcal{B}_n$ whose Walsh support $S_f$ is an affine subspace is called {\em trivial}, otherwise  $f$ is  called {\em nontrivial} if $S_f$ is not an affine subspace.
%those $s$-plateaued functions having $S_f$ which is not an affine subspace are called {\em nontrivial}.
\end{defi}
%
%For convenience, we provide the following example in which we also provide the definition of $f^*$ and $h^*$ in terms of relation (\ref{DPL}).
%
In order to proceed further with our analysis, we recall some  notation introduced in \cite{Secondary}. More precisely, for an arbitrary $s$-plateaued function $f$ defined on $\mathbb{F}^n_2$ let $S_f\subset \mathbb{F}^n_2$ ($\#S_f=2^{n-s}$) be its Walsh support ordered as $S_f=\{\omega_{0},\ldots,\omega_{2^{n-s}-1}\}$. The \emph{sequence profile} of $S_f$,
%\textbf{Samir: I think that "indicator set" term is not good, since we have already an indicator function. Maybe to say "supporting set of f"? Or some other term? The same refers to Section 5.1, since we have indicator set there...},
 which is a multi-set of $2^n$ sequences of length $2^{n-s}$ induced by $S_f$, is defined as %\cite{Secondary}
%By $\Phi_{f}$ we denote the multi-set of $2^k$ functions whose sequences are exactly given as $((-1)^{u\cdot \omega_{i_0}},\ldots,(-1)^{u\cdot \omega_{i_{2^{k-s}-1}}}),$ where $\omega_{i_j}\in S_f$ and $u\in \mathbb{F}^{k}_2$, i.e.,
\begin{eqnarray}\label{RSf}
\Phi_{f}=\{\phi_u:\mathbb{F}^{n-s}_2\rightarrow \mathbb{F}_2\;:\; \chi_{\phi_u}=((-1)^{u\cdot \omega_{0}},(-1)^{u\cdot \omega_{1}},\ldots,(-1)^{u\cdot \omega_{2^{n-s}-1}}),\;\omega_i\in S_f,\; u\in \mathbb{F}^{n}_2\}.
\end{eqnarray}
As noted in \cite{Secondary}, $\Phi_{f}$ depends on the ordering of $S_f$ and it is spanned by the functions $\phi_{b_1},\ldots,\phi_{b_n}$, i.e., $\Phi_{f}=\langle \phi_{b_1},\ldots,\phi_{b_n}\rangle$, where $b_1,\ldots,b_n$ is the canonical basis of $\mathbb{F}^n_2$ ($b_i$ contains the non-zero coordinate at the $i$-th position).

%Since we do not know what is the relation between  duals of $f$ which are defined with respect to different representations of $S_f$ when different vectors $v\in S_f$ are chosen, then by "dual" we always mean some function $f^*:\mathbb{F}^{n-s}_2$ which is defined as $f^*(x_i)\leftrightsquigarrow f^*(\omega_i)$, where $S_f=\{\omega_0,\ldots,\omega_{2^{n-s}-1}\}$ is given in "some" ordering (Remark \ref{rem:dual}).
%\item From the previous point, in relation (\ref{eq:equiv22}) by "dual of $h$" we mean that $h^*:\mathbb{F}^{n-s}_2\rightarrow \mathbb{F}_2$ is defined as $h^*(x_i)\leftrightsquigarrow h^*(z_i)$, where $S_h=\{z_0,\ldots,z_{2^{n-s}-1}\}$ is ordered such that $z_i=c+\omega_iA^T$ ($\omega_i\in S_f$). Thus, the reason why the function $h^*(x_i)\leftrightsquigarrow h^*(z_i)$ is considered to be the "dual of $h$" instead of some function, say $\widehat{h}^*:\mathbb{F}^{n-s}_2\rightarrow \mathbb{F}_2$ defined as $\widehat{h}^*(x_i)\leftrightsquigarrow \widehat{h}^*(w_i)$,  where $S_h=c+S_fA^T=\widehat{v}+\widehat{E}=\{w_0,\ldots,w_{2^{n-s}-1}\}$ (for some lexicographically ordered $\widehat{E}$ and $\widehat{v}\in \mathbb{F}^n_2$), is due to the previous point, i.e., due to the fact that we do not know how functions $h^*$ and $\widehat{h}^*$ are related. The analysis of the relation between these functions is left as an open problem. This open question forced us to use the invertible matrix $M$ in definition (\ref{DPL}) so that we can adopt the concept of the "dual" with the notion of EA-equivalence in relation (\ref{eq:equiv22}).
%\end{itemize}
The following result  combines the relations (\ref{eq:identif}), (\ref{eq:equiv22}) and (\ref{RSf}) to specify EA-(in)equivalence between two plateaued functions.
\begin{theo}\label{th:eq}
Let $f,h:\mathbb{F}^n_2\rightarrow \mathbb{F}_2$ be $s$-plateaued functions whose Walsh supports are given as $S_f=\{\omega_0,\ldots,\omega_{2^{n-s}-1}\}$ and $S_h=\{z_0,\ldots,z_{2^{n-s}-1}\}$. Suppose that the duals $\overline{f}^*, \overline{h}^*:\mathbb{F}^{n-s}_2\rightarrow \mathbb{F}_2$ are defined using (\ref{eq:identif}).
%$f^*(x_i)\leftrightsquigarrow f^*(\omega_i)$ and $h^*(x_i)\leftrightsquigarrow h^*(z_i)$ $(x_i\in \mathbb{F}^{n-s}_2)$.
Then:
\begin{enumerate}[i)]
\item If there does not exist a vector $c\in \mathbb{F}^n_2$ and a matrix $A\in GL(n,\mathbb{F}_2)$ such that $S_h=c+S_fA^T$, then $f$ and $h$ are EA-inequivalent.
\item $f$ and $h$ are EA-equivalent (in terms of (\ref{eq:EAeq})) if and only if there exist  $c,b\in \mathbb{F}^n_2$, $A\in GL(n,\mathbb{F}_2)$ and $\varepsilon \in \mathbb{F}_2$ such that $S_h=c+S_fA^T$ $($thus $z_i=c+\omega_iA^T)$ and $\overline{h}^*(x_i)=\overline{f}^*(x_i)+\phi_{b}(x_i)+\varepsilon$, for all $i \in [0,2^{n-s}-1]$. %, where $\phi_b(x_i)\leftrightsquigarrow \phi_b(\omega_i)$ for all $i\in[0,2^{n-s}-1]$.
\end{enumerate}
\end{theo}
\proof While $i)$ follows trivially from (\ref{eq:equiv}), we focus on $ii)$. From (\ref{eq:equiv}), we have that
$$h^*(z_i)=f^*((z_i+c)A^{-T})+(z_i+c)\cdot bA^{-1}+\varepsilon,\;\;\;\omega_i=(z_i+c)A^{-T}\in S_f.$$
Since  $z_i\in S_h$, then from $(z_i+c)A^{-T}\in S_f$ we can write $z_i=c+\omega_iA^T$, for some (unique) $\omega_i\in S_f$. Thus, the equality above can be written as
\begin{eqnarray*}
h^*(z_i)=h^*(c+\omega_iA^T)&=&f^*(\omega_i)+(c+\omega_iA^T+c)\cdot bA^{-1}+\varepsilon\\
&=&f^*(\omega_i)+\omega_iA^T\cdot bA^{-1}+\varepsilon=f^*(\omega_i)+\omega_i\cdot bA^{-1}A+\varepsilon\\
&=&f^*(\omega_i)+\omega_i\cdot b+\varepsilon\stackrel{(\ref{RSf})}{=}f^*(\omega_i)+\phi_b(\omega_i)+\varepsilon,
\end{eqnarray*}
which completes the proof due to relations (\ref{eq:identif}), (\ref{eq:equiv22}) and (\ref{RSf}).\qed  % and considering $\phi_b(\omega_i)\leftrightsquigarrow \phi_b(x_i)$, where $\phi_b:\mathbb{F}^{n-s}_2\rightarrow \mathbb{F}_2$ is defined by (\ref{RSf}).\qed}
%\textbf{Samir: I added this small proof due to Fengrong's comment that the computation is not clear to him.}\\\\
%
%
%\begin{rem}\label{rem:ordering2}
%Note that if $f$ and $h$ are EA-equivalent plateaued functions, with $S_f=v+E$, then there exist a unique  $A \in GL(n,\FB)$ and  $c \in \FB^n$ for which it holds that $S_h=c+S_fA^T$.
%\end{rem}
By Theorem \ref{th:eq}-$ i)$, we clearly have that trivial and nontrivial $s$-plateaued functions in $n$ variables cannot be EA-equivalent. In the case of nontrivial plateaued functions, we notice that their Walsh supports in general may contain different number of linearly independent vectors. This fact can be used as  an efficient indicator for EA-inequivalence  since affine mappings preserve the cardinality of linearly independent vectors and the equality $S_h=c+S_fA^T$ will not  always be satisfied (thus we can distinguish EA-inequivalent functions $f$ and $h$). Nevertheless, in the case when this equality holds (which also applies to trivial plateaued functions), then the EA-equivalence of $f$ and $h$  strictly depends on the relation between its duals $\overline{h}^*$ and $\overline{f}^*$, which is given in Theorem \ref{th:eq}-$(ii)$ as $\overline{h}^*=\overline{f}^*+\phi_b+\varepsilon$.
%{\bf Enes: Shorten if possible !!}

Now we focus on the analysis of trivial plateaued functions. We firstly provide the following result which slightly extends \cite[Lemma 3.1]{Secondary} by involving invertible matrices, where for conciseness we omit mentioning that $E$ is lexicographically ordered.
%for self-completeness we include all assumptions (although they resemble on those in \textbf{Notation})}. %, and later we show that Corollary \ref{cor:affineS_f} still holds if $f^*$ is defined by (\ref{DPL}) with respect to its Walsh support represented as $S_f=v\oplus EM$ when $E$ is a  lexicographically ordered linear subspace.
\begin{prop}\label{prop:Hrow}
Let $S=\{\omega_0,\ldots,\omega_{2^m-1}\}\subseteq\mathbb{F}^n_2$ be an affine subspace represented as $S=v+E$, where $E=\{e_0,\ldots,e_{2^m-1}\}\subset \mathbb{F}^n_2$ is a linear subspace.
% is a lexicographically ordered linear subspace.
% (i.e., $\omega_i=v+e_i$, $i\in[0,2^m-1]$).
Let for an arbitrary  $M\in GL(n,\mathbb{F}_2)$ and $c\in \mathbb{F}^n_2$ the set $\widetilde{S}$ be given by $\widetilde{S}=c+SM=\{\alpha_0,\ldots,\alpha_{2^m-1}\}$ ($\alpha_i=c+(v+e_i)M$, $i\in[0,2^m-1]$). Then, for any $u\in \mathbb{F}^n_2$, it holds that
$$((-1)^{u\cdot \alpha_0},\ldots, (-1)^{u\cdot \alpha_{2^m-1}})=(-1)^{\varepsilon_u}H^{(r_u)}_{2^{m}},$$
for some $0\leq r_u\leq 2^m-1$ and $\varepsilon_u\in \mathbb{F}_2$. In addition,  it holds that $\{T_{\ell}:\ell\in \mathcal{L}_m\}\subseteq \{(u\cdot e_0,\ldots,u\cdot e_{2^m-1}):u\in \mathbb{F}^n_2\}$. %, which means that  $\mathcal{L}_m$ is contained in a multi-set of $m$-variable linear functions whose truth tables are $\{(u\cdot e_0,\ldots,u\cdot e_{2^m-1}):u\in \mathbb{F}^k_2\}$.
%In addition, when $u$ runs through  $\mathbb{F}^n_2$
%then $r_u$ is surjective on $\{0,\ldots,2^m-1\}.$
\end{prop}
\begin{proof} Denoting by $c'=c+vM\in \mathbb{F}^n_2$, we have that $\alpha_i=c+(v+e_i)M=c'+e_iM$ for all $i\in[0,2^m-1]$. Consequently, for any $u\in \mathbb{F}^n_2$, we have that $u\cdot \alpha_i=u\cdot c'+u\cdot e_iM=u\cdot c'+uM^T\cdot e_i$, and thus by \cite[Lemma 3.1]{Secondary} it holds that
$$((-1)^{u\cdot \alpha_0},\ldots, (-1)^{u\cdot \alpha_{2^m-1}})=(-1)^{u\cdot c'}((-1)^{uM^T\cdot e_0},\ldots,(-1)^{uM^T\cdot e_{2^m-1}})=(-1)^{\varepsilon_u}H^{(r_u)}_{2^m},$$
for some $0\leq r_u\leq 2^m-1$ and $\varepsilon_u\in \mathbb{F}_2$. The rest of the statement follows by \cite[Lemma 3.1]{Secondary}.\qed
%The second part follows trivially from the fact that $\dim(E)=m.$ \qed
\end{proof}
%
%\begin{rem}
%Note that in the case that $f$ is a trivial plateaued function constructed by Theorem \ref{th:eq} ($S_f,S_h$ are affine subspaces), then $\phi_b$ for all $b\in \mathbb{F}^n_2$ are linear/affine sequences (functions) and the inequality  $g \neq f^*+\phi_b+\varepsilon$ means that the duals $h^*=g$ and $f^*$ (defined via (\ref{eq:equiv22})) are not related in an affine manner.
%\end{rem}
%\begin{rem}\label{rem:bent}
%Note that for a trivial $s$-plateaued function $f:\mathbb{F}^n_2\rightarrow \mathbb{F}_2$, whose support is given as $S_f=v+E$, by \cite[Theorem %3.1-$(ii)$]{Secondary} we have that $\overline{f}^*$ is a bent function on $\mathbb{F}^{n-s}_2$. However, the dual $\overline{h}^*$ of a trivial $s$-plateaued function $h:\mathbb{F}^n_2\rightarrow \mathbb{F}_2$ is also bent, due to relation (\ref{eq:equiv22}).
%\end{rem}
%\begin{prop}\label{prop:hbent}
%Let $h:\mathbb{F}^n_2\rightarrow \mathbb{F}_2$ be a trivial $s$-plateaued function with $S_h=\{z_0,\ldots,z_{2^{n-s}-1}\}=c+S_fA^T$, where %$S_f=v+E=\{\omega_0,\ldots,\omega_{2^{n-s}-1}\}$ and furthermore $z_i=c+\omega_iA^T$, for some $c\in \mathbb{F}^n_2$ and $A\in GL(n,\mathbb{F}_2)$. %and $E$ is a %lexicographically ordered linear subspace.
%Then $h^*$ defined on $\mathbb{F}^{n-s}_2$ as $h^*(x_i)\leftrightsquigarrow h^*(z_i)$ is bent.
%%\begin{enumerate}[i)]
%%\item $h^*$ defined on $\mathbb{F}^{n-s}_2$ as $h^*(x_i)\leftrightsquigarrow h^*(z_i)$ is bent.
%%\item For a trivial $s$-plateaued function $\widehat{h}:\mathbb{F}^n_2\rightarrow \mathbb{F}_2$, let $S_{\widehat{h}}=\widehat{c}+S_fM=c+S_fA^T=S_h$ for some %$M\neq A$ $(M\in GL(n,\mathbb{F}_2))$ and $\widehat{c}\in \mathbb{F}^n_2$. Then the dual function $\widehat{h}^*(\widehat{z}_i)\leftrightsquigarrow %\widehat{h}^*(x_i)$ (defined on vectors $\widehat{z}_i=\widehat{c}+\omega_iM\in S_{\widehat{h}}$) is EA-equivalent to $h^*(x_i)\leftrightsquigarrow h^*(z_i)$.
%%\end{enumerate}
%\end{prop}
%\proof %The proof uses the same arguments as \cite[Corollary 1]{Secondary} (see also \cite[Remark 3.1]{Secondary}).
%For an arbitrary $u\in \mathbb{F}^n_2$, using  relation (\ref{WHT}) we have that %$W_h(u)=2^{\frac{n+s}{2}}(-1)^{h^*(u)}$ implies that
%$$(-1)^{h(u)}=2^{-n}\sum_{z_i\in S_h}W_h(z_i)(-1)^{u\cdot z_i}=2^{-\frac{n-s}{2}}\sum_{z_i\in S_h}(-1)^{h^*(z_i)+ u\cdot z_i},$$
%which is equivalent to
%$$\sum_{z_i\in S_h}(-1)^{h^*(z_i)+u\cdot z_i}=(-1)^{u\cdot c+uA\cdot v}\sum^{2^{n-s}-1}_{i=0}(-1)^{h^*(z_i)+uA\cdot e_i}=2^{\frac{n-s}{2}}(-1)^{h^*(u)}.$$
%%{\bf This part above needs to be expanded - I needed a few minutes to verify it and it is correct but do not forget that reviewers may ask questions here, for instance (10) is written in a nicer way ... Also, in the mid sum above you keep $z_i$ as the variable of $h^*$ but replacing $z_i$ otherwise ?????  I think you should extend any proof considerably, especially Corollary 3.1,  so that we can check whether they are correct and not speculate about it - then later we can shorten these !! Put some question marks to any statement if you are not 100\% sure !!!! \\ Maybe we should be more formal about this $h^*(x_i)\leftrightsquigarrow h^*(z_i)$ and introduce a mapping $P$ which takes care of correspondence between subset and $\FB^{n-s}$ ?? The same remark is valid as for our CF paper- we do not know its explicit definition - maybe in case of affine subspaces !} \\
%%{\bf Isn't it $2^{\frac{n+s}{2}}(-1)^{h^*(u)}$ in the last term above ??}\\
%By \cite[Lemma 3.1]{Secondary} (or Proposition \ref{prop:Hrow}), we have that $\{(uA\cdot e_0,\ldots,uA\cdot e_{2^{n-s}-1}):v\in \mathbb{F}^{n}_2,\;e_i\in E\}$ ($E$ ordered lexicographically) contains truth tables of all linear functions defined on $\mathbb{F}^{n-s}_2$, i.e., $uA\cdot e_i$ corresponds to the values of a linear function $\ell(x_i)$ ($\ell\in \mathcal{L}_{n-s}$).
%Then the equality above combined with $h^*(x_i)\leftrightsquigarrow h^*(z_i)$ means that $h^*$ is at bent distance to all linear functions in $n-s$ variables. In other words, $h^*$ is bent.\qed
%
%\textbf{Fengrong: I feel we can delete Proposition \ref{prop:Hrow} since from Theorem 3-ii) we have know the result. }
%
%{\bf Enes: Can we indeed do that ?? Some results rely on this Prop. 3.1 but maybe it is not essential here ?? We need to get a well structured manuscript avoiding too many small results !!}
%%%%%%%%%%%%%%%%%%%%%%%%%%%%%%%%%%%%%%%%%%%%%%%%%%
%\begin{prop}\label{prop:hbent2}
%Let $h:\mathbb{F}^n_2\rightarrow \mathbb{F}_2$ be a trivial $s$-plateaued function with $S_h=\{z_0,\ldots,z_{2^{n-s}-1}\}$ $=c+S_fA^T$, for some $c\in \mathbb{F}^n_2$ and $A\in GL(n,\mathbb{F}_2)$ $(z_i=c+\omega_iA^T,$ $\omega_i\in S_f)$. For a trivial $s$-plateaued function $\widehat{h}:\mathbb{F}^n_2\rightarrow \mathbb{F}_2$, let $S_{\widehat{h}}=\widehat{c}+S_fM=c+S_fA^T=S_h$ for some $M\neq A$ $(M\in GL(n,\mathbb{F}_2))$ and $\widehat{c}\in \mathbb{F}^n_2$. Then the dual function $\widehat{h}^*(\widehat{z}_i)\leftrightsquigarrow \widehat{h}^*(x_i)$ (defined on vectors $\widehat{z}_i=\widehat{c}+\omega_iM\in S_{\widehat{h}}$) is EA-equivalent to $h^*(x_i)\leftrightsquigarrow h^*(z_i)$.
%\end{prop}
%\proof For a matrix $K\in GL(n,\mathbb{F}_2)$ with $A^TK=M$ and a (suitable) vector $\lambda\in \mathbb{F}^n_2$ we have  $$\lambda+z_iK=\lambda+(c+\omega_iA^T)K=(\lambda+cK)+\omega_iM=\widehat{c}+\omega_iM=\widehat{z}_i,$$ i.e.,  $\widehat{z}_i=\lambda+z_iK$ are related in an affine manner. Consequently, $\widehat{h}^*(\widehat{z}_i)=h^*(L(z_i))$ ($L$ is an affine mapping on $S_h$), i.e., $\widehat{h}^*$ and $h^*$ are affine equivalent.\qed
%
%Due to the fact for an affine subspace $S_f$ it holds that the sequence profile $\Phi_f$ contains only linear/affine functions (Proposition \ref{prop:Hrow}), then  Theorem \ref{theo:EA} implies the following two results.
%{\bf This statement below needs to be heavily rewritten !!! Define the binary representation outside (maybe we already have it defined) and just use this in the Proposition !!}

%Before we provide a technical result, we firstly set some preparatory notation. %By $(i)_2$ ($i\in[0,2^{m}-1]$), we denote the binary representation of length $m$ of an integer $i$.
%
Suppose that a linear subspace $E=\{e_0,\ldots,e_{2^{n-s}-1}\}\subset {\Bbb F}_2^n$ ($\dim(E)=n-s$) and ${\Bbb F}_2^{n-s}=\{x_0,x_1,x_2,\ldots,x_{2^{n-s}-1}\}$ are ordered lexicographically.  From the proof of Lemma 3.1-$(ii)$ given in \cite{Secondary}, the linear mapping $\Theta(x_i)=e_i$, with $\Theta: {\Bbb F}_2^{n-s} \rightarrow E$,  can be described via a matrix $R=R_{(n-s)\times n}$ so that
\begin{eqnarray}\label{eq:Q}
\Theta(x_i)=x_iR=x_i\left(
                                                                                        \begin{array}{c}
                                                                                          e_{2^{n-s-1}} \\
                                                                                          \vdots \\
                                                                                          e_2 \\
                                                                                          e_1 \\
                                                                                        \end{array}
                                                                                      \right)=e_i,\;\;\;\;i\in[0,2^{n-s}-1],
\end{eqnarray}
where the rows of  $R$, $e_{2^j}$ for $j =0,\ldots, n-s-1$, are basis vectors of $E$. We have the following technical result useful for establishing  EA-equivalence of plateaued functions.
%
%let their basis respectively be denoted as $\mathbf{B}_{{\Bbb F}_2^{n-s}}=(x_1,x_2,x_4,\ldots,x_{2^{n-s-1}})$ (as the canonical basis) and $\mathbf{B}_E=(e_1,e_2,e_4,\ldots,e_{2^{n-s-1}})$ (see \cite[Lemma 3.1-$(i)$]{Secondary}). In terms of these bases, the sets $E$ and ${\Bbb F}_2^{n-s}$ can be written as
%\begin{eqnarray}\label{eq:E}\nonumber
%E&=&\{0_n,(1)_2\cdot \mathbf{B}_E,\ldots, (2^{n-s}-1)_2\cdot \mathbf{B}_E\}=\{(i)_2\cdot \mathbf{B}_E: i\in[0,2^{n-s}-1]\},\\
%{\Bbb F}_2^{n-s}&=&\{(i)_2\cdot \mathbf{B}_{{\Bbb F}_2^{n-s}}: i\in[0,2^{n-s}-1]\}.
%\end{eqnarray}
%To clarify the notation, let $n-s=3$. Then the previous notation means that for $(3)_2=(0,1,1)$ and basis $\mathbf{B}_E=(\alpha^{(1)},\alpha^{(2)},\alpha^{(3)})$ we have that %$(3)_2\cdot \mathbf{B}_E=\alpha^{(2)}+\alpha^{(3)}.$ The following result concerns the existence of a mapping ${\Bbb F}_2^{n-s} \rightarrow S$, where $S=\gamma +E$.
%%Similarly, for a basis $\Omega_{{\Bbb F}_2^{n-s}}=(\beta^{(1)},\beta^{(2)},\ldots,\beta^{(n-s)})$ of ${\Bbb F}_2^{n-s}$, the space ${\Bbb F}_2^{n-s}$ we order
%%as ${\Bbb F}_2^{n-s}=\{(i)_2\cdot (\beta^{(1)},\beta^{(2)},\ldots,\beta^{(n-s)}): i\in[0,2^{n-s}-1]\}$.
\begin{lemma}\label{prop:Hrow1}
Let a linear subspace $E\subset {\Bbb F}_2^n$ and ${\Bbb F}_2^{n-s}$ be lexicographically ordered, and let the mapping $\Theta:{\Bbb F}_2^{n-s} \rightarrow E$ be defined by (\ref{eq:Q}). Then for an arbitrary (fixed) matrix $B\in GL(n-s,\mathbb{F}_2)$ and vector $t\in \mathbb{F}^{n-s}_2$, there exists a vector $\gamma\in E$ and a matrix $D\in GL(n,\mathbb{F}_2)$ which is invariant on $E$ (that is $E=ED$) such that
 $$\Theta(x_iB+t)=e_iD+\gamma,$$
 holds for all $i\in[0,2^{n-s}-1]$.
\end{lemma}
\begin{proof}
Since the rows of  $R$ are basis vectors $e_1,e_2,e_4,\ldots,e_{2^{n-s-1}}\in E$, then clearly we can always find a set of vectors $\lambda_1,\ldots,\lambda_s\in \mathbb{F}^n_2$ such that $\{e_1,\ldots,e_{2^{n-s-1}},\lambda_1,\ldots,\lambda_s\}$ is a basis of $\mathbb{F}^n_2$. Now, let $\Lambda=\Lambda_{s\times n}=(\lambda_1,\ldots,\lambda_s)^T$ denotes the matrix whose rows are $\lambda_i$ vectors. Then, for a matrix $B\in GL(n-s,\mathbb{F}_2)$ and an arbitrary $x_i\in \mathbb{F}^{n-s}_2$ it holds that
\begin{eqnarray}\label{eq:ED}
\Theta(x_iB)=(x_iB)R=(x_i,\textbf{0}_s)\left(
                                           \begin{array}{c}
                                             BR \\
                                             \Lambda \\
                                           \end{array}
                                         \right)_{n\times n}=(x_i,\textbf{0}_s)Q=e_i\in E,
\end{eqnarray}
since $\tilde{x}_i=x_iB\in \mathbb{F}^{n-s}_2$ and $\Theta(\mathbb{F}^{n-s}_2)\in E$. Here, the block matrix $Q=\left(
                                           \begin{array}{c}
                                             BR \\
                                             \Lambda \\
                                           \end{array}
                                         \right)_{n\times n}$ is clearly invertible, i.e., $Q\in GL(n,\mathbb{F}_2)$, since the matrix $BR$ (of size $(n-s)\times n$) contains rows which are linear combinations of vectors $e_{2^j}$, $j\in[1,n-s-1]$ (which are linearly independent of $\lambda_i$).
Moreover, denoting by $U=\left(
                                           \begin{array}{c}
                                             R \\
                                             \Lambda \\
                                           \end{array}
                                         \right)_{n\times n}\in GL(n,\mathbb{F}_2)$, we have that
$$\Theta(x_i)=x_iR=(x_i,\textbf{0}_s)U=e_i,$$ implies that $(x_i,\textbf{0}_s)=e_iU^{-1}$ holds for all $i\in[0,2^{n-s}-1].$ Consequently, for an arbitrary vector $t\in \mathbb{F}^{n-s}_2$ we have that
\begin{eqnarray*}
\Theta(x_iB+t)=(x_iB+t)R=x_i(BR)+tR=(x_i,\textbf{0}_s)Q+tR=e_iU^{-1}Q+tR=e_iD+\gamma\in E,
\end{eqnarray*}
where $D=U^{-1}Q\in GL(n,\mathbb{F}_2)$ and $tR=\gamma\in E$ due to the facts that $\Theta(x_iB)=x_i(BR)=e_iD\in E$ holds by (\ref{eq:ED}). Also, $\Theta(x_iB+t)\in \Theta(\mathbb{F}^{n-s}_2)=E$ since $E$ is a linear subspace.\qed
%%%%%%%%%%%%%%%%%
%\begin{rem}
%In definition of the mapping $\phi$ in Proposition \ref{prop:Hrow1}, the vectors $t$ and $\widehat{c}$ are neglected. Thus we only consider the vector $(i)_2$ which is present %on both sides, where the basis $\mathbf{B}_E$ and $\mathbf{B}_{\mathbb{F}^{n-s}_2}$ are fixed. In the case that these bases are not fixed, then for instance one can easily find another basis in $E$, say $\widehat{\mathbf{B}}_E$, such that $\phi((i)_2\cdot \mathbf{B}_{\mathbb{F}^{n-s}_2}+t)=(i)_2\cdot \widehat{\mathbf{B}}_E+\gamma,$ and thus the mapping $\phi$ is not well defined (since it can map one input to two different outputs).
%\end{rem}
%%%%%%%%%%%%%%%%%%
%Regarding a trivial plateaued function (say $h:\mathbb{F}^{n}_2\rightarrow \mathbb{F}_2$), it is important to emphasize that its dual $h^*$ (defined by (\ref{DPL})) is a bent function if and only if the linear space $E\subset \mathbb{F}^n_2$ (in representation $S_h=v+EM$) is ordered as in (\ref{eq:E}) (due to Proposition \ref{prop:Hrow}). In other words, as long as our ordering of the Walsh support $S_h=\{z_0,\ldots,z_{2^{n-s}-1}\}$ (when written as a matrix of size $2^{n-s}\times n$) has columns which correspond to truth tables of affine/linear functions in $n-s$ variables, then the corresponding values of the dual $h^*$ (when taken to be as values on $\mathbb{F}^{n-s}_2$ ordered as in \ref{eq:E})
%%%%%%%%%%%%%%%%%%%%%%%%%%%%%%%%%%%%%%%%%%%%%%%%%%%%%%%%%%%%%%%%%%%%%%%
%Now a matrix $A\in GL(n,\mathbb{F}_2)$ be chosen such that
%$$\alpha^{(j)}A+\gamma= \alpha^{(i_{j})},\;\;\;j\in[1,n-s],$$
%for some vector $\gamma\in \mathbb{F}^n_2.$  Now, let $(\alpha^{(1)},\alpha^{(2)},\ldots,\alpha^{(n-s)},\alpha^{(n-s+1)}, \alpha^{(n)})$ be a basis of ${\Bbb %F}_2^n$. Note that this basis contains the vectors of $\Omega_E$. and for $j=1,2,\ldots,n$ and $i_j=j $ when $j>n-s$ let
% \begin{equation}\label{equ matrix}
%  \alpha^{(j)}A+\gamma= \alpha^{(i_{j})}.
%  \end{equation}
% % where $j=1,2,\ldots,n$ and $i_j=j $ when $j>n-s$.
%  Now if we can show that
%  $(\alpha^{(i_{1})},\alpha^{(i_{2})},\ldots,\alpha^{(i_{n})})$ (that is, $(\alpha^{(i_{1})},\alpha^{(i_{2})},\ldots,\alpha^{(i_{{n-s}})},\alpha^{({{n-s+1}})}, %\ldots, \alpha^{({{n}})}  )$) is a basis of ${\Bbb F}_2^n$, then
%  the matrix $A$ can be obtained.
%  We also know $(\alpha^{(i_1)},\alpha^{(i_2)},\ldots,\alpha^{(i_{n-s)}})$ is  a basis of $E$. Hence,   $$(\alpha^{(i_{1})},\alpha^{(i_{2})},\ldots,\alpha^{(i_{{n-s}})},\alpha^{({{n-s+1}})}, \ldots, \alpha^{({{n}})}  )$$ is a basis of ${\Bbb F}_2^n$.
\end{proof}
%%%%%%%%%%%%%%%%%%%%%%%
Recall that any two EA-equivalent trivial plateaued functions $f$ and $h$ have the Walsh supports related in an affine manner as $S_h=c+S_fM$ (for some invertible matrix $M$ and vector $c$). If there do not exist $M$ and $c$ such that $S_h=c+S_fM$ holds, then by Theorem \ref{th:eq}-$(i)$ (or relation (\ref{eq:equiv})) the functions $f$ and $h$ are not equivalent. In other words, the EA-equivalence between $f$ and $h$ is reasonable to consider only under the assumption that $S_h$ and $S_f$ are related in an affine manner.
The following result characterizes the EA-equivalence between trivial plateaued functions.
\begin{theo}\label{th:EA1}
Let $f,h:\mathbb{F}^n_2\rightarrow \mathbb{F}_2$ be two trivial $s$-plateaued functions whose Walsh supports are related as $S_h=c+S_fM$, for some matrix $M\in GL(n,\mathbb{F}_2)$ and $c\in \mathbb{F}^n_2.$ Representing $S_f=v+E=\{\omega_i=v+e_i:e_i\in E\}$ for a lexicographically ordered linear space $E=\{e_0,\ldots,e_{2^{n-s}-1}\}$, let the functions $\overline{f}^*$ and $\overline{h}^*$ be defined as
$$\overline{f}^*(x_i)=f^*(\omega_i)\;\;\;\text{and}\;\;\;\overline{h}^*(x_i)=h^*(z_i),\;\;\;(i\in[0,2^{n-s}-1]),$$
where $z_i=c+\omega_iM\in S_h$. Then $f$ and $h$ are EA-equivalent if and only if their duals $\overline{f}^*,\overline{h}^*:\mathbb{F}^{n-s}_2\rightarrow\mathbb{F}_2$ %(defined as functions on $\mathbb{F}^{n-s}_2$ via (\ref{eq:equiv22}) by setting $z_i\leftrightsquigarrow \omega_i\leftrightsquigarrow  x_i$, $x_i\in \mathbb{F}^{n-s}_2$)
are EA-equivalent bent functions. %, where for lexicographically ordered linear space $E=\{e_0,\ldots,e_{2^{n-s}-1}\}$ and $S_f=v+E=\{\omega_i=v+e_i:e_i\in E\}$, the functions $\overline{f}^*$ and $\overline{h}^*$ are defined as $\overline{f}^*(x_i)=f^*(\omega_i)$ and $\overline{h}^*(x_i)=h^*(c+\omega_iM)$ $(i\in[0,2^{n-s}-1]).$
\end{theo}
\begin{proof} $(\Rightarrow)$ By \cite[Theorem 3.1-$(ii)$]{Secondary} and relations (\ref{eq:identif})-(\ref{eq:equiv22}) it is clear that $\overline{f}^*,\overline{h}^*:\mathbb{F}^{n-s}_2\rightarrow\mathbb{F}_2$  are EA-equivalent bent functions.

$(\Leftarrow)$ In this part, we assume that $f$ and $h$ are two fixed trivial plateaued functions and thus the values of $h^*$ (on $S_h$) and $f^*$ (on $S_f$) are fixed in advance.
%
%Without loss of generality, suppose that $S_f$ is given as $S_f=v+E$, with lexicographically ordered $E= \{e_0,\ldots,e_{2^{n-s}-1}\}$. Note that \cite[Lemma 3.1-$(i)$]{Secondary} implies the fact that $E$ is ordered as in (\ref{eq:E}) with the basis set $\mathbf{B}_E=(e_1,e_2,e_4,\ldots,e_{2^{n-s-1}-1}).$ With respect to $S_f=v+E$, the dual $f^*:\mathbb{F}^{n-s}_2\rightarrow \mathbb{F}_2$ is defined as
%$$\overline{f}^*(x_i)=f^*(v+e_i),\;\;\;x_i\in \mathbb{F}^{n-s}_2,,\;e_i\in S_h.$$
%Note that $\overline{f}^*$ is bent due to \cite[Theorem 3.1-$(ii)$]{Secondary}. Furthermore, since the Walsh supports $S_h$ and $S_f$ are affine subspaces in $\mathbb{F}^{n}_2$, then there exist $M\in GL(n,\mathbb{F}_2)$ and $\gamma\in \mathbb{F}^n_2$ such that $S_h=\gamma+EM$ (i.e., $S_h$ and $S_f$ are related in an affine manner). Note that $S_h$ is of the form $S_h=\gamma+\widetilde{E}$ with $\widetilde{E}=EM$, and $\widetilde{E}=(e_0M,e_1M,\ldots,e_{2^{n-s}-1}M)$ is ordered as in (\ref{eq:E}) for basis $\mathbf{B}_{\widetilde{E}}=\mathbf{B}_E M$.  According to this relation, we denote by $z_i=\gamma+e_iM$, for $i\in[0,2^{n-s}-1]$, and define $\overline{h}^*:\mathbb{F}^{n-s}_2\rightarrow \mathbb{F}_2$ as
%$$\overline{h}^*(x_i)=h^*(z_i),\;\;\;x_i\in \mathbb{F}^{n-s}_2,\;z_i\in S_h.$$
%%%%%%%%%%%%%%%%%%%%
Due to Proposition \ref{prop:Hrow} and Theorem \ref{theo:plateH}-$(ii)$ (which is given later on in Section \ref{sec:ineq}), it is not difficult to see that both duals $\overline{h}^*$ and $\overline{f}^*$ are bent functions.

Now, let us assume that $\overline{f}^*,\overline{h}^*$  are affine equivalent, that is for some $B\in GL(n-s,\mathbb{F}_2)$, $t,r\in \mathbb{F}^{n-s}_2$ and $\kappa\in \mathbb{F}_2$, we have
\begin{eqnarray}\label{eq:w0}
\overline{h}^*(x_i)=\overline{f}^*(x_iB+t)+r\cdot x_i+\kappa, \;\;\;\;\; x_i\in \mathbb{F}^{n-s}_2.
\end{eqnarray}
%In order to proceed further, we consider the following steps:\\\\
%\textbf{Step 1:} Following the notation of Proposition \ref{prop:Hrow1}, let us represent $E=ED+\gamma$ ($\gamma\in E$). Taking a vector $p\in \mathbb{F}^n_2$ such that $v=\gamma D^{-1}+p$, let us define the  mapping $\Psi:E\rightarrow S_f$ as $\Psi(e)=eD^{-1}+p.$ For an arbitrary vector $e_iD+\gamma\in E$ we have
%$$\Psi(e_iD+\gamma)=(e_iD+\gamma)D^{-1}+p=e_i+(\gamma D^{-1}+p)=e_i+v\in v+E=S_f.$$
%Now, since the mapping $\Theta:\mathbb{F}^{n-s}_2\rightarrow E=ED+\gamma$ and $\Psi:ED+\gamma\rightarrow S_f$, and both are affine mappings, then their composition is an affine mapping, and it maps $\mathbb{F}^{n-s}_2$ to $S_f$, i.e., we have that $\Psi \circ \Theta:\mathbb{F}^{n-s}_2\rightarrow S_f$ such that $(\Psi\circ \Theta)(x_iB+t)=v+e_i,$ $i\in[0,2^{n-s}-1].$
%
%Consequently, the value of $\overline{f}^*$ can be symbolically identified with $f^*=\overline{f}^*\circ\Psi\circ \Theta$, i.e., we have
%\begin{eqnarray}\label{eq:f*}
%\overline{f}^*(x_iB+t)=\overline{f}^*((\Psi\circ \Theta)(x_iB+t))=f^*(v+e_i)=f^*(\omega_i),\;\;\;i\in[0,2^{n-s}-1].
%\end{eqnarray}
%
%\textbf{Step 2:}
Since $S_f$ is an affine subspace, then by Proposition \ref{prop:Hrow1} there exist  $b\in \mathbb{F}^n_2$ and $\varepsilon\in \mathbb{F}_2$ such that
\begin{eqnarray}\label{eq:w1}
b\cdot \omega_i+\varepsilon=r\cdot x_i+\kappa,
\end{eqnarray}
where $(b\cdot e_0,\ldots,b\cdot e_{2^{n-s}-1})=(r\cdot x_0,\ldots,r\cdot x_{2^{n-s}-1})$ ($x_i\in \mathbb{F}^{n-s}_2$) is truth table of a linear function in $n-s$ variables, and $b\cdot v+\varepsilon=\kappa$.

%%%%%
%%%%%
Furthermore, since by Lemma \ref{prop:Hrow1} we have that $\Theta(x_iB+t)=e_iD+\gamma\in E$, then clearly we can identify vectors $x_iB+t\in \mathbb{F}^{n-s}_2$ with vectors $v+(e_iD+\gamma)\in v+E=S_f$, and thus we have that
\begin{eqnarray}\label{eq:w2}
\overline{f}^*(x_iB+t)=f^*(v+(e_iD+\gamma)),\;\;\;\;i\in[0,2^{n-s}-1].
\end{eqnarray}
This is possible due to the fact that the definition of $\overline{f}^*$ on $\mathbb{F}^{n-s}_2$ (given by (\ref{DPL})) is strictly given with respect to vectors $e_i\in E$ in $f^*(v+e_i)$, where the vector $v$ is ignored.
%
 %Moreover, by $\omega_i=v+e_i$ we have that $v+(e_iD+\gamma)=v+\gamma+vD+\omega_iD$, and using the substitution $t_i=v+\gamma+vD+\omega_iD\in S_f$ we can write
%$$\omega_i=(t_i+v+\gamma+vD)D^{-1}.$$
Combining the previous computation with Lemma \ref{prop:Hrow1} and relations (\ref{eq:w1})-(\ref{eq:w2}), we have that
\begin{eqnarray}\label{eq:w3}\nonumber
\overline{f}^*(x_iB+t)+r\cdot x_i+\kappa \hskip -2mm &=&\hskip -2mm f^*(v+(e_iD+\gamma))+b\cdot \omega_i+\varepsilon=f^*(v+\gamma+vD+\omega_iD)+b\cdot \omega_i+\varepsilon\\
&=&f^*(t_i)+b\cdot (t_i+v+\gamma+vD)D^{-1}+\varepsilon,
\end{eqnarray}
where we have used the substitution $t_i=v+\gamma+vD+\omega_iD\in S_f$ which gives that $\omega_i=(t_i+v+\gamma+vD)D^{-1}.$

Consequently, by relations (\ref{eq:w0})-(\ref{eq:w3}), for arbitrary $x\in \mathbb{F}^n_2$ it holds that
 \begin{equation*}%\label{equ IWHT}
 \begin{array}{rl}
 (-1)^{h(x)}=&2^{-n}\sum\limits_{z_i\in S_h}W_h(z_i)(-1)^{x\cdot z_i}=2^{\frac{s-n}{2}}\sum\limits_{z_i\in S_h}(-1)^{h^*(z_i)+x\cdot z_i}\\
=&2^{\frac{s-n}{2}}\sum\limits_{e_i\in E}(-1)^{f^*(v+e_iD+\gamma)+b\cdot (t_i+v+\gamma+vD)D^{-1}+x\cdot (c+vM+e_iM)+\varepsilon}\\
\stackrel{t_i=v+e_iD+\gamma}{=}&2^{\frac{s-n}{2}}\sum\limits_{t_i\in S_f}(-1)^{f^*(t_i)+b\cdot (t_i+v+\gamma+vD)D^{-1}+x\cdot [c+(t_i+v+\gamma+vD)D^{-1}M]+\varepsilon}\\
=&2^{-n}\sum\limits_{t_i\in S_f}W_f(t_i)(-1)^{t_i\cdot [bD^{-T}+x(D^{-1}M)^{T}]+ b\cdot (v+\gamma+vD)D^{-1}+ x\cdot [c+(v+\gamma+vD)D^{-1}M] +\varepsilon}\\
{=}& (-1)^{f(bD^{-T}+x(D^{-1}M)^{T})+x\cdot [c+(v+\gamma+vD)D^{-1}M]+[b\cdot (v+\gamma+vD)D^{-1} +\varepsilon]},\\
 \end{array}
 \end{equation*}
which means that $f$ and $h$ are affine equivalent.\qed

\begin{rem}\label{rem:ordering}
Note that Theorem \ref{th:EA1} also holds  when the set $E$ is ordered as $E=\{\hat{e}_0,\hat{e}_1,\ldots,\hat{e}_{2^{n-s}-1}\}$, where the vectors $\hat{e}_j$, for any $i\in[0,n-s-1]$, satisfy the recursion $\hat{e}_j=\hat{e}_{2^i}+\hat{e}_{j-2^i}$ for all $2^i\leq j\leq 2^{i+1},$ in which case the set $\{\hat{e}_{2^i}:i\in[1,n-s-1]\}$ is a basis of $E$. Regarding the equality $S_h=c+S_fM$, it is clear that the corresponding ordering of $S_h=\{\hat{z}_0,\ldots,\hat{z}_{2^{n-s}-1}\}$ with $\hat{z}_j=c+(v+\hat{e}_j)M$ (where $v+\hat{e}_j\in S_f$) satisfies the same recursion $\hat{z}_j=\hat{z}_{2^i}+\hat{z}_{j-2^i}$. Thus, by Lemma 3.1-$(i)$ and Theorem 3.1 given in \cite{Secondary} (or Theorem \ref{theo:plateH}),  both duals $\overline{f}^*(x_i)=f^*(v+\hat{e}_j)$ and $\overline{h}^*(x_i)=f^*(\hat{z}_j)$ are still bent functions on $\mathbb{F}^{n-s}_2$. Another proof in this context (without using a restriction on lexicographic ordering of $E$) is given in Appendix.
\end{rem}
%%%%%%%%%%%%%%%%%%%%%%%%%%%%%%%%%%%%%%%%%%%%%%%%%%%%%5
%%%%%%%%%%%%%%%%%%%%%%%%%%%%%%%%%%%%%%%%%%%%%%%%%%%%%
Note that second part of the proof of Theorem \ref{th:EA1} embeds a spectral equivalence which is used in Section \ref{sec:ineq} (cf. Theorem \ref{theo:plateH}) to provide an efficient method of constructing inequivalent trivial plateaued functions through suitably  specified duals.
%%%%%%%%%%%%%%%%%%%%%%%
%From Proposition \ref{prop:Hrow1}, there clearly exist a matrix $A\in GL(n,\mathbb{F}^n_2)$, $b,c\in \mathbb{F}^n_2$ and $\varepsilon \in \mathbb{F}_2$ such that for $h^*(z_i)=h^*(c+\omega_iA^T)=\overline{h}^*(y_iB^{-1}+tB^{-1})=\overline{h}^*(x_i)$ it holds that $$h^*(z_i)=f^*(\omega_i)+b\cdot z_i+\varepsilon=f^*(\omega_i)+\phi_b(z_i)+\varepsilon,$$ as a consequence of Proposition \ref{prop:Hrow} and Remark \ref{rem:ordering}. This simply means that the affine transformation $\mathbb{F}^{n-s}_2B^{-1}+tB^{-1}$ of the space $\mathbb{F}^{n-s}_2$ (in terms of indices) can be identified with $c+S_fA^T=S_h$ for some matrix $A$ and vector $c$.
%
%%Consequently, by relations (\ref{eq:equiv})-(\ref{eq:equiv22}), %(using $\phi_b\leftrightsquigarrow \ell\in \mathcal{A}_{n-s}$ due to Proposition \ref{prop:Hrow}),
%this implies the equivalence between $f$ and $h$.\qed
\end{proof}
%\textbf{\textcolor[rgb]{0.98,0.00,0.00}{Samir:} Fengrong proposed the proof with the Carlet's result on partially bent functions. Here is my version of it:}\\\\
%\proof Since $f$ is a trivial plateaued function, then by \cite[Section 2]{CarletPartial} there exist two linear subspaces $G,D\subset \mathbb{F}^n_2$, vector $t\in \mathbb{F}^n_2$ and a bent function $\xi:R\rightarrow \mathbb{F}_2$ such that $G\oplus D=\mathbb{F}^n_2$ ($G\cap D=\{\textbf{0}_n\}$) and
%$$f(x+y)=\xi(x)+t\cdot y,\;\;\;\;x\in G,\;y\in D.$$
%Thus, for arbitrary $u\in \mathbb{F}^n_2$ it is not difficult to see that
%$$W_f(u)=\left\{\begin{array}{cc}
%                  0, & u\not\in t+D^{\perp} \\
%                  2^{\frac{n+s}{2}}(-1)^{\xi^*(u)}, &  u\in t+D^{\perp}=S_f
%                \end{array}
%\right..$$
%Clearly, here we have that $f^*=\xi^*$. Similarly, by relation (\ref{eq:EAeq}) for $z\in \mathbb{F}^n_2$ we have that
%$$W_h(z)=\left\{\begin{array}{cc}
%                  0, & z\not\in S_f \\
%                  2^{\frac{n+s}{2}}(-1)^{h^*(z)}=2^{\frac{n+s}{2}}(-1)^{\xi^*((c+z)A^{-T})+(c+z)\cdot bA^{-1}+\varepsilon}, &  (c+z)A^{-T}\in S_f
%                \end{array}
%\right..$$
%The previous computation implies that
%$$h^*(z_i)=h^*(c+\omega_iA^T)=\xi^*(\omega_i)+b\cdot\omega_i+\varepsilon,\;\;\;\;i\in[0,2^{n-s}-1],$$
%and since by Proposition \ref{prop:Hrow} we have that $b\cdot\omega_i +\varepsilon\leftrightsquigarrow \ell(x_i)$ with $\ell\in \mathcal{A}_{n-s}$, then by $f^*=\xi^*$ we have that $h^*(z_i)=h^*(c+\omega_iA^T)=f^*(\omega_i)+\ell(x_i).$....\\\\
%\textbf{\textcolor[rgb]{0.98,0.00,0.00}{Samir:} I stopped here, due to the following reasons:
%\begin{itemize}
%\item First, notice that the use of the Carlet's result on partially bent functions provides the proof that looks the same as the already given earlier. Everything narrows down to the same equality, that is $h^*(z_i)=h^*(c+\omega_iA^T)=f^*(\omega_i)+b\cdot\omega_i+\varepsilon$.
%\item Notice first that even having a bent function $\xi$ (from Carlet's result) we do not avoid "$\leftrightsquigarrow$", since $\xi$ is considered to be defined on the set $G\subset \mathbb{F}^n_2$, and thus saying that $\xi$ is a bent function $\mathbb{F}^{n-s}_2$, we need to identify $G$ with $\mathbb{F}^{n-s}_2$.
%\item On the picture provided by Fengrong, it is written that $f^{(b)}$ is a function in variable $x^{(1)}\in \mathbb{F}^{n-s}_2$, but from Carlet's result it is evident that we are speaking about the function $\xi$ (i.e., its dual), and it is defined on $G$ which is  a subspace of $\mathbb{F}^n_2$!
%\item Why we even consider "viewing $\xi$ as a function on $\mathbb{F}^{n-s}_2$" -- this is due to the statement of Theorem 3.2, which says that $f^*$ and $h^*$ are two equivalent \textcolor[rgb]{0.98,0.00,0.00}{bent functions on $\mathbb{F}^{n-s}_2$}. Now, in this part one has to know that considering a function $\xi:G\rightarrow \mathbb{F}_2$ as a function from $\mathbb{F}^{n-s}_2$ to $\mathbb{F}_2$ is NOT something trivial, and one has to consider "some ordering" and Proposition 3.1. I hope that we all understand the importance of orderings in our article, and how it affects the whole story.
%\item Considering the equality, $h^*(z_i)=h^*(c+\omega_iA^T)=f^*(\omega_i)+b\cdot\omega_i+\varepsilon$ -- here one can not skip the story about the transition from Walsh supports $S_h$ or $S_f$ to $\mathbb{F}^{n-s}_2$ in order to say that \textcolor[rgb]{0.98,0.00,0.00}{$h^*$ and $f^*$ are bent functions on $\mathbb{F}^{n-s}_2$}. This requires the identification "$\leftrightsquigarrow$" (that is relation (\ref{DPL})), and in the proof that I provided earlier, we simply have to discuss all details regarding the "transition from $S_h$ or $S_f$ to $\mathbb{F}^{n-s}_2$". This means that all details related to this part that are provided in my proof -- have actually to be repeated in the "new" proof above...Hence, we can not skip the details given in the old proof.
%\item Just one more comment. Imagine that in $h^*(z_i)=h^*(c+\omega_iA^T)=f^*(\omega_i)+b\cdot\omega_i+\varepsilon$ we have that
%$$(b\cdot \omega_0,\ldots,b\cdot\omega_{2^{n-s}-1})$$
%is not truth table of any linear function. What can we say about equivalence of $h^*$ and $f^*$ in that case?? Nothing....since it does not fit to definition of EA-equivalence -- and in that case the main problem is the ordering of $S_f=\{\omega_0,\ldots,\omega_{2^{n-s}-1}\}$, and therefore one has to have proper orderings in relation (\ref{DPL}). In general, various orderings can work, but Proposition 3.1 has to be there always ($S_h$ and $S_f$ are affine subspaces), since we need to have that $b\cdot\omega_i+\varepsilon$ is a linear function over indices $i\in[0,2^{n-s}-1].$
%\item FINALLY: The statement of Theorem 3.2 and its proof are given from the aspect of spectral method and Theorem 3.3 given later on. I was always aware of the Carlet's result on partially bent functions, but as far as Theorem 3.3 is concerned, I find it more suitable for proving Theorem 3.2, due to the fact that Theorem 3.2 is considering duals $f^*$ and $h^*$ as functions on $\mathbb{F}^{n-s}_2$ (but not on some subspaces of $\mathbb{F}^n_2$). And what is the main reason behind this approach in general? It is the following fact. Imagine that I say that the dual of a trivial plateaued function is bent, say $f^*=g:\mathbb{F}^{n-s}_2\rightarrow \mathbb{F}_2$). Do you know how to construct a plaeaued function from it? No...it is not possible except in the case that you know which value $g(x_i)$ has to be considered as $g(\omega_i)$ ($\omega_i\in S_f\subset \mathbb{F}^n_2$) in the spectrum, but that requires proper ordering of the Walsh support $S_f$, since not all orderings of $S_f$ will work in general...
%\end{itemize}
%}


%%%%%%%%%%%%%%%%%
%\textbf{\textcolor[rgb]{0.98,0.00,0.00}{More comments:}\\\\
%1) Let $L:\mathbb{F}^{n-s}_2\rightarrow \mathbb{F}^{n-s}_2$ be an arbitrary affine mapping (it is bijection on $\mathbb{F}^{n-s}_2$) with $L:y_i\rightarrow d_j$. Does there always exist an affine mapping on $\mathbb{F}^n_2$ that it sends $\omega_i\leftrightsquigarrow y_i$ to some $z_i\leftrightsquigarrow d_j$ (so, this mapping will send $\omega_i$ to $z_i$ in the same way as $L$ sends $y_i$ to $d_j$ in terms of indices!!)? In my opinion this holds, and I do not see reason why it should not...so this gives the possibility of the blue text above (in the proof) to be always true, even if we fix $S_h$ and $S_f$ in advance!!!!!\\\\
% 2) Is there anything that may go wrong in the part above? We fix the ordering on $S_f$ and we are looking for a vector $c$ and matrix $A$ that will provide $z_i=c+\omega_iA^T$ in the same way (IN TERMS OF INDICES) as the mapping $L$ sends $\omega_i\leftrightsquigarrow y_i$ to some $z_i\leftrightsquigarrow d_j$....so still there is not problem here in my view...}\\\\
%%%%%%%%%%%
%{\bf The above proof is quite unclear in terms of this relation $\leftrightsquigarrow$ !! We need to be more clear here since nobody will understand this !!} \\
%We have the following result.
%\begin{cor}\label{cor:EA2}
%Let $f,h:\mathbb{F}^n_2\rightarrow \mathbb{F}_2$ be trivial $s$-plateaued functions with affine Walsh supports related as $S_h=c+S_fA^T$ (for some $c\in \mathbb{F}^n_2$ and $A\in GL(n,\mathbb{F}_2))$. If their duals $f^*$ and $h^*$ $\in \mathcal{B}_{n-s}$ (defined via (\ref{eq:equiv22})) are related in a non-linear manner, i.e., $h^*(z_i)=f^*(\omega_i)+g(\omega_i)$ ($g(x_i)\leftrightsquigarrow g(\omega_i)\in \mathcal{B}_{n-s}$ is a non-linear function), then $f$ and $h$ are EA-inequivalent.
%\end{cor}
%\proof Suppose that $f$ and $h$ are EA-equivalent, i.e., for some $B\in GL(n,\mathbb{F}_2)$, $d,p\in \mathbb{F}^n_2$ and $\tau\in \mathbb{F}_2$  let
%$$h(x)=f(xB+d)+p\cdot x+\tau,\;\;\;\;x\in \mathbb{F}^n_2.$$
%By relation (\ref{eq:equiv}) we have that $S_h=p+S_fB^T$ (where $\widehat{z}_i=p+\omega_iB^T$, $\omega_i\in S_f$) and let the dual of $h$ defined in vectors $\widehat{z}_i$ be denoted by
%\begin{eqnarray}\label{eq:hcapa}
%\widehat{h}^*(\widehat{z}_i)=f^*(\omega_i)+d\cdot \omega_i+\tau.
%\end{eqnarray}
%Now, for a matrix $K\in GL(n,\mathbb{F}_2)$ with $A^TK=B^T$ and a (suitable) vector $\lambda\in \mathbb{F}^n_2$ we have
%\begin{eqnarray}\label{eq:zi}
%\lambda+z_iK=\lambda+(c+\omega_iA^T)K=(\lambda+cK)+\omega_iB^T=p+\omega_iB^T=\widehat{z}_i,
%\end{eqnarray}
%i.e.,  $\widehat{z}_i=\lambda+z_iK$ are related in an affine manner. Consequently, since $S_h=c+S_fA^T$ and $S_h=p+S_fB^T$, then the values of $\widehat{h}^*(\widehat{z}_i)$ are actually equal to values of $h^*$ at vectors $\lambda+z_iK=\widehat{z}_i$ (due to (\ref{eq:zi})), i.e., we have that $\widehat{h}^*(\widehat{z}_i)=h^*(\lambda+z_iK).$ %This simply means that $\widehat{h}^*$ and $h^*$ are affine equivalent.
%
%On the other hand, by (\ref{eq:hcapa}) and the assumption that $h^*$ is defined as $h^*(z_i)=f^*(\omega_i)+g(\omega_i)$, we have that
%\begin{eqnarray*}
%\left\{\begin{array}{l}
%         h^*(z_i)=f^*(\omega_i)+g(\omega_i) \\
%          \widehat{h}^*(\widehat{z}_i)=h^*(\lambda+z_iK)=f^*(\omega_i)+d\cdot \omega_i+\tau
%       \end{array}
%\right.
%&\Leftrightarrow&\;\;\; h^*(z_i)+h^*(\lambda+z_iK)=g(\omega_i)+d\cdot \omega_i+\tau.
%\end{eqnarray*}
%is equivalent to
%$$h^*(z_i)+h^*(\lambda+z_iK)=g(\omega_i)+d\cdot \omega_i+\tau,\;\;\;;\;\;i\in[0,2^{n-s}-1].$$
%Since our identification between $S_h$ and $\mathbb{F}^{n-s}_2$ is given in terms of indices, and the fact that $\lambda$ and $K$ are fixed, it is clear that for %some $\theta\in \mathbb{F}^{n-s}_2$ and matrix $R\in GL(n-s,\mathbb{F}_2)$ the previous equality can be written as
%$$h^*(x_i)+h^*(x_iR+\theta)=g(x_i)+\ell(x_i),$$
%where we identified $h^*(x_i+\theta)\leftrightsquigarrow h^*(\lambda+z_iK)$ (since $x_i\leftrightsquigarrow z_i\in S_h$ and $\lambda+z_iK\in S_h$), $g(x_i)\leftrightsquigarrow g(\omega_i)$ and $\ell(x_i)\leftrightsquigarrow d\cdot \omega_i+\tau$, where $\ell\in \mathcal{A}_{n-s}$ due to Proposition \ref{prop:Hrow}.\\\\
%\textbf{\textcolor[rgb]{0.98,0.00,0.00}{Samir:}  So, in order to get some "contradiction", then I have the following idea. Let us say that for arbitrary fixed $R$ and $\theta$ there exists an hyperplane, say $H$, such that $H\cap (HR+\theta)=\emptyset$, then because $h^*(x_i)$ and $h^*(x_iR+\theta)$ are bent functions, then the restrictions $h^*|_{H}$ and $h^*|_{HR+\theta}$ are two disjoint spectra semi-bent functions, and the sum of every two such semi-bent functions is either affine or balanced function (this can be proven in a few lines). And then, if $g+\ell$ is not balanced or affine for every $\ell\in \mathcal{A}_{n-s}$, then we get that $f$ and $h$ are not equivalent. Thus, in our Corollary 3.1, we will just add this condition "$g+\ell$ is not balanced or affine for every $\ell\in \mathcal{A}_{n-s}$"....The main question, does for every $R$ and $\theta$ there exists a hyperplane $H$ such that $H\cap (HR+\theta)=\emptyset$?? (I think that it should be the case...is there any explicit method to show this?)}

%By Theorem \ref{th:EA1} and Proposition \ref{prop:Hrow} we have that  $h^*$ and $f^*$ are EA-equivalent, that is, $h^*(z_i)=f^*(\omega_i)+b\cdot \omega_i+\varepsilon$, where  $z_i=c+\omega_iA^T$ and $b\cdot \omega_i=\phi_b(\omega_i)\leftrightsquigarrow \ell(x_i)\in \mathcal{A}_{n-s}$. Consequently, we have that $h^*(x_i)\leftrightsquigarrow h^*(z_i)$ and $f^*(x_i)\leftrightsquigarrow f^*(\omega_i)$ ($x_i\in \mathbb{F}^{n-s}_2$) are related in affine manner, i.e., $h^*+f^*=\ell$, which is a contradiction to the fact that $h^*+f^*=g$ and $g$ is non-linear. \qed
%%%%%%%%%%%%%%%%%%%%%%%%%%%%%%%%%%%%%%%%
%Since this conclusion holds regardless of the choice of $A$, then $f$ and $h$ are EA-inequivalent.\qed
%\begin{rem}\label{rem:nonlinear}
%Note that  $h^*+f^*=\ell$ in the proof above holds for every matrix $A\in GL(n,\mathbb{F}_2)$ and $c\in \mathbb{F}^n_2$ with $S_h=c+ S_fA^T$, and any ordering of $S_f$ with the property that $\phi_b(\omega_i)\leftrightsquigarrow \ell(x_i)\in \mathcal{A}_{n-s}$. However, if for a fixed matrix $A$ (and ordering of $S_f$) we have that $h^*$ and $f^*$ are related in non-linear manner, then $f$ and $h$ are EA-inequivalent (where $f$ and $h$ are trivial plateaued functions).
%\end{rem}

%Using Proposition \ref{prop:hbent} we have that the EA-inequivalence between two bent functions $f^*,h^*:\mathbb{F}^{n-s}_2\rightarrow\mathbb{F}_2$ (which are duals of trivial $s$-plateaued functions $f,g:\mathbb{F}^{n}_2\rightarrow\mathbb{F}_2$), by Proposition \ref{prop:EA1} actually implies the EA-inequivalence between $s$-plateaued functions functions $f$ and $h$. This observation is formalized with the following result.
%\begin{cor}\label{cor:EA2}
%Let $f,h:\mathbb{F}^n_2\rightarrow \mathbb{F}_2$ be trivial $s$-plateaued functions with affine Walsh supports related as $S_h=c+S_fA^T$ (for some $c\in \mathbb{F}^n_2$ and $A\in GL(\mathbb{F}^n_2)$). If their duals $f^*$ and $h^*$ (defined via (\ref{eq:equiv22})) are inequivalent bent functions on $\mathbb{F}^{n-s}_2$, then $f$ and $h$ are EA-inequivalent.
%\end{cor}

%{\bf In the whole section we have not used this representation $S_h=v +EM^T$ for the purpose of establishing equivalence between $f$ and $h$ ??? We devote half a page to introduce this representation which includes $M$ fro this purpose and then we use it for other purposes ???}\\\\
%\textbf{\textcolor[rgb]{0.98,0.00,0.00}{Samir:} At the end of page 7, it has been fixed that $S_h$ is in the form $S_h=\widehat{v}+EM$, where $M=A^T$ and $\widehat{v}=c+vA^T$, which is actually obtained by representing $S_f$ as $S_f=v+E$. The representation $S_h=\widehat{v}+EM$ one clearly obtains from equality $S_h=c+S_fA^T$, where we use $S_f=v+E$.}

\subsection{Generic construction of EA-inequivalent plateaued functions}\label{sec:ineq}

A generic construction method of plateaued functions, which specifies these functions in their Walsh spectra, has been introduced in \cite[Section 3]{Secondary}. In order to distinguish it from the direct or shortly \emph{ANF-construction methods}\footnote{An  \emph{ANF-construction method} refers to any construction technique which specifies a form or ANF of the function, for which it can be shown to posses the plateaued property.}, the construction technique given in \cite[Section 3]{Secondary} will be called a \emph{spectral method}\footnote{A \emph{spectral method} refers to any construction method which specifies a suitable Walsh spectrum (its Walsh support and the signs of nonzero spectral values) which after applying the inverse WHT gives a  Boolean function with prescribed properties.}.
%An example of a trivial plateaued function constructed by the spectral method was given by Example 3.1 in \cite{Secondary}. In addition, Example 3.2 in \cite{Secondary} it has been indicated that the non-trivial plateaued functions can be also constructed by the spectral method, if we utilize  \cite[Theorem 3.1-$(i)$]{Secondary}.
Utilizing Proposition \ref{prop:Hrow}, the following result generalizes \cite[Theorem 3.1]{Secondary} and shows how one can efficiently construct plateaued functions (using the spectral method) by using the representation of a Walsh support as in (\ref{DPL}).
%Although the proof is elementary, it is given for self-completeness.
\begin{theo}\label{theo:plateH}
%\begin{theo} \label{theo:nonafine}
Let  $S_f=v+ EM=\{\omega_0,\ldots,\omega_{2^{n-s}-1}\} \subset \FB^n$, %with $\#S_f=2^{n-s}$,
for some $v \in \mathbb{F}^n_2$, $M\in GL(n,\mathbb{F}_2)$ and subset $E=\{e_0,e_1, \ldots, e_{2^{n-s}-1}\}\subset \mathbb{F}^n_2$. For a function $g:\FB^{n-s}\rightarrow \mathbb{F}_2$ such that  $wt(g)=2^{n-s-1}+ 2^{\frac{n-s}{2}-1}$ or $wt(g)=2^{n-s-1}-2^{\frac{n-s}{2}-1}$, let the Walsh spectrum of $f$ be defined (by identifying $x_i \in \FB^{n-s}$ and $\omega_i \in S_f$ through $e_i \in E$) as
\begin{eqnarray}\label{eq:walshX}
W_f(u)= \left \{  \begin{array}{ll}
  2^{\frac{n+s}{2}}(-1)^{g(x_i)} & \textnormal{ for } u= v + e_iM \in S_f,\\
 0 & u \not \in S_f. \\
\end{array}  \right.
\end{eqnarray}
Then:
\begin{enumerate}[i)]
\item $f$ is an $s$-plateaued function if and only if $g$ is at bent distance to $\Phi_f$ defined by (\ref{RSf}).
\item If $E\subset \mathbb{F}^n_2$ is a linear subspace, then $f$ is an $s$-plateaued function if and only if $g$ is a bent function on $\mathbb{F}^{n-s}_2$.
\end{enumerate}
\end{theo}
\proof $i)$ By relations (\ref{WHT}) and (\ref{eq:walshX}), for arbitrary $u\in \mathbb{F}^n_2$, we have that
\begin{eqnarray}\label{eq:g1}
2^n(-1)^{f(u)}=\sum_{\omega\in S_f}W_f(\omega)(-1)^{u\cdot \omega}=2^{\frac{n+s}{2}}(-1)^{u\cdot v}\sum_{e\in E}(-1)^{f^*(v+eM)+uM^T\cdot e}
\end{eqnarray}
Since we have the correspondence  $f^*(\omega_i)=f^*(v+e_iM)=g(x_i)$ and $uM^T\cdot e_i=\phi_{uM^T}(x_i)\in \Phi_f$ ($x_i\in \mathbb{F}^{n-s}_2$, $e_i\in E$, $i\in[0,2^{n-s}-1]$), then using (\ref{eq:g1})  we get
\begin{eqnarray}\label{eq:g2}
2^{\frac{n+s}{2}}(-1)^{u\cdot v}\sum_{e\in E}(-1)^{f^*(v+eM)+uM^T\cdot e}=2^{\frac{n+s}{2}}(-1)^{u\cdot v}\sum_{x_i\in \mathbb{F}^{n-s}_2}(-1)^{g(x_i)+\phi_{uM^T}(x_i)}.
\end{eqnarray}
Thus,  we have that $f$ is $s$-plateaued if and only if $\sum_{x_i\in \mathbb{F}^{n-s}_2}(-1)^{g(x_i)+\phi_{uM^T}(x_i)}=2^{\frac{n-s}{2}} (-1)^{\tau_u}$, for some $\tau_u\in \{0,1\}$. However,  the bent distance between $g$ and $\phi_{uM^T}\in\Phi_f$ in (\ref{eq:g2}) is in accordance with relation (\ref{eq:g1}) since these two relations imply that $(-1)^{f(u)}=(-1)^{u\cdot v+ \tau_u}$ holds. Thus, the inverse WHT applied to the spectrum given by (\ref{eq:walshX}) generates a Boolean function which is $s$-plateaued.

$ii)$ If $E$ is a linear subspace (ordered lexicographically),  by Proposition \ref{prop:Hrow} the sequence profile $\Phi_f$ contains sequences of all linear functions (or their affine equivalents) and  $g$ has to be a bent function.\qed

Essentially, combining Theorem \ref{theo:plateH}  with results given in Section \ref{sec:equivalence} gives us the possibility  to construct inequivalent plateaued functions regardless of whether they are trivial or not. For instance, the construction of inequivalent trivial plateaued functions can be deduced as follows. We first remark  that for any two trivial $s$-plateaued functions $f,h:\mathbb{F}^{n}_2\rightarrow \mathbb{F}_2$ one can always relate their Walsh supports as $S_h=c+S_fA^T$, for  a matrix $A\in GL(n,\mathbb{F}_2)$ and some vector $c\in \mathbb{F}^n_2$. Then, one simply employs EA-inequivalent bent functions in $n-s$ variables as duals for $f$ and $h$ and the inequivalence follows by Theorem \ref{th:EA1}.

On the other hand, the construction of inequivalent non-trivial plateaued functions (by Theorem \ref{theo:plateH}) is achieved by setting non-affine Walsh supports of different dimensions (in combination with suitable duals).
%%%%%%%%%%%%%
%Then, by Theorem \ref{cor:EA2}, it is sufficient to define their duals $f^*$ and $h^*$ so that they are not related in an affine manner as $h^*=f^*+\ell$, for some $\ell\in \mathcal{A}_{n-s}$. Another possibility is to utilize Theorem \ref{th:EA1} and employ EA-inequivalent bent functions in $n-s$ variables as duals for $f$ and $h$.
%
The following example illustrates the former case, thus employing Theorem \ref{th:EA1} for this purpose.
 %
\begin{ex}\label{ex:ineq1}
Let $f:\mathbb{F}^7_2\rightarrow \mathbb{F}_2$ ($n=7$) be a semi-bent function ($s=1$) defined as $f(x_1,\ldots,x_7)=x_1 x_4 + x_2 x_5 + x_3 x_6 + x_4 x_5 x_6$. One can verify that the Walsh support of $f$ is $S_f=v+E=\mathbb{F}^6_2\times \{0\}\subset \mathbb{F}^7_2$, where $v=\textbf{0}_7$ and $e_i\in E$ are given as $e_i=(x_i,0)$ with $x_i\in \mathbb{F}^6_2$ (and $\mathbb{F}^6_2$ is ordered lexicographically). In addition, the dual $f^*(\omega_i)=f^*(v+e_i)\stackrel{(\ref{DPL})}{=}\overline{f}^*(x_i)$ ($\overline{f}^*:\mathbb{F}^6_2\rightarrow \mathbb{F}_2$) is given as
$$\overline{f}^*(x_1,\ldots,x_6)=(x_1,x_2,x_3)\cdot (x_4,x_5,x_6)+x_1x_2x_3.$$

To construct a (trivial) semi-bent function $h:\mathbb{F}^5_2\rightarrow \mathbb{F}_2$ which is EA-inequivalent to $f$,  by Theorem \ref{th:EA1}, it is sufficient to take a bent function $\overline{h}^*:\mathbb{F}^6_2\rightarrow \mathbb{F}_2$ as its dual which is affine inequivalent to $\overline{f}^*$. For instance, let us define $\overline{h}^*$ as
$$\overline{h}^*(x_1,\ldots,x_6)=(x_1,x_2,x_3)\cdot (x_4,x_5,x_6).$$
Clearly, $\overline{f}^*$ and $\overline{h}^*$ are inequivalent since they are of different degrees.  Furthermore, taking that $S_h=(1,0,1,0,1,1,1)+S_f=(1,0,1,0,1,1,1)+\mathbb{F}^6_2\times \{0\}$, then by Theorem \ref{theo:plateH} (considering $h^*(z_i)=h^*((1,0,1,0,1,1,1)+(x_i,0))= \overline{h}^*(x_i)$, $x_i\in \mathbb{F}^6_2$) we obtain the function
$$h(x_1,\ldots,x_7)=x_1 + x_3 + x_1 x_4 + x_5 + x_2 x_5 + x_6 + x_3 x_6 + x_7.$$
Clearly, $f$ and $h$ are inequivalent due to $deg(f)\neq deg(h)$.
\end{ex}
%\textbf{\textcolor[rgb]{0.98,0.00,0.00}{Samir:} Theoretically, functions $f$ and $h$ are inequivalent in the example above.\\\\ \textcolor[rgb]{0.98,0.00,0.00}{Question for Fengrong:} Can we verify this by computer search? (Maybe your students can check this?)}\\\\
Some generic construction methods,  using the spectral approach of Theorem \ref{theo:plateH}, of nontrivial plateaued functions are presented in Section \ref{sec:nonaffine}. Two EA-inequivalent nontrivial plateaued functions can be constructed by ensuring that their Walsh supports contain different number of linearly independent vectors (thus applying Theorem \ref{th:eq}-$(i)$). Alternatively, if their supports have the same number of linearly independent vectors then the idea from Example \ref{ex:ineq1} can be used where additionally the condition of Theorem \ref{th:eq}-$(ii)$ on duals has to be satisfied.





\subsection{Plateaued functions without linear structures}\label{sec:linear}


It has already been observed that trivial plateaued functions essentially correspond to partially bent functions {\cite{Zheng,Ayca2}} introduced by Carlet in \cite{CarletPartial}. Furthermore, it is well-known that partially bent functions  always admit linear structures. Therefore, it is of interest to analyze whether  nontrivial plateaued functions admit linear structures which extends the work of Zhang and Zheng \cite{Zheng} where a design method of nontrivial plateaued functions without linear structures were given in terms of sequences.
%\\{\bf We have to be careful here !!! It might be the case that only trivial plateaued functions admit linear structures  !! Can we essentially have any nontrivial plateaued function with linear structures ??
%Maybe we should provide an example of nontrivial functions having linear structures to motivate this section even better !!}\\\\
%\textbf{\textcolor[rgb]{0.98,0.00,0.00}{Samir:} An example is given at the end of this subsection.}\\\\
%
A useful concept related to a Boolean function $f \in \FB^n$ is the so-called autocorrelation spectrum whose spectral values are given by  $\triangle_f(a)=\sum_{x\in \FB^n}(-1)^{f(x)+f(x+a)}$, for $a \in {\FB^n}^*$. The property of being partially bent can be stated as follows.

\begin{defi}\label{def partially bent}\cite{CarletPartial}
  A function $f\in B_n$ with the Walsh support $S_f$  which satisfies the equality
  $\# S_f\cdot \# \mathfrak{R}=2^n$, where $\mathfrak{R}=\{u\in \mathbb{F}^n_2:\; \Delta_f(u)\neq 0\}$,
  %$(2^n-Z_{\triangle_f})(2^n-Z_{W_f})=2^n$
  is called partially-bent.
 % where $\mathfrak{R}=\{u\in \mathbb{F}^n_2:\; \Delta_f(u)\neq 0\}.$
%  where $Z_{\triangle_f}=||\{a\in \FB^n|\triangle_f(a)=0\}|| $ and $Z_{W_f}=||\{u\in \FB^n|W_f(u)=0\}||$.
\end{defi}
In has been shown  \cite[Corollary 3.1]{CCA}, using somewhat different notation, that the autocorrelation values can be computed in terms of spectral values over the Walsh support.
%The following result, which uses somewhat different notation, is actually a special case of \cite[Theorem 3.1]{CCA}.
\begin{lemma}\cite{CCA}\label{lemma:connection}
Let $f \in \mathcal{B}_n$ be an arbitrary functions with Walsh support $S_f$ represented as $S_f=v+ E$, for some subset $E$ of $\FB^n$ and $v\in S_f$. Then
\begin{equation}\label{eq:connect}
\triangle_f(a)=2^{-n}(-1)^{v\cdot a}\sum\limits_{u\in E}W_f^2(u+v)(-1)^{u\cdot a}.
\end{equation}
%is an affine subspace of dimension $n-s$  if and only if $f$ is a partially-bent function in $n$ variables.
\end{lemma}
%\begin{proof}
%%Since $f$ is a plateaued function with Walsh spectra $\{0, \pm 2^{\frac{n+s}{2}}\}$, we have $\# S_f= 2^{n-s}$.
%By using the inverse WHT,  for any $a\in \FB^n$, we have \begin{equation*}\begin{array}{rl}\label{equ partially bent}
% \triangle_f(a)&=\sum\limits_{x\in \FB^n}(-1)^{f(x)+f(x+a)}\\
% &=\sum\limits_{x\in \FB^n}(-1)^{f(x)}(-1)^{f(x+a)}\\
% &=\sum\limits_{x\in \FB^n}2^{-2n}\sum\limits_{u\in \FB^n}W_f(u)(-1)^{u\cdot x}\sum\limits_{v\in \FB^n}W_f(v)(-1)^{v\cdot (x+a)} \;\; \textnormal{ (by inverse %WHT)}\\
% &=\sum\limits_{x\in \FB^n}2^{-2n}\sum\limits_{v\in \FB^n}\sum\limits_{u\in \FB^n}W_f(u)(-1)^{u\cdot x}W_f(u+v)(-1)^{(v+u)\cdot (x+a)}\\
% &=\sum\limits_{x\in \FB^n}2^{-2n}\sum\limits_{v\in \FB^n}\sum\limits_{u\in \FB^n}W_f(u)(-1)^{u\cdot x}W_f(u+v)(-1)^{(v+u)\cdot (x+a)}\\
%  &=\sum\limits_{x\in \FB^n}2^{-2n}\sum\limits_{u\in \FB^n}W_f^2(u)(-1)^{u\cdot a}\\
%& =2^{-n}\sum\limits_{u\in \FB^n}W_f^2(u)(-1)^{u\cdot a}\\
%& =2^{-n}\sum\limits_{u\in S_f}W_f^2(u)(-1)^{u\cdot a}\\
%& =2^{-n}(-1)^{\alpha\cdot a}\sum\limits_{u\in E}W_f^2(u+\alpha)(-1)^{u\cdot a}.
% \end{array}
% \end{equation*}
%%If $S_f=\alpha + E$ is an affine subspace of dimension $n-s$, from (\ref{equ partially bent}), we have $\# \mathfrak{R}=2^{n-(n-s)}$, that is, $f$ is a partially-bent function in $n$ variables.
%%
%%If $f$ is a partially-bent function in $n$ variables, we have $\# \mathfrak{R}=2^{s}$ and $\mathfrak{R} $ is a linear subspace of dimension $s$. From (\ref{equ partially bent}), we know $E\subseteq \mathfrak{R}^{\perp}$, where $\mathfrak{R}^{\perp}$ denotes the dual of $\mathfrak{R}$. We also know $\#E= \# S_f= 2^{n-s}=\#\mathfrak{R}^{\perp}$. Hence, $S_f=\alpha + E$ is an affine subspace of dimension $n-s$.
%\qed
%\end{proof}
The following result specifies the existence of linear structures for an arbitrary Boolean function, thus including plateaued functions as a special case.
 \begin{theo}\label{theo baselinear}
Let $f \in \mathcal{B}_n$  whose Walsh spectrum
%with Walsh spectra $\{0, \pm 2^{\frac{n+s}{2}}\}$.
is given by  $S_f= v+E\subset \FB^n$, where  $v\in S_f$.
Let $m$ denote the maximum number  of linearly independent elements of $E$ and  $\Lambda=\{a\in \FB^n: f(x)+f(x+a)=constant, x\in \FB^n \}$. Then $$m+\dim(\Lambda)=n.$$
\end{theo}
\begin{proof}
  If $a $ is a linear structure of $f$,
  then $\triangle_f(a)=2^n ~\emph{or} ~-2^n$.  By relation (\ref{eq:connect}), this is equivalent to \begin{equation}\label{equ lsup00}
\sum\limits_{u\in \FB^n}W_f^2(u)(-1)^{u\cdot a}= (-1)^{v \cdot a}\sum\limits_{u\in E}W_f^2(u+v)(-1)^{u\cdot a}=2^{2n} ~\emph{or} ~-2^{2n}.\end{equation}
Using the Parseval's equality $ \sum\limits_{u\in \FB^n}W_f^2(u)=2^{2n}$ in (\ref{equ lsup00}), we have that
%From the Parseval's equality, we have   \begin{equation}\label{equ lsup01}
%  \sum\limits_{u\in \FB^n}W_f^2(u)=2^{2n}.\end{equation}
%Then, combining (\ref{equ lsup00}) and (\ref{equ lsup01}),
%we have
$$\sum\limits_{u\in E}(-1)^{u\cdot a}=\# S_f\; \textnormal{ or }\;-\# S_f. $$
Consequently, using the fact that $S_f=v+E$ and $v\in S_f$, we have that $0_n\in E$ and thus the sum on the left-hand side above contains the term $(-1)^{0_n \cdot a}=1$.
%
%We also know $S_f=\alpha +E$ and $\alpha\in S_f$, then $0_n\in E$, that is, $(-1)^{0_n \cdot a}=1$.
Hence, $$ \sum\limits_{u\in E}(-1)^{u\cdot a}=\# S_f.$$

 On the other hand,  $\sum\limits_{u\in E}(-1)^{u\cdot a}=\# S_f$ if and only if $E\subseteq \{a\}^{\perp}$.  Similarly,  we have $E\subseteq \Lambda^{\perp}$, that is,
  \begin{equation}\label{equ lsup 1}\dim(E_*)\leq n-\dim(\Lambda),
   \end{equation}where $E_*$ denotes the minimum subspace including $E$.
From (\ref{eq:connect}), we know that  $|\triangle_f(a)|=2^n$ for any $a\in E_*^{\perp} $. Thus, we have $ E_*^{\perp}\subseteq\Lambda$, that is,
\begin{equation}\label{equ lsup 2} \dim(\Lambda)\geq n-\dim(E_*).
   \end{equation}
  Combining (\ref{equ lsup 1}) and (\ref{equ lsup 2}),  we have $m+\dim(\Lambda)=n.$ \qed

\end{proof}

 \begin{cor}\label{theo:baselinear2}
Let $f:\mathbb{F}^n_2\rightarrow \mathbb{F}_2$ be an arbitrary $s$-plateaued function with its Walsh spectrum given as  $S_f= v+E\subset \FB^n$, where  $v\in S_f$.
%Let $m$ denote the maximum number  of linearly independent elements of $E$ and  $\Lambda=\{a\in \FB^n: f(x)+f(x+a)=constant, x\in \FB^n \}$. Then $$m+\dim(\Lambda)=n.$$
Then $f$ does not contain (non-zero) linear structures if and only if the number of linearly independent vectors of $E$ is equal to $n.$
\end{cor}
\begin{proof}
Let us assume that $f$ has a non-zero linear structure $a\in \mathbb{F}^n_2$. Then by relation (\ref{eq:connect}) we have that
$$\Delta_f(a)=2^{-n}(-1)^{v\cdot a}2^{n+s}\sum_{u\in E}(-1)^{u\cdot a}=2^{s}(-1)^{v\cdot a}\sum_{u\in E}(-1)^{u\cdot a},$$
and thus $a$ is a linear structure of $f$ (that is $|\Delta_f(a)|=2^n$) if and only if $\sum_{u\in E}(-1)^{u\cdot a}=\pm 2^{n-s}$. However, using the fact that $\# E=2^{n-s}$, we have that
\begin{eqnarray}\label{eq:phi}
\pm 2^{n-s}=\sum_{u\in E}(-1)^{u\cdot a}=2^{n-s}-2wt(\varphi_a),
\end{eqnarray}
where the truth table of the function $\varphi_a$ is given as $T_{\varphi_a}=(u\cdot e_0,\ldots,u\cdot e_{2^{n-s}-1})$ (and $E=\{e_0,\ldots,e_{2^{n-s}-1}\}$ can be considered in any ordering). However, (\ref{eq:phi}) holds if and only if $\varphi_a$, viewed as a function on $\mathbb{F}^{n-s}_2$, is a constant function for $a\neq \textbf{0}_n$ (thus equal to $0$ or $1$).

 Equivalently, representing $E$ as a matrix, using $e_0,\ldots,e_{2^{n-s}-1}$ as its rows, then some of its columns are either constant or there is a linear dependency among its columns. Thus, the number $E$ of linearly independent columns of $E$ is less than $n$.
 %thus this number is less than $n$.
\qed
\end{proof}
It is rather challenging to conjecture that only trivial plateaued functions admit linear structures which is however not true as illustrated in the example below.
\begin{ex}
Let us consider the function $f:\mathbb{F}^6_2\rightarrow \mathbb{F}_2$ $(n=6)$ given as $$f(x_1,\ldots,x_6)=x_1 (x_3 + x_2 x_5) + x_2 (x_4 + x_6).$$
One can verify that the autocorrelation spectrum of $f$ is given by
$$(\Delta_f(u_0),\ldots,\Delta_f(u_{63}))=(64, 0, 32, 32, 0, 64, 32, 32, 0, 0, 32, -32, 0, 0, -32, 32,\bf{0}_{48}),$$
where $\bf{0}_{48}$ means that the last $48$ coefficients are equal to $0$. Also, the Walsh support of $f$ is given as $S_f=\mathbb{F}^4_2\wr T_g\wr T_{\ell}$, where $g,\ell:\mathbb{F}^4_2\rightarrow \mathbb{F}_2$ are given as $g(x_1,\ldots,x_4)=x_3x_4$ and $\ell(x_1,\ldots,x_4)=x_4$, and thus $f$ is a non-trivial plateaued function. Since $S_f$ contains only $5$ linearly independent vectors, then $dim(\Lambda)=n-m=6-5=1$ (Theorem \ref{theo baselinear}) and we have that $\Lambda=\{{\bf 0}_6,(0, 0, 0, 1, 0, 1)\}.$
%Notice that only for $u=(0, 0, 0, 1, 0, 1)$ it holds that $\Delta_f(u)=64$.
%One can verify that the autocorrelation coefficients of $f$ take the values $\Delta_f(u)\in\{\pm 32, 64\}$ ($u\in \mathbb{F}^6_2$), where only for $u\in \{\textbf{0}_6,(0, 0, 0, 1, 0, 1)\}$ it holds that $\Delta_f(u)=64.$
\end{ex}
Clearly, if the Walsh support of $f$ contains $n$ linearly independent vectors then $\dim(\Lambda)=0$ and such a function cannot have linear structures. This result will be useful in Section \ref{sec:nontrivial} to draw a conclusion  that the constructed nontrivial plateaued functions are without linear structures.
\begin{rem}
The above result implies that a trivial $s$-plateaued function on $\FB^n$ (which corresponds to a partially bent function) always admit linear structures of dimension $s$, see also \cite{Zheng}. This is a consequence of the fact that $\dim(S_f)=n-s=m$ and thus $\dim(\Lambda)=s$.
\end{rem}
%\textbf{\textcolor[rgb]{0.98,0.00,0.00}{Samir:} Just comment -- The fact that for partially bent functions it holds that $\dim(\Lambda)=s$  is known in the paper of Zheng - "On plateaued functions", 2001.}
%I feel we can construct plateaued functions without nonzero linear structure by using Theorem \ref{th:Sf}.  From Theorem \ref{th:Sf}, we can directly construct plateaued functions without nonzero linear structure since $t_i$ is arbitrary function.  The function in Example 4.1 has no nonzero linear structure.


\section{Constructions of (non)trivial disjoint spectra plateaued functions} \label{sec:nontrivial}

%\textbf{\textcolor[rgb]{0.98,0.00,0.00}{Samir:} I would delete the  blue sentence below, and in the last sentence the part "In this direction" would change to "In this section"}\\\\
%%
%\textcolor[rgb]{0.00,0.00,1.00}{In this section we provide generic constructions of (disjoint spectra) plateaued functions with (affine or) non-affine Walsh supports.}
A construction method of  trivial disjoint spectra plateaued functions is easily established by using Theorem \ref{theo:plateH} and setting the Walsh supports to be disjoint affine subspaces (where duals can be arbitrary bent functions).
%Nevertheless,  in Section \ref{sec:MM} we demonstrate  that the same task can be accomplished in the ANF domain using the $\mathcal{MM}$ class of functions.
The design of nontrivial plateaued functions and in particular the problem of deriving  disjoint spectra nontrivial $s$-plateaued functions of maximum cardinality appears to be a more difficult task. In this direction, we provide some efficient design methods  of  nontrivial plateaued functions  which also might have  interesting implications on the design of bent functions.%\\ \\
%{\bf Enes: Check out in this section whether we have a proper use of notation !! What I mean here is whether we use $f^*$ or we need to use $\overline{f}^*$ ?? No weird arrow however !!}

\subsection{Disjoint spectra plateaued functions  via $\mathcal{MM}$ class}\label{sec:MM}

For convenience and due to the consequences related to generic design methods of bent functions, we briefly discuss the possibility of obtaining trivial disjoint spectra plateaued functions of maximal cardinality. Since any trivial $s$-plateaued function has an affine subspace $S_f=\alpha + E$ of dimension $n-s$ as its Walsh support, it is natural to consider the space $\FB^n$ as a union of disjoint cosets, thus $\FB^n=\cup_{\alpha_i \in \mathcal{G}} \alpha_i +E$ ($\mathcal{G}$ is orthogonal complement of $E$), for some linear subspace $E$ and elements $\alpha_i$ such that $\alpha_i + E \cap \alpha_j +E = \emptyset$. Then, clearly $\#\mathcal{G}=2^{s}$ and any function $f_i \in B_n$ whose support is defined on the affine subspace $S_{f_i}=\alpha_i +E$ is a disjoint trivial plateaued function to $f_j$, for $\alpha_i \neq \alpha_j$. It is important to notice that such a set of disjoint spectra plateaued functions (of maximal cardinality) can be defined using  arbitrary  bent duals for any $f_i$ which additionally may belong  to different classes of bent functions. This leads to an important research problem that can be stated as follows.
\begin{op}
Assume that a  set of disjoint spectra trivial $s$-plateaued functions on $\FB^n$ of (maximal) cardinality $2^s$ is constructed using the dual bent functions (by means of Theorem \ref{theo:plateH}) which are taken from different classes of bent functions on $\FB^{n-s}$.  To which class of bent functions a bent function on $\FB^{n+s}$, defined as a concatenation of these $2^s$ functions, do belong?
\end{op}
%{\bf \textcolor[rgb]{0.98,0.00,0.00}{Enes:} We need to do something about this problem to get a strong paper !!!} \\ \\
 The above approach, specifying the bent dual in spectral domain, provides much more flexibility  compared to a standard design in the ANF domain. The only generic class of Boolean functions that admits an efficient construction method of disjoint spectra plateaued functions of maximal cardinality in the ANF domain is the Marioana-McFarland ($\mathcal{MM}$) class \cite{MMclass}. This calls  was initially introduced as a primary class of bent functions though it can be used for plateaued functions as well.

Let $n=s+k$ where $s > n/2$ and denote the set of linear functions in $s$ variables as $\mathcal{L}_s=\{a_1x_1 + \ldots + a_s x_s: a_i \in \mathbb{F}_2\}$. Let $f(y,x)= \phi(y) \cdot x + g(y)$, where $x=(x_1,\ldots,x_s) \in \FB^s$, $y=(y_1,\ldots,y_k) \in \FB^k$, and $\phi : \FB^k \rightarrow \FB^s$ be an injective mapping. Since $s>k$ we can partition the vector space $\FB^s$ of cardinality $2^s$ into $2^{s-k}$ disjoint subsets $E_i$ each of cardinality $2^k$. That is, $\FB^s=\cup_i^{2^{s-k}}E_i$, where $E_i \cap E_j = \emptyset$.
Then we can define  a set of functions
\begin{equation} \label{eq:mmdisjoint}
f^{(i)}(y,x)=\phi_i(y) \cdot x + g_i(y), \;\;  i=1, \ldots, 2^{s-k},
\end{equation}
where  $\phi_i : \FB^k \rightarrow E_i \subset \FB^s$ is a bijection and $g_i$ is arbitrary. It is straightforward to verify that each $f_i$ is $s$-plateaued (see also \cite{CarletQ}) and that $f_i$ are disjoint spectra functions.
%  and construct $f^{(i)}(y,x)=\phi_i(y) \cdot x + g_i(y)$.

%{\color{red}Fengrong: I feel we need to specify $x_i\in F_2$ or $x_i\in F_2^n$.}
%for each such subset $\mathcal{L}_s^i,$ where $i=1, ..., 2^{s-k}$, we can construct a plateaued function $$f^{(i)}(y,x)= \sum_{\sigma \in \FB^k}\prod_{j=1}^k(y_j+1 + \sigma_j)l^{(i)}_{[\sigma]}(x) + g^{(i)}(y),$$ where $l^{(i)}{[\sigma]}(x)$ uniquely identifies a linear function from the set $\mathcal{L}_s^i=\{l_0^{(i)}, ..., l_{2^k-1}^{(i)}\}$, where $[\sigma]$ is simply the integer representation of $\sigma \in \FB^k$.

\begin{prop}\label{prop:disspec_affine}
The set of functions $f^{(i)}$, for $i=1, \ldots, 2^{s-k}$, defined by (\ref{eq:mmdisjoint}) is a set of disjoint spectra plateaued functions. Furthermore, if the sets $E_i$ used to define the functions $\phi_i : \FB^k \rightarrow E_i \subset \FB^s$ are not affine subspaces then $f_i$ are nontrivial disjoint spectra plateaued functions.
\end{prop}

\begin{rem}
%The above method of splitting the space of linear functions into disjoint subsets does not necessarily implies that these subsets are affine subspaces.
It can be easily checked that if $E_i$ are  affine subspaces then  the concatenation of $f^{(1)}, \ldots,f^{(2^{s-k})}$  will again provide a bent function in the $\mathcal{MM}$ class.
%\textcolor[rgb]{0.00,0.00,1.00}{The same is essentially true (which is slightly more difficult to verify) when $E_i$ are not affine subspaces.}{\bf Correct ????}
%On the other hand, it is not completely clear whether taking $E_i$ to be arbitrary disjoint subsets may lead to bent functions (which are a concatenation of $f_i$ for $i=1,\ldots,2^{s-k}$) that lie outside the $\mathcal{MM}$ class.
\end{rem}
%\textbf{\textcolor[rgb]{0.98,0.00,0.00}{Samir:} The blue sentence I would rather delete. If we consider a bent function, say $x\cdot \phi(y)+g(y)$, then it is clearly a concatenation of linear functions on affine subspaces.....but using the non-affine sets $E_i$ above, I do not see clearly whether this is still the case...thus I suggest the deletion.}
%{\bf The above observation is a bit tricky !! We say that we can simply split the space into disjoint SUBSETS which essentially covers nontrivial plateaued functions !! So the whole generic approach given in the next section can then be viewed as a special case of the above comment ?! We should either discuss the above approach in more detail - is it really true that there are no plateaued conditions using the above MM based approach ?? It is clear that these functions are disjoint spectra functions but how we ensure that they are plateaued for arbitrary space division into disjoint subsets ?}
%The  analysis regarding the classification of bent functions obtained by means of disjoint spectra trivial plateaued functions, whose duals are taken from different classes of bent functions, is omitted here. Nevertheless, there are strong intuitive indications that this approach may give rise to new classes of bent functions if the duals (bent functions) are selected properly.
%In the first place, this conjecture is based on the fact that the form of duals for two main (explicit) primary classes ($\mathcal{MM}$ and $\mathcal{PS}_{ap}$) of bent functions are known. Then, concatenating disjoint spectra $s$-plateaued functions whose duals have different ANF forms does not necessarily imply that the overall dual of the resulting bent function has the form corresponding to those of the primary classes.
%\begin{ex}\label{ex:possiblyNEQ}
%In this example we want to construct a bent function on $\FB^5$ by using two disjoint spectra semi-bent functions (thus $1$-plateaued). We can consider the same function as given in Example 3.1 and define the other one using a disjoint affine subspace than the one used to define $S_f$.
%content...
%\end{ex}
%The above part is trivial but we may keep it so far for clarity. This should be equivalent to the approach already given, thus considering the cosets of $\FB^{n-s}$ as already mentioned. \\
%{\bf Our main goal is to end up with something different using different supports for our Walsh supports (possible different degrees, classes etc. ) We may easily have a nice discussion here as follows: \\ \\
%-- We know that the dual of a function in MM class is again a bent function in MM !\\ \\
%-- We know that the dual of a function in PS is again in PS class \\ \\
%-- We know that a concatenation of functions in MM class is again a function in MM class (this can be easily proved and maybe we should give a proof here for convenience ! unless this is well-known then we need a reference) \\ \\
%-- It follows that taking disjoint spectra plateaued functions in MM class gives a bent function in MM class \\ \\
%-- It should be clear that taking these affine subspaces of disjoint supports on which we define different Walsh supports by mixing e.g. MM and PS class gives a function which is not in MM \\ \\
%{\color{red} Fengrong: I think this is a good idea. I had also this idea that a bent function is obtained by using different  disjoint spectra functions :-). However, I do not know how to do.  If we select the initial functions which were constructed by a secondary construction, then the bent functions which were obtained by concatenating the initial functions can be also  constructed by the secondary construction. For example, if $f_1(x,y_1,y_2,y_3)=g(x)+y_1y_2+y_3,f_2(x,y_1,y_2,y_3)=g(x)+y_1y_2$, then $f_1||f_2$ is a bent function $g(x)+y_1y_2+y_3y_4$, which is obtained by the direct sum of $g$ and $y_1y_2+y_3y_4$.   }
%-- These $2^s$ subfunctions can be specified explicitly and using the second order derivative and  Mesnager easy method (see \cite{SihemTransl}) we can show that such a function is not in MM. Well, the completed MM class is much harder and whether we can show that in an elegant way is an open question - we should avoid these lengthy proofs as in our IEEE paper that uses Rothaus method ! }
%\textbf{\textcolor[rgb]{0.98,0.00,0.00}{Samir:} Probably I can add some computations in this section, regarding the bent functions which are obtained by concatenating functions $f^{(i)}(y,x)=\phi_i(y) \cdot x + g_i(y)$ (but I need to check my computations more carefully!! - so I will add it later on...)}


%The results given in the previous section were mainly based on the direct correspondence between the distribution of zeros and ones of a bent function $g$ which then was used to assign the spectral values of $f$, thus to specify properly an affine subspace $S_f$ which then satisfies the Titsworth theorem. On the other hand, when specifying $S_f$ which is not an affine subspace we are faced with the following problem. The condition that $\sum_{u \in \FB^n} W_f(u)W_f(u+v)=0$, for any $v \in {\FB^n}^*$ cannot be easily handled as before. Even though $W_f(u) \neq 0$ if and only if $u \in S_f$ which also implies that we consider $\sum_{u \in S_f} W_f(u)W_f(u+v)=0$, the set $S_f$ is not closed under addition and $u + v$ may or may not belong to $S_f$.
%
%Furthermore, using a bent function $g$ to specify the signs of the nonzero values in $S_f$ does not ensure that the condition $\sum_{u \in \FB^n} W_f(u)W_f(u+v)=0$ is met, since in general the mapping from $S_f$ to $\FB^{n-s}$ is nonlinear and not well-defined.
%
%In a recent work a useful characterization of plateaued functions represented using the so-called compositional form was given \cite{Secondary}.  Let $\mathfrak{f}:\mathbb{F}^n_2 \rightarrow \mathbb{F}_2$ be an arbitrary Boolean function given in the $CF$-representation as
%\begin{eqnarray}\label{F2}
%\mathfrak{f}(x)=f(H(x))=f(h_1(x),\ldots,h_k(x)),
%\end{eqnarray}
%with the form $f:\mathbb{F}^k_2 \rightarrow \mathbb{F}_2$ and vectorial function $H=(h_1,\ldots,h_k):\mathbb{F}^n_2 \rightarrow \mathbb{F}^k_2$. Since in  (\ref{F2}) one can view $\mathfrak{f}(x)=f(H(x))$ as a composition of two vectorial functions, the result in \cite[Proposition 9.1]{CarletBoolean} implies that the WHT of $\mathfrak{f}$ is given by \cite{Secondary}
%\begin{eqnarray}\label{mainF}
%W_\mathfrak{f}(u)&=&\sum_{x\in \mathbb{F}^n_2}(-1)^{f(h_1(x),\ldots,h_k(x))+ u\cdot x}=2^{-k}\sum_{\omega\in \mathbb{F}^k_2}W_f(\omega)W_{\omega\cdot (h_1,\ldots,h_k)}(u),\;\;\;u\in \mathbb{F}^n_2.
%\end{eqnarray}

\subsection{Construction of nontrivial plateaued functions}\label{sec:nonaffine}

An example of a plateaued function with non-affine Walsh support (which is not partially bent and thus nontrivial), has been given in \cite[Example 3.2]{Secondary}. In difference to the approach taken in the previous section, using again  the $\mathcal{MM}$ class of bent functions, we will  provide a generic construction of nontrivial plateaued functions in the spectral domain. Furthermore, we show that the classes $\mathcal{C}$ and $\mathcal{D}$ \cite{CarletTNC} can also be used for the same purpose, but with certain limitations (Remark \ref{rem:CD}).
Before we present a general framework, we illustrate the main construction idea of nontrivial plateaued functions through an example. %For convenience, for any two sets $A=\{\alpha_1,\ldots,\alpha_{r}\}$ and $B=\{\beta_1,\ldots,\beta_{r}\}$, we define  $A\wr B=\{(\alpha_i,\beta_i): i=1,\ldots,r\}$.
In what follows,  the truth table $T_g$ of a function $g\in \mathcal{B}_n$ will be viewed as a column vector of length $2^n$.
\begin{ex}\label{ex1}
Using Theorem \ref{theo:plateH}, one can construct a plateaued function with non-affine Walsh support as follows. Let us define the  columns of the Walsh support $S_f$ as
$$S_f=T_{\phi_{b_1}}\wr\ldots\wr T_{\phi_{b_4}}\wr T_{\phi_{b_5}}=\mathbb{F}^4_2\wr T_{\phi_{b_5}},$$
where $\phi_{b_5}:\mathbb{F}^4_2\rightarrow \mathbb{F}_2$ is defined as $\phi_{b_5}(x_1,\ldots,x_4)=x_3x_4$.  Let also the dual $\overline{f}^*=g$ be defined as $g(x_1,\ldots,x_4)=x_1x_3+ x_2x_4\in \mathcal{MM}$, see Table \ref{tab:ex1}. Notice that the vectors $x=(x_1,\ldots,x_4)$ are sorted lexicographically.
\begin{table}[H]
\scriptsize
\centering
\begin{tabular}{|r|c|c|}
  \hline
  % after \\: \hline or \cline{col1-col2} \cline{col3-col4} ...
  \makecell{$S_f=\mathbb{F}^4_2\wr T_{\phi_{b_5}}$} & $f^*(\omega_i)=\overline{f}^*(x_i)=g(x_i),$ $x_i\in \mathbb{F}^4_2$\\ \hline
  $\omega_0=(0, 0, 0, 0, 0)$  & 0 \\
  $\omega_1=(0, 0, 0, 1, 0)$ &  0 \\
  $\omega_2=(0, 0, 1, 0, 0)$  & 0 \\
  $\omega_3=(0, 0, 1, 1, 1)$  & 0 \\
  $\omega_4=(0, 1, 0, 0, 0)$  & 0 \\
  $\omega_5=(0, 1, 0, 1, 0)$  & 1 \\
  $\omega_6=(0, 1, 1, 0, 0)$  & 0 \\
  $\omega_7=(0, 1, 1, 1, 1)$  & 1 \\
  $\omega_8=(1, 0, 0, 0, 0)$  & 0 \\
  $\omega_9=(1, 0, 0, 1, 0)$  & 0 \\
$\omega_{10}=(1, 0, 1, 0, 0)$  & 1 \\
$\omega_{11}=(1, 0, 1, 1, 1)$  & 1 \\
$\omega_{12}=(1, 1, 0, 0, 0)$  & 0 \\
$\omega_{13}=(1, 1, 0, 1, 0)$  & 1 \\
$\omega_{14}=(1, 1, 1, 0, 0)$  & 1 \\
$\omega_{15}=(1, 1, 1, 1, 1)$  & 0 \\
  \hline
\end{tabular}
\caption{Values of $f^*$ with respect to $S_f$.}
\label{tab:ex1}
\end{table}
Then, by applying the inverse WHT where  the non-zero Walsh spectral values are given with respect to Table \ref{tab:ex1} and using (\ref{eq:walshX}), we obtain a semi-bent function $f$ given as
$$f(x_1,\ldots,x_5)=x_1 x_3 + x_2 x_4 + x_1 x_2 x_5.$$
It can be verified that $\#\mathfrak{R}=5$, hence $f$ is not a partially bent function.
\end{ex}
\begin{rem} The function $f$ in Example \ref{ex1} does not admit linear structures. This is the consequence of Theorem \ref{theo baselinear} and the fact that $S_f$ contains five linearly independent vectors, thus spanning $\FB^5$.
\end{rem}
In general, the construction of $S_f$ by means of Theorem \ref{theo:plateH} has to satisfy two important conditions. More precisely, $S_f$ must not  be a multiset and its  columns (when $S_f$ is represented as a matrix of size $2^{n-s} \times n$) must lie at bent distance to the dual $\overline{f}^*=g$. The  requirement that $S_f=\mathbb{F}^4_2\wr T_{\phi_{b_5}}$ in Example \ref{ex1} is not a multi-set is apparently fulfilled regardless of the choice of  $\phi_{b_5}$. The second condition is satisfied since the first four columns of $S_f$ are sequences or linear functions (\cite[Lemma 3.1]{Secondary}). Furthermore, $g+\phi_{b_5}$ is a bent function in $n-s$ variables which belongs to the $\mathcal{MM}$ class and therefore $g$ and $\phi_{b_5}$ are at bent distance.

In order to construct a plateaued function $f$ on $\FB^n$ (of any amplitude) we will employ the $\mathcal{MM}$ class of bent functions on $\FB^{n-s}$ and define   its dual as $\overline{f}^*=g\in \mathcal{MM}$. This approach will be later extended  to cover the cases  when $\overline{f}^*=g$ belongs to $\mathcal{C}$ or $\mathcal{D}$ \cite{CarletTNC} classes of bent functions. The main idea behind  our approach is to specify  the Walsh support of $f$ as  $S_f=\mathbb{F}^{n-s}_2\wr S$, for some suitable set $S$ of cardinality $2^{n-s}$ whose elements belong to $\FB^{s}$.

To construct  nontrivial plateaued functions one can employ  bent functions from the  $\mathcal{MM}$ class defined by \begin{equation} \label{def:bentmm}
g(x,y)=x\cdot \psi(y)+ t(y); \;\;x,y\in \mathbb{F}^{(n-s)/2}_2,
\end{equation}
where $\psi$ is an arbitrary permutation on $\mathbb{F}^{(n-s)/2}_2$ and $t \in \mathcal{B}_{(n-s)/2}$ is arbitrary.
\begin{theo}\label{th:Sf}
Let $g:\mathbb{F}^{n-s}_2\rightarrow \mathbb{F}_2$ be defined by (\ref{def:bentmm}).
%an arbitrary bent function in the $\mathcal{MM}$ class, given as $g(x,y)=x\cdot \psi(y)+ t(y)$, $x,y\in \mathbb{F}^{(n-s)/2}_2$, $t\in \mathcal{B}_{n-s}$ ($\psi$ is an arbitrary permutation on $\mathbb{F}^{(n-s)/2}_2$).
For an arbitrary matrix $M\in GL(n-s,\mathbb{F}_2)$ and vector $c\in \mathbb{F}^{n-s}_2$, let $S_f$ be defined as
\begin{eqnarray}\label{Sff}
S_f=T_{\phi_{b_1}}\wr\ldots \wr T_{\phi_{b_{n-s}}} \wr T_{\phi_{b_{n-s+1}}}\wr \ldots\wr T_{\phi_{b_n}}=(c+EM)\wr T_{t_1}\wr\ldots\wr T_{t_s},
\end{eqnarray}
where $t_i=t_i(y)\in \mathcal{B}_{(n-s)/2}$ $(i\in[1,s])$ are arbitrary and $E=\mathbb{F}^{n-s}_2$. Then
%\begin{enumerate}[i)]
%\item
$f:\mathbb{F}^n_2\rightarrow \mathbb{F}_2$ constructed using Theorem \ref{theo:plateH}, with $S_f$ defined by (\ref{Sff}) and taking $\overline{f}^*=g$,  is an $s$-plateaued function.
%\item For any permutation $\sigma:\{1,\ldots,n\}\rightarrow\{1,\ldots,n\}$, it holds that $f$ is an $s$-plateaued function if its dual is $g$ and the Walsh support of $f$ is given by {\bf (This is a very strange formulation ???) }
%$$S^{(\sigma)}_f=T_{\phi_{b_{\sigma(1)}}}\wr\ldots\wr T_{\phi_{b_{\sigma(n)}}}.$$
%where $T_{\phi_{b_1}},\ldots,T_{\phi_{b_n}}$ are defined as in (\ref{Sff}).
%\end{enumerate}
\end{theo}
%{\bf This notation $c+\mathbb{F}^{n-s}_2M$ is slightly confusing and needs to be changed or explained (defined)!!} \\
\proof  By Proposition \ref{prop:Hrow}, we have that the columns corresponding to $c+EM$  are  affine functions in $n-s$ variables. Thus, we may write $\phi_{b_i}=\ell_i$  for some $\ell_i\in \mathcal{A}_{n-s}$, where $i\in[1,n-s]$.
 Since $g$ is a bent function in the $\mathcal{MM}$ class, then for any $v\in \mathbb{F}^k_2$
 $$g(x,y)+ v\cdot (\phi_{b_1}(x,y),\ldots,\phi_{b_n}(x,y))=g(x,y)+ v\cdot(\ell_1(x,y),\ldots,\ell_{n-s}(x,y),t_1(y),\ldots,t_s(y))$$
is also a bent function. By Theorem \ref{theo:plateH}-$(i)$, this  means that $g=\overline{f}^*$ is at bent distance to functions $v\cdot (\phi_{b_1},\ldots,\phi_{b_n})$, for every $v\in \mathbb{F}^n_2$. Thus, $f$ is $s$-plateaued.\qed

%In ordered to compare functions constructed by Theorem \ref{th:Sf} with plateaued functions in Proposition \ref{prop:disspec_affine} derived from the $\mathcal{MM}$ class in the ANF domain, we consider the WHT of functions given by (\ref{eq:mmdisjoint}). For $(u,v)\in \mathbb{F}^s_2\times \mathbb{F}^k_2$ we have:
%\begin{eqnarray}\label{eq:MMSf}\nonumber
%W_{f^{(i)}}(u,v)&=&\sum_{y\in \mathbb{F}^k_2}(-1)^{g_i(y)+v\cdot y}\sum_{x\in \mathbb{F}^s_2}(-1)^{x\cdot(\phi_i(y)+u)}\\
%&=&\left\{\begin{array}{cc}
%              0, & \phi^{-1}(u)=\emptyset \\
%              2^s(-1)^{g_i(\phi^{-1}(u))+v\cdot \phi^{-1}(u)}, & \phi^{-1}(u)\neq\emptyset
%   \end{array}
%\right..
%\end{eqnarray}
%From (\ref{eq:MMSf}) we have that the Walsh support of $f^{(i)}$ is given as $S_{f^{(i)}}=Im(\phi)\times \mathbb{F}^k_2=E_i\times \mathbb{F}^k_2$. However, it is not difficult to see that the Walsh supports given by (\ref{Sff})  actually include these particular forms of Walsh supports. For instance, notice that the Walsh support of $f$ in Example \ref{ex1} can be written as
%$$S_f=\mathbb{F}^2_2\times \left(
%                             \begin{array}{ccc}
%                               0 & 0 & 0 \\
%                               0 & 1 & 0 \\
%                               1 & 0 & 0 \\
%                               1 & 1 & 1 \\
%                             \end{array}
%                           \right).
%$$
%Moreover, since the dual of $f^{(i)}$ in (\ref{eq:MMSf}) is from the $\mathcal{MM}$ class of bent functions, we conclude that Theorem \ref{th:Sf} provides functions which belong to the $\mathcal{MM}$ class of plateaued functions.{\bf This is not clear ?? You mention the dual of (19) and then conclude that Theorem 4.1 gives MM functions ??}
%%%%%%%%%%%%%%%%%%%%%%%%%%%%%%%%%
%\begin{rem}
%Notice that the only case when Theorem \ref{th:Sf} does not provide plateaued functions from the $\mathcal{MM}$ class is when at least one $t_i$ functions depends on both variables $x,y\in \mathbb{F}^{(n-s)/2}_2$ and not only on $y$. {\bf This is very confusing ?? Maybe not true at all but anyway hard to follow ... It might be the case that these functions are still equivalent and in MM class ?? Maybe we delete the above part completely }
%\end{rem}
%{\bf In general, this Section is very heavy and hard to follow !! It looks like we are trying to do something at every price and it just becomes too complicated and lacks precision. The main question is given (19) where $E_i$ are arbitrary subsets of cardinality $2^k$ how do we argue in smart and concise way that the spectral method is more general .... My suggestion is not to overload this paper as we did with CF paper but rather try to give fewer results which are simpler to follow and have a structural flow throughout the paper ! }
%However, the more important fact is that Theorem \ref{th:Sf} also gives the possibility of specifying new plateaued functions which do not belong to the  $\mathcal{MM}$ class of plateaued functions. %For instance, this can be achieved by defining the dual function $f^*=g$ to be a bent functions from either $\mathcal{C}$ or $\mathcal{D}$ classes of bent functions, and setting suitable functions $t_i$ in (\ref{Sff}). We provide the following result.
Using the same arguments as in the proof of Theorem \ref{th:Sf}, one can easily employ $\mathcal{C}$ and $\mathcal{D}$ classes of bent functions in order to construct semi-bent functions ($s=1$) as follows.
\begin{theo}\label{th:CD}
Let $g(x,y)=x\cdot \psi(y)$, $x,y\in \mathbb{F}^{n/2}_2$, be a bent function, $n$ is even.
% given as $g(x,y)=x\cdot \psi(y)$, $x,y\in \mathbb{F}^{(n-1)/2}_2$ ($\psi$ is a permutation on $\mathbb{F}^{(n-1)/2}_2$).
For an arbitrary matrix $M\in GL(n,\mathbb{F}_2)$ and vector $c\in \mathbb{F}^{n}_2$, let $S_f=(c+EM)\wr T_{\mu},$ where $E=\mathbb{F}^{n}_2$ is ordered lexicographically and $\mu\in \mathcal{B}_{n}$. Then:
 \begin{enumerate}[i)]
\item  If for subspaces $E_1,E_2\subset \mathbb{F}^{n/2}_2$ we have $\mu(x,y)=\phi_{E_1}(x)\phi_{E_2}(y)$ and $\psi$ satisfies the equality $\psi(E_2)=E^{\perp}_1$, then the function $f:\mathbb{F}^n_2\rightarrow \mathbb{F}_2$ with Walsh support $S_f$ and the dual $\overline{f}^*=g$ is a semi-bent function.
\item For a subspace $L\subset \mathbb{F}^{n}_2$ let $\mu(x,y)=\phi_L(x)$. If $\psi^{-1}(v+L^{\perp})$ is an affine subspace for all $v\in \mathbb{F}^{n}_2$, then the function $f:\mathbb{F}^n_2\rightarrow \mathbb{F}_2$ with Walsh support $S_f$ and the dual $\overline{f}^*=g$ is a semi-bent function.
\end{enumerate}
\end{theo}
%{\bf For several reasons $E_1$ and $E_2$ cannot be subsets but strictly linear subspaces (possibly affine) !! So, in this context we need to be clear what does this implies ! I cannot see the benefits of this result at all ! We should focus on (19) and either prove or state that the dual of $f_i$ in Proposition 4.1 is an MM function whereas it is not necessarily true for our spectral design. We must avoid such statements as Theorem 4.2 - it is too long and complicated .... Remark bellow should be removed it simply is an overload - again less results but precise and stated in a nice way ... }
\begin{rem}\label{rem:CD}
Note that Theorem \ref{th:CD} can be modified so that it provides $s$-plateaued functions with $s>1$, similarly as Theorem \ref{th:Sf}. However, in that case (setting $g$ as in Theorem \ref{th:CD}) in relation (\ref{Sff}) we need to set functions $t_i$ to be equal to $t_i(x,y)=\phi_{E_1}(x)\phi_{E_2}(y)$ or $t_i(x,y)=\phi_L(x)$ for all $i\in[1,s]$. The other possibility is to find a space of functions $t_i$ which have these forms and for which the permutation $\psi$ satisfies the necessary conditions imposed by the definitions of $\mathcal{D}$ and $\mathcal{C}$ classes (since the dual $\overline{f}^*=g$ has to be at bent distance to $\Phi_f$).
\end{rem}
%Note that $S_f$ is not a multi-set regardless of the choice of functions $t_1,\ldots,t_s$ due to the fact that $c+\mathbb{F}^{n-s}_2M$ is not a multiset.\\
%{\bf This notation $c+\mathbb{F}^{n-s}_2M$ is slightly confusing and needs to be changed or explained !!}
%
%What remains to be proven is that $S_d$ is not a multi-set in $\mathbb{F}^k_2$, i.e., that $S_d$ is a proper subset in $\mathbb{F}^k_2$. Let $\omega_i=(\omega'_i,\omega''_i)$ and $\omega_j=(\omega'_j,\omega''_j)$ be two arbitrary vectors from $S_d=\{\omega_0,\ldots,\omega_{2^{k-s}-1}\}\subset \mathbb{F}^k_2$ with $i\neq j$, where $\omega'_i,\omega'_j\in \mathbb{F}^{k-s}_2$. If we assume that $\omega_i=\omega_j$, it means that $\omega'_i=\omega'_j$ and $\omega''_i=\omega''_j$. However, $i\neq j$ means that we can not have that  $\omega'_i=\omega'_j$, since $\omega'_i,\omega'_j\in \mathbb{F}^{k-s}_2$. Consequently, $S_d$ is not a multi-set, regardless of selected functions $g_i(y)\in \mathcal{B}_{k-s}$.
%
%In order to obtain $s$-plateaued function $f$ with Walsh support $S_f$ and dual $g$, we construct Walsh spectrum as in relation (\ref{eq:walshX}) and then the inverse WHT formula given in relation (\ref{WHT}) provides the function $f$.
%$ii)$ Any reordering of the columns of $S_f$ does not violate the property that $S_f$ is not a multi-set and the bent distance to $g$ is preserved. \qed
% the part $c+\mathbb{F}^{n-s}_2M$ in $S_f$ will not violate the property of $S_f$ not being a multiset, the proof is completed.\qed
%%%%%%%%%%%%%%%%%%%%%%%%%%%%%%%%%%%%%%%%%%%%%%%%%%%%%%%%%%%%%%%
%\begin{rem}\label{grem}
%Notice in Theorem \ref{th:Sf}, if we have that functions $t_i(y)\in \mathcal{B}_{n-s}$ are linear, then $S_f$ is an affine subspace (Proposition \ref{prop:Hrow} and \cite[Lemma 3.1]{Secondary}). This means that by a simple choice of at least one nonlinear functions $t_i(y)$ we have that $S_f$ will not be an affine subspace of $\mathbb{F}^n_2$. Furthermore, by  a proper selection of $t_i$ it can be ensured that the set $S_f$ contains the basis of $\FB^n$ so that $f$ in Theorem \ref{th:Sf} does not admit linear structures.
%%%%%%%%%%%%%%%%%%%%%%%%%%%%%%%%%%%%%%%%%%%%%%%%%%%%%%%%%%%%%%%%
%\textcolor[rgb]{0.98,0.00,0.00}{Also note that by choosing a suitable permutation $\psi$ in $g(x,y)=x\cdot \psi(y)$ ($t=0$) and functions $t_i$ in $S_f$ ($i\in [1,s]$) one can construct plateaued functions $f$ by utilizing $\mathcal{D}$, $\mathcal{C}$ \cite{CarletTNC} and $\mathcal{PS}$ \cite{Dillon} classes of bent functions.}
 %Notice that the dual $g$ can be taken form other classes of bent functions such as $\mathcal{D}$ and $\mathcal{C}$ \cite{CarletTNC} classes of bent functions but the conditions are then more complicated.
%\end{rem}
%\textcolor[rgb]{0.98,0.00,0.00}{Using the same arguments as in the proof of Theorem \ref{th:Sf}, one can easily employ $\mathcal{D}$ and $\mathcal{C}$ \cite{CarletTNC} classes of bent functions in order to construct semi-bent functions ($s=1$) as follows.
%\begin{theo}\label{th:CD}
%Let $g:\mathbb{F}^{n-1}_2\rightarrow \mathbb{F}_2$ be a bent function given as $g(x,y)=x\cdot \psi(y)$, $x,y\in \mathbb{F}^{(n-1)/2}_2$ ($\psi$ is a permutation on $\mathbb{F}^{(n-1)/2}_2$). For an arbitrary matrix $M\in GL(n-1,\mathbb{F}_2)$ and vector $c\in \mathbb{F}^{n-1}_2$, let $S_f$ be defined as $S_f=(c+EM)\wr T_{\mu},$ where $E=\mathbb{F}^{n-1}_2$ is ordered lexicographically and $\mu\in \mathcal{B}_{n-1}$. Then, if for subsets $E_1,E_2\subset \mathbb{F}^{(n-1)/2}_2$ (or $L\subset \mathbb{F}^{n-1}_2$) we have that $\mu(x,y)=\phi_{E_1}(x)\phi_{E_2}(y)$ (or $\mu(x,y)=\phi_L(x)$) and $\psi$ satisfies the equality $\psi(E_2)=E^{\perp}_1$ (or $\psi^{-1}(v+L^{\perp})$ is an affine subspace for all $v\in \mathbb{F}^{n-1}_2$), then the function $f:\mathbb{F}^n_2\rightarrow \mathbb{F}_2$ with Walsh support $S_f$ and the dual $f^*=g$ is a semi-bent function.
%\end{theo}
%\begin{rem}\label{rem:CD}
%Note that Theorem \ref{th:CD} can be modified so that it provides $s$-plateaued functions with $s>1$, similarly as Theorem \ref{th:Sf}. However, in that case (setting $g$ as in Theorem \ref{th:CD}) in relation (\ref{Sff}) we need to set functions $t_i$ to be equal to $t_i(x,y)=\phi_{E_1}(x)\phi_{E_2}(y)$ or $t_i(x,y)=\phi_L(x)$ for all $i\in[1,s]$. The other possibility is to find a space of functions $t_i$ which have these forms and for which the permutation $\psi$ satisfies the necessary conditions imposed by the definitions of $\mathcal{D}$ and $\mathcal{C}$ classes (due to the fact that the dual $f^*=g$ has to be at bent distance to $\Phi_f$).
%\end{rem}
%}
The above approach  can be easily generalized using  the following result.
\begin{theo}\label{theo:vectorial}
Let $H=(h_1,\ldots,h_{\lambda}):\mathbb{F}^{n-s}_2\rightarrow \mathbb{F}^{\lambda}_2$, where  $\lambda\leq (n-s)/2$ and $n-s$ is even, be a vectorial bent function. Let the dual of $f$  be specified as $\overline{f}^*=g=h_i$, for some fixed $i\in[1,\lambda]$. In addition, for an arbitrary matrix $M\in GL(n-s,\mathbb{F}_2)$ and vector $c\in \mathbb{F}^{n-s}_2$, let $S_f$ be defined as
$$S_f=(c+EM)\wr T_{t_1}\wr\ldots\wr T_{t_m} \wr T_{h_1}\wr\ldots\wr T_{h_r},$$
where $t_i\in \mathcal{A}_{n-s}$, $m+r=s$ with $r\leq \lambda-1$ and $E=\mathbb{F}^{n-s}_2$ is ordered lexicographically. In addition, $h_1,\ldots,h_{i-1},h_{i+1},\ldots, h_r \neq g$.
% (functions $h_j$ in $S_f$ are not equal to $g=h_i$).
Then,  $f:\mathbb{F}^n_2\rightarrow \mathbb{F}_2$, whose  Walsh support is $S_f$ and the dual $\overline{f}^*=g=h_i$,
%, or $S^{(\sigma)}_f=T_{\phi_{b_{\sigma(1)}}}\wr\ldots\wr T_{\phi_{b_{\sigma(n)}}}$ ($\sigma$ is permutation on $\{1,\ldots,n\}$), and dual $f^*=g$,
is an $s$-plateaued function.%{\bf This is contradictory, first we have $g \neq h_i$ and then $g=h_i$, also which $i$ here ??}
\end{theo}
%\begin{rem}
%It is not difficult to see that disjoint plateuaed functions to those constructed by Corollary \ref{cor:vectorial} are obtained by taking a Walsh support as
%\end{rem}
By Proposition \ref{prop:Hrow}, if the affine subspace $c+EM$ (with $E=\mathbb{F}^{n-s}_2$)
%ordered lexicographically, when
is written as a matrix of size $2^{n-s}\times (n-s)$, then its columns  correspond to linear/affine functions on $\mathbb{F}^{n-s}_2$. An important  question in this context is whether there exists (many) plateaued functions with the property that $S_f=T_{\phi_{b_1}}\wr \ldots \wr T_{\phi_{b_n}}$ ($\phi_{b_i}$ defined by (\ref{RSf})) contains strictly less than $n-s$ columns which correspond to affine functions on $\mathbb{F}^{n-s}_2$.
%This leads to the following problem.
\begin{op}
Provide (generic) constructions of $s$-plateaued functions in $n$ variables whose Walsh support, when written as a matrix of order $2^{n-s}\times n$, contains more than $s$ columns which correspond to nonlinear  functions on $\mathbb{F}^{n-s}_2.$
\end{op}
\begin{rem}
Recall that  $S_f\subset \mathbb{F}^n_2$ of cardinality $2^{n-s}$ can be used as a Walsh support for an $s$-plateaued function if and only if
%it is not a multi-set and if
there exists a function  $f^*\in \mathcal{B}_{n-s}$ of weight $2^{n-s-1}\pm 2^{(n-s)/2-1}$ which is at bent distance to all  linear combinations of columns of $S_f$ (Theorem \ref{theo:plateH}).
\end{rem}

\subsection{Designing disjoint spectra nontrivial plateaued functions}


The problem of finding  nontrivial  $s$-plateaued functions with disjoint spectra is intrinsically more difficult than for trivial plateaued functions.
% constructed by Theorem \ref{th:Sf}, the question is whether the Walsh supports of these functions are pairwise disjoint.
 The reason is   that taking $S_f$ not to be affine subspace than the cosets of $S_f$ are in general not disjoint, that is, $(v+ S_f)\cap S_f \neq \emptyset$ for some $v\in\mathbb{F}^n_2$.
 % if we take in $S_f$ (in relation (\ref{Sff})) that at least one function $t_i$ is nonlinear, then not for all vectors $v\in\mathbb{F}^n_2$ it will hold that $(v+ S_f)\cap S_f=\emptyset$, since $S_f$ is not an affine subspace.
Nevertheless, the class of nontrivial plateaued functions given in Theorems \ref{th:Sf}-\ref{theo:vectorial} also gives rise to disjoint spectra nontrivial plateaued functions.
\begin{lemma}\label{lemma:Q}
Let $f$ be an $s$-plateaued functions constructed by any of Theorems \ref{th:Sf}-\ref{theo:vectorial} with $S_f \subset \FB^n$, where $\#S_f=2^{n-s}$. Define the set $Q$ of cardinality $2^s$ by
\begin{eqnarray}\label{Q}
Q=\{c\cdot (b_{n-s+1},b_{n-s+2},\ldots,b_{n}):c\in \mathbb{F}^s_2\},
\end{eqnarray}
where $b_i=(0,\ldots,0,1,0,\ldots,0)\in \mathbb{F}^n_2$ ($1$  at the $i$-th position).  Then, the sets $q+S_f$, $q\in Q$, partition the space $\mathbb{F}^n_2$, i.e., $$\mathbb{F}^n_2=\bigcup_{q\in Q} (q + S_f), \;\;  (q_i+ S_f)\cap(q_j+ S_f)=\emptyset \;\; \textnormal{ for } q_i\neq q_j,  \;\; q_i,q_j\in Q.$$
\end{lemma}
\proof In order to prove the statement for all mentioned theorems, we assume that $S_f$ is given as in (\ref{Sff}) and additionally  we take that $t_i=t_i(x,y)$ (thus depending also on $x$). Then, such a Walsh support $S_f$ covers all the cases in Theorem \ref{th:Sf}-\ref{theo:vectorial}.

It is sufficient to prove that $(q_i+ S_f)\cap(q_j+ S_f)=\emptyset$, for $q_i\neq q_j$ ($q_i,q_j\in Q$), since $\# Q=2^s$ and $\# S_f=2^{n-s}$ will then imply that $\mathbb{F}^n_2=\bigcup_{q\in Q} (q+ S_f)$.  Assume, on contrary,  that for some $q_i\neq q_j$ there exists  $w\in (q_i+ S_f)\cap(q_j+ S_f)$, where $S_f$ is given by relation (\ref{Sff}).
% and $q_i=(q^{(i)}_1,\ldots,q^{(i)}_s)$ and $q_j=(q^{(j)}_1,\ldots,q^{(j)}_s)$ are two arbitrary vectors from $Q$. If we assume that there exists a vector $w\in (q_i+ S_f)\cap(q_j+ S_f)$, where $S_f$ is given by relation (\ref{Sff}),
Then, denoting  $q_i=(q^{(i)}_1,\ldots,q^{(i)}_s)$ and $q_j=(q^{(j)}_1,\ldots,q^{(j)}_s)$, for some $x',x'',y',y''\in \mathbb{F}^{(n-s)/2}_2$ and $M\in GL(n,\mathbb{F}^{n-s}_2)$ it holds that
\begin{eqnarray*}
\left\{\begin{array}{c}
         w=(c+(x',y')M, t_1(x',y')+ q^{(i)}_1,\ldots,t_s(x',y')+ q^{(i)}_s)\in q_i+ S_f \\
         w=(c+(x'',y'')M, t_1(x'',y'')+ q^{(j)}_1,\ldots,t_s(x'',y'')+ q^{(j)}_s)\in q_j+ S_f
       \end{array}
\right..
\end{eqnarray*}
This implies that $x'=x''$, $y'=y''$ and $t_r(x',y')+ q^{(i)}_r=t_r(x'',y'')+ q^{(j)}_r$, for all $r\in[1,s]$, and consequently $q_i=q_j$. This contradicts the assumption that $q_i \neq q_j$ (since $t_r(x',y')=t_r(x'',y'')$), which completes the proof.\qed
%{\bf I would remove both Proposition 4.2 and Remark 4.5 !!!  Proposition 4.2 talks about two functions and we need $2^s$ functions ... I find it a small distracting result with no significance and its statement is ENORMOUS and confusing !! We need to say one sentence for clarity "We now employ Lemma 4.1 to provide generic design method of disjoint spectra nontrivial plateued functions of maximal cardinality !!!} \\
%\textcolor[rgb]{0.98,0.00,0.00}{\begin{prop}\label{prop:Q2}
%Let $f:\mathbb{F}^n_2\rightarrow \mathbb{F}_2$ be an $s$-plateaued function with Walsh support $S_f=(c+EM)\wr S\subset \mathbb{F}^n_2$ ($\#S_F=2^{n-s}$, $E=\mathbb{F}^{n-s}_2$ is ordered lexicographically, $M\in GL(n-s,\mathbb{F}_2)$), where $S$ (when written as a matrix) has columns which correspond to nonlinear functions on $\mathbb{F}^{n-s}_2$. Suppose that for arbitrary nonzero vector $\delta\in \mathbb{F}^s_2$ and $\beta=(\textbf{0}_{n-s},\delta)\in \mathbb{F}^n_2$, the set $S_d\subset\mathbb{F}^n_2$ is given as $S_d=(c+EM)\wr(\beta+S)$. Then an $s$-plateaued function $d:\mathbb{F}^n_2\rightarrow \mathbb{F}_2$ constructed by Theorem \ref{theo:plateH} with Walsh support $S_d$ and dual $d^*=f^*$ is disjoint spectra function to $f$.
%\end{prop}
%\proof Using the same arguments as in the proof of Lemma \ref{lemma:Q}, it is not difficult to see that $S_f\cap S_d=\emptyset$. Consequently, setting $d^*=f^*$ in Theorem \ref{theo:plateH} we obtain a nontrivial $s$-plateaued function which is disjoint spectra to $f$.\qed}
%\begin{rem}
%Clearly, if the plateaued function $f$ has the Walsh support given by Theorem \ref{th:Sf}-(ii), then for any permutation $\sigma:\{1,\ldots,n\}\rightarrow \{1,\ldots,n\}$ it holds that the direct sum  $S^{(\sigma)}_d \oplus Q^{(\sigma)}=\mathbb{F}^n_2$,
%% (i.e., $S^{(\sigma)}_d\oplus  Q^{(\sigma)}=\mathbb{F}^n_2$),
%where $Q^{(\sigma)}=\{c\cdot (b_{\sigma(n-s+1)},\ldots, b_{\sigma(n)}):c\in \mathbb{F}^s_2\}$.
%\end{rem}
Now employing Theorems \ref{th:Sf}-\ref{th:CD} and Lemma \ref{lemma:Q} we derive a construction method of disjoint spectra nontrivial plateaued functions (possibly without linear structures). Additionally, we show that these functions can be efficiently utilized for construction of bent functions.
%
%bent functions $\mathfrak{f}:\mathbb{F}^m_2\rightarrow \mathbb{F}_2$ based on nontrivial disjoint spectra plateaued functions.
%
%{\bf We should reformulate this sentence and first say that it gives nontrivial disjoint spectra plateaued functions and then mention (maybe separately) that these functions can be used to construct a bent function - though we have said that several times ?? What is much better here in the second part, since Example 4.2 does not work well, is to prove that these functions are without linear structures !! We have to be careful here since MM class admits linear structures, so we get (easily) these nontrivial plateaued functions without linear structures, concatenate these and get a bent function in MM class which has linear structures ?! Probably best to skip Example 4.2 since somehow we have done nothing just going around in circles ...}
\begin{theo}\label{th:gencon}
Let $V$ be a $n$-dimensional subspace of $\mathbb{F}^{m}_2$ ($m$ even), such that $W\oplus V=\mathbb{F}^m_2$,
where $W=\{a_0,a_1,\ldots,a_{2^{m-n}-1}\}$ and $Q=\{q_0,q_1,\ldots,q_{2^s-1}\}$ is defined by (\ref{Q}) ($s=m-n$, $q_0=\textbf{0}_n,$ $a_0=\textbf{0}_k$). Let $f:\mathbb{F}^n_2\rightarrow \mathbb{F}_2$ be  $s$-plateaued  constructed by Theorem \ref{th:Sf} or \ref{th:CD}. Then:
\begin{enumerate}[i)]
\item Let  $f_0,\ldots,f_{2^{s}-1}:\mathbb{F}^n_2\rightarrow \mathbb{F}_2$ ($f_0=f$) be $s$-plateaued functions constructed by Theorem \ref{theo:plateH}, with Walsh supports $S_{f_i}=S_f+ q_i\subset \mathbb{F}^n_2$ and bent duals $\overline{f}^*_i$ taken from the $\mathcal{MM}$ or $\mathcal{C}$/$\mathcal{D}$ class  if Theorem \ref{th:Sf} or Theorem \ref{th:CD} is used, respectively. Then $f_0,\ldots,f_{2^{s}-1}$ are pairwise disjoint spectra functions.
\item $f_1,\ldots,f_{2^s-1}$ do not admit linear structures if and only if $f$ does not admit linear structures, that is when $dim(S_{f})=n$.
\item Suppose that the restrictions of a function $\mathfrak{f}:\mathbb{F}^m_2\rightarrow \mathbb{F}_2$ are defined by
 $$\mathfrak{f}|_{a_i+ V}(a_i+x)=f_i(x),\;\;\;i\in[0,2^{m-n}-1],\; x\in V,$$
 where $f_i$ are constructed as in $i)$. Then, $\mathfrak{f}$ is a bent function on $\mathbb{F}^m_2$.
 \end{enumerate}
\end{theo}
\proof $i)$ By Lemma \ref{lemma:Q}, for any $i\neq j$ ($i,j\in [0,2^s-1]$) it holds that $S_{f_i}\cap S_{f_j}=\emptyset$, and thus $f_0,\ldots,f_{2^s-1}$ are disjoint spectra functions.

$ii)$ By Corollary \ref{theo:baselinear2},  $f_i$ does not admit linear structures if and only if $\dim(S_{f_i})=\dim(S_f)=n$.

$iii)$ For an arbitrary $u\in \mathbb{F}^n_2$,  the WHT of $\mathfrak{f}$ is given by
\begin{eqnarray*}\label{maineq}\nonumber
W_\mathfrak{f}(u)=\sum_{a_i\in W}(-1)^{u\cdot a_i}\sum_{x\in V}(-1)^{f_i(x)+ u\cdot x}=\sum_{a_i\in W}(-1)^{u\cdot a_i}W_{f_i}(\vartheta_u),
\end{eqnarray*}
where $\vartheta_u\in \mathbb{F}^{n}_2$ is a vector for which it holds that $(\vartheta_u\cdot y_0,\ldots,\vartheta_u\cdot y_{2^n-1})=(u\cdot x_0,\ldots,u\cdot x_{2^n-1})$, where $y_i\in \mathbb{F}^n_2$, $x_i\in V\subset\mathbb{F}^m_2$. Recall that by Proposition \ref{prop:Hrow} for lexicographically ordered space $V=\{x_0,\ldots,x_{2^n-1}\}$ it holds that $(u\cdot x_0,\ldots,u\cdot x_{2^n-1})$ is a sequence of a linear function.
%
%defines the linear function $\ell:\mathbb{F}^n_2\rightarrow \mathbb{F}_2$, given as $\ell(x)=\vartheta_u\cdot x$ ($x\in \mathbb{F}^{n}_2$), {\bf (This is very very unclear ??? Introducing $\ell$ and other notation  is just too much ?? Why not use $W_{f_i}(u)$ simply ?? )} whose truth table is given as $\{u\cdot x:x\in V\}$ when the subspace $V$ is ordered lexicographically (Proposition \ref{prop:Hrow}).
Consequently, using the fact that $f_i$ are disjoint spectra $s$-plateaued functions on $\mathbb{F}^n_2$, we have that $W_{\mathfrak{f}}(u)=\pm 2^{\frac{n+s}{2}}=\pm 2^{\frac{m}{2}}$, which completes the proof. \qed

%\textbf{\textbf{Samir:} For some reason I am not sure whether $q_i+S_f$ will preserve the linear independence (of vectors in $S_f$)?? I have written above that $dim(S_{f_i})=dim(S_f)$, but when we check results in linear algebra, all of them are related to linear invertible maps. But what with shifting a non-affine subspace? Intuitively it must be the case that it holds, as I have written, but I am still not sure why :) (proving it in a classical way does not go so smooth, or there is some trick?)}

%\begin{rem}
%Note that in Theorem \ref{th:gencon} we define $f_i$ with respect to relation (\ref{DPL}) by setting the subspace $V\subset \mathbb{F}^m_2$ to be ordered lexicographically (and thus in (\ref{DPL}) we have $v=\textbf{0}_m$ and $M=I_{m\times m}$).
%\end{rem}
%{\bf In general, we need to avoid having dozens of remark !!  Having too many may indicate that we are clumsy in our statements and need to clarify further each now and then !!} \\
%The construction details are illustrated in the following example.
%\begin{ex}\label{ex2}
%Let us construct a bent function $\mathfrak{f}$  on $\mathbb{F}^8_2$ ($m=8$), where the restrictions of $\mathfrak{f}$ are defined on cosets of the $6$-dimensional subspace $V\subset\mathbb{F}^8_2$ ($n=6$, and thus $s=m-n=2$) spanned by vectors $b_3,\ldots,b_8$ (where $\{b_1,\ldots,b_8\}$ is a  canonical basis of $\FB^8$), i.e.,
%$$V=span\{b_3,\ldots,b_8\}=\textbf{0}_{2}\times\mathbb{F}^6_2.$$
%Clearly, the space $W$ is given as $W=span\{b_1,b_2\}=\{a_0,\ldots,a_{3}\}$ ($b_1,b_2\in \mathbb{F}^8_2$) and $W\oplus V=\mathbb{F}^8_2.$ The restrictions $\mathfrak{f}|_{a_i+ V}$ are viewed as functions defined on $\mathbb{F}^6_2$ (being defined on the cosets of $V$), and they need to be disjoint spectra $2$-plateaued functions (i.e., semi-bent functions).
%
%Using Theorem \ref{th:Sf}-(i) as a design method,  we first specify the Walsh support $S_{f}$ ($f=f_0$) of size $2^{n-s}=2^4$ as $S_{f}=\mathbb{F}^4_2\wr T_{t_1}\wr T_{t_2}$, where we choose two (different) arbitrary functions \textcolor[rgb]{0.98,0.00,0.00}{$t_1$ and $t_2$ to be defined on $\mathbb{F}^4_2$ as }
%$$t_1(x_1,x_2,x_3,x_4)=x_3 x_4+ 1,\;\;\;t_2(x_1,x_2,x_3,x_4)=x_3 x_4+ x_4.$$
%%
%% Hence, let us consider
%%$$S_{f|_{V}}=S_{f|_{a_0+ V}}=\mathbb{F}^4_2\wr T_{g_1}\wr T_{g_2}=
%%                                                                         \begin{array}{cccccc}
%%                                                                           T_{\phi_{e_1}} & T_{\phi_{e_2}} & T_{\phi_{e_3}} & T_{\phi_{e_4}} & T_{g_1} & T_{g_2} %\vspace{2mm} \\
%%                                                                           0 & 0 & 0 & 0 & 1 & 0 \\
%%                                                                           0 & 0 & 0 & 1 & 1 & 0 \\
%%                                                                           0 & 0 & 1 & 0 & 1 & 1 \\
%%                                                                           0 & 0 & 1 & 1 & 0 & 0 \\
%%                                                                           0 & 1 & 0 & 0 & 1 & 0 \\
%%                                                                           0 & 1 & 0 & 1 & 1 & 0 \\
%%                                                                           0 & 1 & 1 & 0 & 1 & 1 \\
%%                                                                           0 & 1 & 1 & 1 & 0 & 0 \\
%%                                                                           1 & 0 & 0 & 0 & 1 & 0 \\
%%                                                                           1 & 0 & 0 & 1 & 1 & 0 \\
%%                                                                           1 & 0 & 1 & 0 & 1 & 1 \\
%%                                                                           1 & 0 & 1 & 1 & 0 & 0 \\
%%                                                                           1 & 1 & 0 & 0 & 1 & 0 \\
%%                                                                           1 & 1 & 0 & 1 & 1 & 0 \\
%%                                                                           1 & 1 & 1 & 0 & 1 & 1 \\
%%                                                                           1 & 1 & 1 & 1 & 0 & 0 \\
%%                                                                         \end{array},
%%$$
%%where $g_1(x_1,x_2,x_3,x_4)=x_3 x_4+ 1$, $g_2(x_1,x_2,x_3,x_4)=x_3 x_4+ x4$ and $\mathbb{F}^4_2=T_{\phi_{e_1}} \wr\ldots\wr T_{\phi_{e_4}}.$
%The dual $f^*$ (viewed as a function on $\mathbb{F}^4_2$) we define as
%$$f^*(x_1,x_2,x_3,x_4)=x_1x_4 +  x_2x_3 + x_3x_4 + x_3.$$
%%and its sequence is given as $\chi_{f^*|_{V}}=(1, 1, -1, 1, 1, 1, 1, -1, 1, -1, -1, -1, 1, -1, 1, 1).$
%
%Now, constructing the Walsh spectrum by (\ref{eq:walshX}) and inverting by (\ref{WHT}) we obtain the plateaued function $f$ given as
%$$f(x_1,\ldots,x_6)=x_1 x_4+ x_5 + x_2 (1 + x_3 + x_6 + x_1 (1 + x_5 + x_6)).$$
%Consequently, we have constructed the restriction of $\mathfrak{f}$ which corresponds to the space $V$, as $\mathfrak{f}|_{V}=f_0=f$, since $a_0=\textbf{0}_m\in W$.
%
%Recall that by relation (\ref{Q}) the space $Q$ is defined as $Q=span\{b_5,b_6\}=\textbf{0}_{2}\times \mathbb{F}^2_2$ ($b_5,b_6\in \mathbb{F}^6_2$), and thus we have that other restrictions of $\mathfrak{f}$ are obtained from the following Walsh supports and duals (which are arbitrary bent functions from the $\mathcal{MM}$-class)
%\begin{eqnarray*}
%\left\{\begin{array}{c}
%         S_{f_1}=S_{\mathfrak{f}|_{a_1+ V}}=S_{\mathfrak{f}|_{(0,1,\textbf{0}_6)+ V}}=q_1+ S_{f}=(\textbf{0}_4,0,1)+ S_{f}=\mathbb{F}^4_2\wr T_{t_1}\wr T_{t_2+ 1}, \\
%         \chi_{f^*_1}=(-1, -1, 1, 1, -1, 1, 1, -1, -1, -1, -1, -1, -1, 1, -1, 1) \\
%         \mathfrak{f}|_{(0,1,\textbf{0}_6)+ V}(x_1,\ldots,x_6)=1 + x_2 x_4 + x_5 + x_6 + x_1 (1 + x_3 + x_6 + x_2 (x_5 + x_6))
%       \end{array}
%\right.,\\\\
%\left\{\begin{array}{c}
%         S_{f_2}=S_{\mathfrak{f}|_{a_2+ V}}=S_{\mathfrak{f}|_{(1,0,\textbf{0}_6)+ V}}=q_2+ S_{f}=(\textbf{0}_4,1,0)+ S_{f}=\mathbb{F}^4_2\wr T_{t_1+ 1}\wr T_{t_2}, \\
%         \chi_{f^*_2}=(1, 1, 1, -1, 1, -1, 1, 1, 1, -1, -1, -1, 1, 1, -1, 1) \\
%         \mathfrak{f}|_{(1,0,\textbf{0}_6)+ V}(x_1,\ldots,x_6)=x_1 (x_3 + x_6) + x_2 (1 + x_3 + x_4 + x_5 + x_1 (1 + x_5 + x_6))
%       \end{array}
%\right., \\\\
%\left\{\begin{array}{c}
%         S_{f_3}=S_{\mathfrak{f}|_{a_3+ V}}=S_{f|_{(1,1,\textbf{0}_6)+ V}}=q_3+ S_{f}=(\textbf{0}_4,1,1)+ S_{f}=\mathbb{F}^4_2\wr T_{t_1+ 1}\wr T_{t_2+ 1}, \\
%         \chi_{f^*_3}=(1, 1, 1, -1, 1, -1, 1, 1, -1, -1, 1, -1, -1, 1, 1, 1) \\
%         \mathfrak{f}|_{(1,1,\textbf{0}_6)+ V}(x_1,\ldots,x_6)=x_3 + x_1 x_3 + x_1 x_6 + x_2 (1 + x_4 + x_5 + x_6 + x_1 (1 + x_5 + x_6))
%       \end{array}
%\right..
%\end{eqnarray*}
%{\bf Enes new comment: Is it possible to get something more complicated in this context - the choice of $t_1$ and $t_2$ are arbitrary so that MM class is not immediately recognized, see also comment below ?!} \\
%%{\bf Enes comment: All these functions above are linear for fixed $x_1,x_2$ !! Probably this is just the MM class though notation here does not allow to identify it immediately ... We need to think a bit more what we are essentially doing ... I will come with some further comments later .. }\\\\
%\textcolor[rgb]{0.98,0.00,0.00}{Note that $S_{\mathfrak{f}|_{(0,0,\textbf{0}_6)+ V}}=q_0+ S_{f}=S_{f}$ (since $q_0=\textbf{0}_6$). Now, the truth table of the function $f$ we construct by concatenating the truth tables of functions $\mathfrak{f}|_{(z,\textbf{0}_6)+V}$ ($z\in \mathbb{F}^2_2$), what we write as}
%$$T_\mathfrak{f}=T_{\mathfrak{f}|_V}||T_{\mathfrak{f}|_{(0,1,\textbf{0}_6)+ V}}||T_{\mathfrak{f}|_{(1,0,\textbf{0}_6)+ V}}||T_{\mathfrak{f}|_{(1,1,\textbf{0}_6)+ V}},$$
%since we took that $V=\textbf{0}_{2}\times \mathbb{F}^6_2$, and we are using the ordering of the space $\mathbb{F}^m_2$ fixed in Section \ref{sec:pre}. The ANF of $\mathfrak{f}$ is given by
%\begin{eqnarray*}
%\mathfrak{f}(x_1,\ldots,x_8)&=&x_1 x_3 x_5 + x_3 x_6 + x_1 x_3 x_6 + x_7 + x_1 x_7 + x_1 x_3 x_8 +
% x_4 (1 + x_5 + x_8 + x_3 (1 + x_7 + x_8)\\
% && + x_1 (x_6 + x_7 + x_8))+ x_2 (1 + x_3 + x_4 + x_3 x_4 + x_3 x_5 + x_4 x_5 + x_3 x_6 + x_4 x_6 + x_8 + x_3 x_8\\
% && + x_4 x_8 + x_1 (1 + x_5 + x_4 (1 + x_6) + x_8 + x_3 (1 + x_4 + x_5 + x_6 + x_8))).
%  \end{eqnarray*}
%\end{ex}
%{\bf For any fixed $(x_1,x_2,x_3,x_4)$ this function $\mathfrak{f}$ is a linear function, thus it is in MM class !! This can be recognized by some skilled reviewer and then the potential of this approach is questionable !! } \\ \\
The above result provides an efficient method of specifying nontrivial disjoint spectra $s$-plateaued functions of maximal cardinality. %{\bf Why is that ?? Nobody is going to understand the arguments why !! We need to be more precise in this context and change Theorem 4.3 to remark or possibly remove it - again too many results is not always good !!!}
%Nevertheless, the method is based on a rather simple partition of the space $\FB^n$ into disjoint subsets.
 In general, the problem of finding {\em extendable sets} of nontrivial disjoint spectra plateaued functions  seems to be quite interesting.
%on disjoint supports without having much structure and then try to extend this set. This leads to the following open problem.
\begin{op}
Assume that we have a set $\{f_1,\ldots,f_r\}$ of nontrivial disjoint spectra $s$-plateaued functions defined on $\mathbb{F}^n_2$, where $r<2^s$. Is it always possible to extend this set with some $s$-plateaued functions $f_{r+1},\ldots,f_{2^s}\in \mathcal{B}_n$ so that $\{f_1,\ldots,f_r,f_{r+1},\ldots,f_{2^s}\}$ is a set of pairwise disjoint spectra $s$-plateaued functions?
\end{op}
%\textbf{Open questions:}
%\begin{enumerate}[1)]
%\item In general, one may consider the construction of plateaued functions $\mathfrak{f}|_{a+ V}$ with affine and non-affine Walsh supports (mixing both). In our construction we have that all plateaued functions $S_{\mathfrak{f}|_{a+ V}}$ constructed by Theorem \ref{th:Sf} have non-affine subspaces.
%%\item In Theorem \ref{th:Sf}, one may consider to take for duals $d^*$ bent functions which are from other classes, PS, C,D, etc...
%\item To which class of bent functions $\mathfrak{f}$ belongs? I do not have any idea how to handle this (I assume that the unproved Theorem 3.2 may help...)
%\item The \textbf{important open problem} in general, which regards "only the construction of plateaued functions", is wether we can provide a generic construction of pairs $f^*$ and $\Phi_f$ so that $f^*$ is at bent distance to $\Phi_f$, and thus $(f^*,S_f)$ rise plateaued functions.
%
%    \textcolor[rgb]{0.98,0.00,0.00}{In other words, can we give a method of finding a vector space of functions which is at bent distance to a fixed bent function $f^*$? (the property of the vector space is that truth tables of its spanning functions must build a non-multiset $S_f$)}. Precisely, the problem is how to find spaces of nonlinear functions on bent distance to $f^*$ (and certainly one can always put linear functions in $S_f$ by adding $\mathbb{F}^{n-s}_2$ as in (\ref{Sff}))? \textbf{If we can provide a good method for this problem, then it will significantly rise the weight of the article, since it would actually provide a generic construction method of plateaued functions with any amplitude and non-affine Walsh support!}
%\end{enumerate}


\section{Applying nonlinear transforms to plateaued functions}\label{sec:nonlperm}


%{\color{red}Fengrong: I feel we must consider  the relationship between the constructed functions and the known ones. For example, whether   the constructed functions can be obtained by restricting bent functions which were constructed in CCDS. }\\\\
In difference to standard approaches, two decades ago Hou and Langevin \cite{HouLang} proposed quite a different framework for specifying the bent properties of Boolean functions. More precisely, considering an arbitrary Boolean functions $f$  one may ask a question what kind of nonlinear permutation $\sigma$ over $\FB^n$ composed to the input variables of $f$ yield a bent function $f \circ \sigma^{-1}$. Quite recently, their approach has been extended in \cite{Houextension2017} where the authors have provided several different methods that result in affine inequivalent bent functions. This method can also be applied to plateaued functions, provided that one can identify some classes of plateaued functions that can be represented in a suitable form. In what follows we recall this method applied to bent functions in detail.

Throughout this section a permutation $\sigma: \FB^n \rightarrow \FB^n$ is represented as a collection of its $n$ coordinate functions so that $\sigma(x)=(\sigma_1(x), \ldots, \sigma_n(x))$, where $\sigma_i:\FB^n \rightarrow \FB$.
For even $n$, denoting by $l(f) =\{g \in \mathcal{B}_n : d_H(f,g)=2^{n-1} \pm 2^{n/2-1}\}$ the set of Boolean functions at the bent distance from $f$ Hou and Langevin showed that  $f \circ \sigma^{-1}$ is bent if and only if the linear span $\langle \sigma_1,\ldots,\sigma_n \rangle =span(\sigma)$ is a subset of $l(f)$, cf. Lemma 3.1 in \cite{HouLang}. Furthermore, they also analyzed a special form of bent functions given as
\begin{equation}
\label{eq:form}
f=x_1f_1 + x_2f_2 +x_1x_2 \alpha + g,
\end{equation}
 where the functions $f_1,f_2,g$ and $\alpha$ (where  $\alpha \in \mathcal{A}_n$) only depend on variables $x_3, \ldots, x_n$.

 The main question is whether similar results can be applied to plateaued functions whose spectra is 3-valued.
\begin{theo}\label{th:sigma}
Let $f \in  \mathcal{B}_n$ be  $s$-plateaued and let $\sigma=(\sigma_1 ,\ldots, \sigma_n): \FB^n \rightarrow    \FB^n$ be a bijection. Then $f \circ \sigma^{-1}$  is an $s$-plateaued   function if and only if  $span(\sigma_1 , \ldots, \sigma_n) \subset lp( f )$, where $lp( f )=\{g \in \mathcal{B}_n: W_{f+g}(u) \in \{0,\pm 2^{\frac{n+s}{2}} \},\;\forall u\in \mathbb{F}^n_2\}$.
\end{theo}
\begin{proof}
%Since $f$ is $s$-plateaued then  $W_f(u) \in \{0,\pm 2^{\frac{n+s}{2}} \}$, for any  $u \in \FB^n$, and clearly $f \circ \sigma^{-1}$ is a Boolean function. It is sufficient to show that $W_{f \circ \sigma^{-1}}(u) \in \{0,\pm 2^{\frac{n+s}{2}} \}$.
%Furthermore, being $s$-plateaued then $W_f(u) \in \{0,\pm 2^{\frac{n+s}{2}} \}$, for any  $u \in \FB^n$.
%Notice that $\#S_f=2^{n-s}$ and we only need to show that $\#S_{f\circ \sigma^{-1}}=2^{n-s}$.
%The distribution of zero and nonzero values in the Walsh spectrum of $f$ remains unchanged for $f \circ \sigma^{-1}$. This  comes form the fact that
For an arbitrary $u\in \mathbb{F}^n_2$ we have that
$$W_{f \circ \sigma^{-1}}(u) = \sum_{x \in \FB^n }(-1)^{f \circ \sigma^{-1}(x)+ u \cdot x} =\sum_{y \in \FB^n }(-1)^{f (y)+ u \cdot \sigma (y)},$$
 and clearly $W_{f \circ \sigma^{-1}}(u) \in \{0,\pm 2^{\frac{n+s}{2}} \}$  if and only if $span(\sigma_1 , \ldots, \sigma_n) \subset lp( f )$.
% if and only if $\alpha \circ \sigma \in lp(f)$  for all $\alpha \in \mathcal{L}_n$.  This is true because $ d_H(f \circ \sigma^{-1}, \alpha)=d_H(f ,\alpha \circ \sigma)$, for all $\alpha \in \mathcal{L}_n$.
\qed
\end{proof}
%{\bf I believe we need to add more clarity to this proof - there is a slight difference between bent and plateaued functions !! However, I believe this result is correct anyway - it can be checked using computer and the design method that we had for bent functions !!} \\ \\
Assuming that $f$ given by (\ref{eq:form}) is a bent function, a nonlinear permutation $\sigma$ on $\FB^n$ considered in \cite{HouLang,Houextension2017} which acts nonlinearly on the variables $x_1$ and $x_2$ was given explicitly as:
\begin{equation} \label{eq:sigma}
\sigma(x_1,x_2,\ldots, x_n)= \big ( (f_1,f_2)+ (x_1,x_2) \left [ \begin{array}{cc} 1 & \alpha + 1 \\ \alpha & 1 \end{array} \right ], x_3, \ldots, x_n \big ).
\end{equation}
Then defining
\begin{equation} \label{eq:siginv}
\tau(x_1,x_2,\ldots, x_n)= \big ( [(f_1,f_2)+ (x_1,x_2)] \left [ \begin{array}{cc} 1 & \alpha + 1 \\ \alpha & 1 \end{array} \right ], x_3, \ldots, x_n \big ),
\end{equation}
one can easily verified that $\sigma \circ \tau= id$ which shows that $\tau=\sigma^{-1}$ and consequently $\sigma$ is a permutation.
Furthermore, assuming that  that $f$ is bent  of the form (\ref{eq:form}) it was shown that $F$ defined  as:
\begin{equation} \label{eq:bentF}
F=(\alpha + 1)f_1f_2 + (x_1+1)f_1 + (x_1+x_2 + \alpha + 1)f_2 + \alpha(x_1+1)x_2 + g,
\end{equation}
is also a bent function. Notice that the bentness of  $f$ is  measured as the distance to linear functions, thus verifying that $d_H(f,l)=2^{n-1} \pm 2^{n/2-1}$ for any $l \in \mathcal{L}_n$. On the other hand, the bentness of $F$ is established by measuring the distance  to linear span of $\sigma$. Thus, to show that $F$ is bent is equivalent to showing  that $d_H(f,l)=2^{n-1} \pm 2^{n/2-1}$, for any $l \in span(\sigma_1,\ldots,\sigma_n)$. Furthermore, we have that $span(\sigma_1,\ldots,\sigma_n) \subset span(\mathcal{A}_n, f_1 + x_2\alpha, f_2 + x_1 \alpha)$.

With the following result, we show that the above approach can be applied to plateaued functions as well.
\begin{theo}\label{theo:Ffsigma}
Let $f \in \mathcal{B}_n$ be an $s$-plateaued function (where $n$ and $s$ are of same parity) be given by (\ref{eq:form}). Then, for the permutation $\sigma$ given by (\ref{eq:sigma}) the function $F=f \circ \sigma^{-1}$ defined by (\ref{eq:bentF}) is also an $s$-plateaued function. Furthermore, $f$ and $F$ are EA-inequivalent.  % which is affine inequivalent to $f$.
\end{theo}
%\textbf{\textcolor[rgb]{0.98,0.00,0.00}{Samir:} Due to Example 3.2, we can not be sure that $F$ and $f$ above are inequivalent...}\\\\
\begin{proof} To show that $F$ is $s$-plateaued, by Theorem \ref{th:sigma} it is sufficient to show that for any $l \in span(\sigma_1,\ldots,\sigma_n)$ we have  $W_{f+l}(u) \in \{0, \pm 2^{\frac{n+s}{2}} \}$, where $u\in \mathbb{F}^n_2$. Since $span(\sigma_1,\ldots,\sigma_n) \subset span(\mathcal{A}_n, f_1 + x_2\alpha, f_2 + x_1 \alpha)$ we assume that $l \in span(\mathcal{A}_n, f_1 + x_2\alpha, f_2 + x_1 \alpha)$ so that $l$ can be written as $l=\epsilon_1(f_1 + x_2\alpha) + \epsilon_2(f_2 + x_1 \alpha) + \gamma$, for $\epsilon_1,\epsilon_2 \in \FB$ and $\gamma \in \mathcal{A}_n$. Then, similarly as in \cite{HouLang}, we can write
$$l= f(x_1, \ldots, x_n) + f(x_1 +\epsilon_1 , x_2+\epsilon_2 , x_3 ,\ldots, x_n)+ \epsilon_1\epsilon_2\alpha + \gamma,$$
which implies that  $f+l= f(x_1 +\epsilon_1 , x_2+\epsilon_2 , x_3 ,\ldots, x_n)+ \epsilon_1\epsilon_2\alpha + \gamma$. However,
since $\alpha \in \mathcal{A}_n$ then $\epsilon_1\epsilon_2\alpha + \gamma \in \mathcal{A}_n$, and therefore the Walsh spectra of $f+l$ is just a (linearly) permuted version of the spectra of $f$. Thus, $F$ is $s$-plateaued as well.  \\
That $f$ and $F$ are EA-inequivalent follows from the fact that $\sigma^{-1}$ is a nonlinear permutation, acting nonlinearly on $x_1$ and $x_2$ through $f_1$ and $f_2$ while keeping the remaining coordinates fixed. Thus, there is no affine permutation such that $F(x)=f(Ax+b)$ (with possible addition of affine function) that can achieve this.
\qed %The fact that $F$ is affine inequivalent to $f$ follows easily form the fact that $\sigma$ is a nonlinear permutation. \qed
\end{proof}
%%%%%%%%%%%%%%%%%%%%%%%%
With the following example we point out that if $f,F:\mathbb{F}^{n}_2\rightarrow \mathbb{F}_2$ satisfy $F(x)=f(H(x))+\ell(x)$, where $\ell\in\mathcal{A}_n$ and $H:\mathbb{F}^{n}_2\rightarrow \mathbb{F}^n_2$ is a non-linear vectorial function (not necessarily a permutation), then it does not mean that $f$ and $F$ are inequivalent.
%This means that functions $F=f \circ \sigma^{-1}$ and $f$ are not necessarily inequivalent plateaued functions in general, though we conjecture that if $f$ can be represented as in (\ref{eq:form}) then $f \circ \sigma^{-1}$ is necessarily EA-inequivalent to $f$.
%\textcolor[rgb]{0.00,0.07,1.00}{unless $f$ can be represented as in (\ref{eq:form}).}\textbf{Samir: This blue part can be critical. The conclusion of the example below is of general nature, and saying  "unless $f$ can be represented as" would mean that the form (\ref{eq:form}) provides inequivalent functions in general, for which we do not have proof...neither the form itself is a proof...}

The following example demonstrates that in general affine transform of the input function can be in general equal to a non-bijective nonlinear action to the input variables.
\begin{ex}\label{ex:inEA2}
Let $f,F:\mathbb{F}^{5}_2\rightarrow \mathbb{F}_2$ be (two trivial plateaued functions) given as $f(x_1,\ldots,x_5)=x_1x_3 + x_2x_4 + x_5$ and $F(x_1,\ldots,x_5)= x_1x_3 + x_2x_4 + x_2x_5 + x_3x_5 + x_4x_5+x_1+x_4.$ In terms of the non-bijective mapping $H:\mathbb{F}^5_2\rightarrow \mathbb{F}^5_2$ defined as (even though strictly speaking $f(H(x))$ is not well-defined)

$$H(x_1,\ldots,x_5)=(x_1,x_2,x_3,x_4, x_2x_5 + x_3x_5 + x_4x_5)$$ the functions $f$ and $H$ are related as
\begin{eqnarray*}
f(H(x))=(x_1x_3 + x_2x_4)+ (x_2x_5 + x_3x_5 + x_4x_5)=F(x_1,\ldots,x_5)+(x_1+x_4).
\end{eqnarray*}
However, for the matrix $A\in GL(5,\mathbb{F}_2)$ given as
$$A=\left(
                                                           \begin{array}{ccccc}
                                                             1 & 0 & 0 & 0 & 0 \\
                                                             0 & 1 & 0 & 0 & 0 \\
                                                             0 & 0 & 1 & 0 & 0 \\
                                                             0 & 0 & 0 & 1 & 0 \\
                                                             1 & 1 & 0 & 1 & 1 \\
                                                           \end{array}
                                                         \right),$$
one can verify that for $H'(x)=(x_1,\ldots,x_5)A=(x_1+x_5,x_2+x_5,x_3,x_4+x_5,x_5)$ it holds that
\begin{eqnarray*}
f(H'(x))=f(xA)&=&(x_1+x_5)x_3+(x_2+x_5)(x_4+x_5)+x_5\\
&=&(x_1x_3+x_3x_5)+(x_2x_4+x_2x_5+x_4x_5+x_5)+x_5\\
&=&x_1 x_3 + x_2 x_4 + x_2 x_5 + x_3 x_5 + x_4 x_5\\
&=&F(x_1,x_2,x_3,x_4,x_5)+(x_1+x_4),
\end{eqnarray*}
which means that $f$ and $F$ are affine equivalent.
\end{ex}
%\begin{rem}
%In general, Theorem \ref{theo:Ffsigma} gives EA-inequivalent plateaued functions $f$ and $F=f\circ\sigma^{-1}$ as also demonstrated in \cite{Houextension2017}, see also Example \ref{ex:F}.
%\end{rem}
%{\bf Question 6: Can we provide a similar proof that $F$ is plateaued in this case ?? I will take a look at this problem later. Anyway, if this is the case or for some other form $f$ and $F$ then we essentially provide AFFINE INEQUIVALENT PLATEAUED FUNCTIONS easily !!!} \\ \\
To illustrate the application of Theorem \ref{theo:Ffsigma}  we consider the semi-bent function $f$ (thus 1-plateaued) in Example 4.1.
%\textcolor[rgb]{1.00,0.00,0.00}{We also show that Theorem \ref{theo:Ffsigma} can also provide inequivalent plateaued functions.}
\begin{ex}\label{ex:F}
Let $f \in \mathcal{B}_5$ be a semi-bent function   given by $f(x_1,\ldots,x_5)=x_1x_2x_5+ x_1x_3+ x_2x_4+ x_5$ which can be written as $x_1 f_1 + x_2f_2 + x_1x_2 \alpha + g$, where $f_1=x_3$, $f_2=x_4$, $\alpha=x_5$ and $g=x_5$. Apparently, the functions $f_1,f_2, \alpha$ and  $g$ do not depend on the variables $x_1,x_2$ and furthermore $\alpha$ is affine which is the main condition for the construction to work. Then,
\begin{eqnarray*}
F(x)&=& (\alpha + 1)f_1f_2 + (x_1+1)f_1 + (x_1+x_2 + \alpha + 1)f_2 + \alpha(x_1+1)x_2 + g \\
&=& (x_5+1)(x_3x_4) + (x_1+1)x_3 + (x_1+x_2+ x_5+1)x_4 + x_5(x_1+1)x_2 + x_5 \\
&=& x_1x_2x_5 + x_3x_4x_5 + x_1x_3 + x_1x_4 + x_2x_4 + x_2x_5 +x_3x_4 + x_4x_5 + x_3 + x_4 + x_5.
\end{eqnarray*}
%\textcolor[rgb]{0.00,0.00,1.00}{Even though $F$ and $f$  are of the same algebraic degree these semi-bent functions are not affine equivalent.} \textbf{\textcolor[rgb]{0.98,0.00,0.00}{Samir:} We need to check this as well. My feeling is that Theorem 3.2 can be extended to non-trivial plateaued function by including the existence of "some suitable orderings of $S_h$ and $S_f$". I need to consider this possibility...}
The function $F$ was checked to be a plateaued (semi-bent) function  with:
$$W_F=\{0, 8, 0, 8, 0, 8, 8, 0, 0, -8, 0, 8, 0, 8, -8, 0, 0, -8, 0, -8, 0, 8, 8, 0, -8, 0, 8, 0, -8, 0, 0, 8\}.$$
The fact that $f$ and $F$ are EA-inequivalent follows from Theorem \ref{theo:Ffsigma}, but can  also be proved using the results of Section \ref{sec:equivalence} and showing the inequivalence of their supports.
%we firstly consider their Walsh supports. Recall that $S_f=\mathbb{F}^4_2\wr T_{\lambda}$, where $\lambda(x_1,\ldots,x_4)=x_3x_4$, and one can verify that $S_F=\mathbb{F}^{4}_2\wr T_{\mu}$, where $\mu(x)=x_1 x_2 + x_3 x_4+1$. In addition, we have that $F^*(x)=x_1 + x_2 + x_1 x_2 + x_1 x_3 + x_2 x_3 + x_2 x_4$,  $x\in \mathbb{F}^4_2.$
%
%By Theorem \ref{th:eq}-$(i)$, there must exist a matrix $M\in GL(5,\mathbb{F}_2)$ and a vector $c\in \mathbb{F}^5_2$ such that $S_f=c+S_FM$. Denoting by $M=[M_1,M_2,\ldots,M_5]$, where $M_i$ are the columns of  $M$, and setting $S_F=\{\omega_0,\ldots,\omega_{2^{4}-1}\}$ (here $n-s=5-1=4$) we have that
% $$S_f=c+S_FM=c+\left(
%                  \begin{array}{c}
%                    \omega_0 \\
%                    \omega_1 \\
%                    \vdots \\
%                    \omega_{2^4-1} \\
%                  \end{array}
%                \right)[M_1,\ldots,M_5]
% =c+\left(
%                  \begin{array}{ccc}
%                    \omega_0\cdot M_1 & \ldots &  \omega_0\cdot M_5 \\
%                    \omega_1\cdot M_1 & \ldots &  \omega_1\cdot M_5 \\
%                    \vdots & \ldots  & \vdots \\
%                    \omega_{2^4-1}\cdot M_1 & \ldots &  \omega_{2^4-1}\cdot M_5  \\
%                  \end{array}
%                \right),
% $$
% where in $\omega_j\cdot M_i$ the column $M_i$ is considered as a row vector. As we can see, the columns of $S_f$ are linear combinations of columns of $S_F$, i.e., for $c=(c_1,\ldots,c_5)\in \mathbb{F}^5_2$ an arbitrary column $T_i$ of $S_f=\mathbb{F}^4_2\wr T_{\lambda}=T_1\wr \ldots \wr T_5$ (viewed as a vector) can be written as
% \begin{eqnarray}\label{eq:coll}
% T_i=(\omega_0\cdot M_i+c_i,\ldots,\omega_{2^4-1}\cdot M_i+c_i),\;\;\;\;\;i\in [1,5].
% \end{eqnarray}
%On the other hand, since  $S_F=\mathbb{F}^{4}_2\wr T_{\mu}$ and $\mu$ is a bent function on $\mathbb{F}^4_2$, then clearly $wt(\mu+\ell)\in\{2^3-2,2^3+2\}=\{6,10\}$ holds for every $\ell\in \mathcal{A}_4.$ Also, since $\lambda$ is a semi-bent function on $\mathbb{F}^4_2$ (that is $W_{\lambda}(u)\in\{0,\pm 2^3\}$, $u\in \mathbb{F}^4_2$), then in general $wt(\lambda+\delta)=2^3-\frac{1}{2}W_{\lambda+\delta}(\textbf{0}_4)\in\{2^3,2^3\pm 4\}=\{4,8,12\}$ holds for all $\delta\in \mathcal{A}_4.$
%
%Consequently, the previous facts on weights imply that the column $T_{\lambda}$ of $S_f$ can not be written in the form (\ref{eq:coll}) for any matrix $M$ and vector $c$, and this conclusion holds for every possible ordering that one can consider on $S_f$ and $S_F$ in general. This means that the Walsh supports $S_f$ and $S_F$ can not be related in an affine manner as $S_f=c+S_FM$, which by Theorem \ref{th:eq}-$(i)$ means that $f$ and $F$ are inequivalent.
\end{ex}
%\begin{rem}
%In terms of extendability of such a plateaued function the question is whether there exist a (unique) plateaued function $F'$ having a  disjoint spectra to $F$. Using computer based search, it was confirmed that there are actually 384  different plateaued functions $F'$ whose spectra is disjoint to that of $F'$. This number is a simple consequence of the fact that there are 28 different bent functions on $\FB^4$ up to addition of an affine function. This then implies that using different bent duals up to addition of a linear function will yield $28 \times 16=384$ different $F'$. {\bf Why this number is not 768 instead ?? Taking only linear functions is not justified. There are 768 bent functions on $\FB^4$ and we can take any such function to define the dual of $F'$ on the disjoint support !! I guess something is wrong with your programs ??}
%\end{rem}
%%%%%%%%%%%%%%%%%%%%%%%%%%%%%%%%%%%%%%%%%%%%%
%Notice that the action of a suitable  nonlinear permutation $\sigma$ roughly corresponds to a nonlinear action on the Walsh spectra of a given plateaued function $f$. Indeed, any linear permutation $\sigma$ would correspond to a standard linear transformation of the inputs  so that $f(x)$ is transformed into its linear equivalent $f(Ax)$, where $A$ is an invertible $n \times n$ binary matrix.
%\begin{rem} %One may verify that the Walsh support of the function $F$ (from Example \ref{ex:F}) is $S_F=\mathbb{F}^{4}_2\wr T_{\mu}$, where $\mu(x)=x_1 x_2 + x_3 x_4+1$, and its dual is  given by
% $S_F=\mathbb{F}^{4}_2\wr T_{\mu}$, where $\mu(x)=x_1 x_2 + x_3 x_4+1$, and its dual
%$F^*(x)=x_1 + x_2 + x_1 x_2 + x_1 x_3 + x_2 x_3 + x_2 x_4$,  $x\in \mathbb{F}^4_2.$
%Since both $\mu$ and $F^*$ are bent functions on $\mathbb{F}^*_2$, this means that the function $F$ is a special case of Corollary \ref{cor:vectorial} (where $c=\textbf{0}_4$, $M=I_{4\times 4},$ $s=1$ and $r=1$ with $\lambda=2$, $h_1=\mu$, $h_2=g=F^*$).
%Note that there exist many disjoint spectra functions to $F$ (by computer search we find $384$ such functions), and their Walsh supports are given exactly as $\mathbb{F}^4_2\wr T_{\mu+1}$, where the duals are bent functions on $\FB^4$ at bent distance to $\mu$.
%\end{rem}
\begin{rem} One may verify that the Walsh support of the function $F$ (from Example \ref{ex:F}) is $S_F=\mathbb{F}^{4}_2\wr T_{\mu}$, where $\mu(x)=x_1 x_2 + x_3 x_4+1$, and its dual is  given by
% $S_F=\mathbb{F}^{4}_2\wr T_{\mu}$, where $\mu(x)=x_1 x_2 + x_3 x_4+1$, and its dual
$F^*(x)=x_1 + x_2 + x_1 x_2 + x_1 x_3 + x_2 x_3 + x_2 x_4$,  $x\in \mathbb{F}^4_2.$
%Since both $\mu$ and $F^*$ are bent functions on $\mathbb{F}^*_2$, this means that the function $F$ is a special case of Corollary \ref{cor:vectorial} (where $c=\textbf{0}_4$, $M=I_{4\times 4},$ $s=1$ and $r=1$ with $\lambda=2$, $h_1=\mu$, $h_2=g=F^*$).
 It is not difficult to check that there exist many disjoint spectra functions to $F$ (by computer search we find $384$ such functions), and their Walsh supports are given exactly as $\mathbb{F}^4_2\wr T_{\mu+1}$, where the duals are bent functions on $\FB^4$ at bent distance to $\mu$.
\end{rem}
%%%
The question whether this nonlinear action can be generalized to provide a set of disjoint
spectra plateaued functions of maximal cardinality seems to be difficult and is left as an open
problem.



\section{Conclusions}\label{sec:OP}

The purpose of this article is to provide further clarity concerning the essential structural
properties of a class of Boolean functions known as plateaued functions. We have provided
an efficient approach of designing these functions in the spectral domain, both trivial and nontrivial ones. The EA-equivalence within and between these classes have been addressed as well as the design of disjoint spectra (regardless of being trivial or not) plateaued functions of maximal cardinality.
 %\textcolor[rgb]{0.98,0.00,0.00}{considering the notion of equivalency and non-triviality of Walsh supports.}
The most challenging problem
that remains open is whether the constructed disjoint spectra plateaued functions of maximal
cardinality (whose concatenation is a bent function) may give rise to some new classes of bent
functions.
%The main benefit of this method is that in terms of EA-equivalence it ensures that the modified function is inequivalent to the initial one which is not easy to achieve using the general equivalence result based on duals.
\\\\
\noindent
{\large \bf Acknowledgment:} Yongzhuang Wei (corresponding author) is supported in part by the
National Key R\&D Program of China (No. 2017YFB0802000)
and in part by the Natural Science Foundation of China (Nos.
61572148, 61872103).  Samir Hod\v zi\' c  is supported in part by the Slovenian Research Agency (research program P3-0384 and Young Researchers Grant). Enes Pasalic is partly supported  by the Slovenian Research Agency (research program P3-0384 and research project J1-9108). For the first two authors, the work is supported in part by H2020 Teaming InnoRenew CoE (grant no. 739574). Fengrong Zhang is supported by Jiangsu Natural Science Foundation (BK20181352).


%{\bf There are two new articles on bent functions in IEEE 2017 - need to check these and add references}
%\newpage
\begin{thebibliography}{10}

\bibitem{Decom}
{\sc A. Canteaut, P. Charpin.}
\newblock Decomposing bent functions.
\newblock {\em IEEE Transactions on Information Theory}, vol. 49, no. 8, 2003.

%\bibitem{AnneCosets}
%{\sc A. Canteaut, C. Carlet, P. Charpin, C. Fontaine.}
%\newblock On cryptographic properties of the cosets of $R(1,m)$
%\newblock {\em IEEE Transactions on Information Theory}, vol. 47, no. 4, 2001.

%\bibitem{Adams}
%{\sc C.M. Adams and S.E. Tavares.}
%\newblock Generating and counting binary bent sequences.
%\newblock {\em  IEEE Transactions on Information Theory}, vol. 36, no. 5, pp. 1170-1173, 1990.

%\bibitem{CarletSC}
%{\sc C. Carlet.}
%\newblock A construction of bent functions.
%\newblock {\em Finite Fields and Applications,  London Mathematical Society. Lecture Series 233, Cambridge University Press},  pp. 47--58, 1996.

\bibitem{CarletAPN}
{\sc C. Carlet.}
\newblock Boolean and vectorial plateaued functions and APN functions.
\newblock {\em IEEE Transactions on Information Theory}, vol. 61, no. 11, pp. 6272--6289, 2015.

%\bibitem{CarletBoolean}
%{\sc C. Carlet.}
%\newblock Boolean models and methods in mathematics, computer science, and engineering.
%\newblock {\em  Encyclopedia of Mathematics and its Applications (No. 134) - Cambridge University Press}, pp. 398 -- 469, 2013. %Available at: http://www.math.univ-paris13.fr/~carlet/chap-vectorial-fcts-corr.pdf
%
%\bibitem{CarletGPS}
%{\sc C. Carlet.}
%\newblock Generalized Partial Spreads.
%\newblock {\em  IEEE Transactions on Information Theory}, vol. 41, no. 5, pp. 1482--1487, 1995.


%\bibitem{CarletRes}
%{\sc C. Carlet.}
%\newblock On bent and highly nonlinear balanced/resilient functions and their algebraic immunities.
%\newblock {\em 16th International Symposium, AAECC-16 - Applied Algebra, Algebraic Algorithms and Error-Correcting Code}, vol. 3857, pp. 1--28 , 2006.
%
%\bibitem{CarletRESB}
%{\sc C. Carlet.}
%\newblock  On the secondary constructions of resilient and bent functions.
%\newblock {\em  Proceedings of the Workshop on Coding, Cryptography and Combinatorics 2003, published by Birkh�user Verlag}, pp. 3--28, 2004.

%\bibitem{CarletOP}
%{\sc C. Carlet.}
%\newblock Open problems on binary Bent functions.
%\newblock {\em  Open Problems in Mathematics and Computational Science - Springer International Publishing}, pp. 203-- 242, 2014.

\bibitem{CarletPartial}
{\sc C. Carlet.}
\newblock Partially bent functions.
\newblock {\em  Designs, Codes and Cryptography}, vol. 3, no. 2, pp. 135--145, 1993.

\bibitem{CarletTNC}
{\sc C. Carlet.}
\newblock Two new classes of bent functions.
\newblock {\em  Advances in Cryptology, EUROCRYPT 93, Lecture Notes in Computer Science}, LNCS, vol. 765, pp. 77--101, 1993.


%\bibitem{Decom}
%{\sc A. Canteaut, P. Charpin.}
%\newblock Decomposing bent functions.
%\newblock {\em IEEE Transactions on Information Theory}, vol. 49, no. 8, 2003.

%\bibitem{CarletN}
%{\sc C. Carlet, H. Dobbertin, G. Leander.}
%\newblock Normal extensions of bent functions.
%\newblock {\em  IEEE Transactions on Information Theory}, vol. 50, no. 11, 2004.
%
%\bibitem{CarletRSF}
%{\sc C. Carlet, G. Gao, W. Liu.}
%\newblock A secondary construction and a transformation on rotation symmetric functions, and their action on bent and semi-bent functions.
%\newblock {\em  Journal of Combinatorial Theory, Series A}, vol. 127, pp. 161--175, 2014.
%
%\bibitem{CarletMess}
%{\sc C. Carlet, S. Mesnager.}
%\newblock On Dillon's class H of bent functions, Niho bent functions and o-polynomials.
%\newblock {\em  Journal of Combinatorial Theory, Series A}, vol. 118, no. 8, pp. 2392--2410, 2011.
%
%\bibitem{CarletSupp}
%{\sc C. Carlet, S. Mesnager.}
%\newblock On the supports of the Walsh transforms of Boolean functions.
%\newblock {\em  Boolean Functions: Cryptography and Applications, BFCA'05}, pp. 65 -- 82, 2005. Available at:  https://eprint.iacr.org/2004/256.pdf

\bibitem{CarletQ}
{\sc C. Carlet, E. Prouff.}
\newblock On plateaued functions and their constructions.
\newblock {\em  International Workshop on Fast Software Encryption, FSE 2003}, pp. 54--73, 2003.



%\bibitem{CarletPiece}
%{\sc C. Carlet, J. L. Yucas.}
%\newblock Piecewise constructions of bent and almost optimal Boolean functions.
%\newblock {\em  Designs, Codes and Cryptography}, vol. 37, no. 3, pp. 449--464, 2005.

\bibitem{SecEnf}
{\sc C. Carlet, F. Zhang, Y. Hu.}
\newblock Secondary constructions of bent functions and their enforcement.
\newblock {\em Advances in Mathematics of Communications}, vol. 6, no. 3, pp. 305 -- 314, 2012.

\bibitem{AycaWilffried}
{\sc A. \c{C}e\c{s}melio\u{g}lu, W. Meidl.}
\newblock A construction of bent functions from plateaued functions.
\newblock {\em Designs, Codes and Cryptography}, vol. 66, no. 1-2, pp. 231--242, 2013.

\bibitem{Ayca2}
{\sc A. \c{C}e\c{s}melio\u{g}lu, G. McGuire, W. Meidl.}
\newblock A construction of weakly and non-weakly regular bent functions.
\newblock {\em Journal of Combinatorial Theory, Series A}, vol. 119, no. 2, pp. 420--429, 2012.

\bibitem{THOMCusick2016}
{\sc T. W. Cusick.}
\newblock Highly nonlinear plateaued functions.
\newblock {\em IET Information Security}, vol. 11 PP, 78-81 no. 2 , 2016.

%\bibitem{Adams}
%{\sc C.M. Adams and S.E. Tavares.}
%\newblock Generating and counting binary bent sequences.
%\newblock {\em  IEEE Transactions on Information Theory}, vol. 36, no. 5, pp. 1170-1173, 1990.

%\bibitem{Buda}
%{\sc Budaghyan L., Carlet C., Pott A.}
%\newblock New classes of almost bent and almost perfect nonlinear functions.
%\newblock {\em  IEEE Transactions on Information Theory}, vol. 52, no. 3, pp. 1141?152, 2006.

%\bibitem{CarletWM}
%{\sc C. Carlet.}
%\newblock On the properties of vectorial functions with plateaued components and their consequences on APN Functions.
%\newblock {\em  Codes, Cryptology, and Information Security -  Lecture Notes in Computer Science Volume 9084},  pp. 63--73, 2015.


%\bibitem{CarletCCZ}
%{\sc C. Carlet, P. Charpin, V. Zinoviev.}
%\newblock Codes, bent functions and permutations suitable for DES-like cryptosystems.
%\newblock {\em Des. Codes Cryptogr. }, vol. 15, no. 2, pp. 125--156, 1998.



%\bibitem{CarletRESB2}
%{\sc C. Carlet, T. Helleseth, A. Kholosha, S. Mesnager.}
%\newblock  On the duals of bent functions with $2^r$ Niho exponents.
%\newblock {\em IEEE International Symposium on  Information Theory}, pp. 703-?07, 2011.

\bibitem{Dillon}
{\sc J. F. Dillon.}
\newblock Elementary Hadamard difference sets, Ph.D. dissertation.
\newblock {\em  University of Maryland, College Park, Md, USA}, 1974.

%\bibitem{Do95}
%{\sc H. Dobbertin.}
%\newblock  Construction of bent functions and balanced Boolean functions with high nonlinearity.
%\newblock {\em Fast Software Encryption '94, LNCS 1008, Springer-Verlag}, pp. 61--74, 1995.

%\bibitem{Dobbertin}
%{\sc H. Dobbertin, G. Leander, A. Canteaut, C. Carlet, P. Felke, P. Gaborit.}
%\newblock  Construction of bent functions via Niho power functions.
%\newblock {\em Journal of Combinatorial Theory, Series A}, vol. 113, no. 5, pp. 779-798, 2006.
%
%
%\bibitem{SHWM}
%{\sc S.  Hod\v zi\'c, W. Meidl, E. Pasalic.}
%\newblock Full characterization of generalized bent functions as (semi)-bent spaces, their dual, and the Gray image.
%\newblock Available at: https://arxiv.org/abs/1605.05713

%\bibitem{SHSB}
%{\sc S.  Hod\v zi\'c, E. Pasalic.}
%\newblock Construction methods for generalized bent functions.
%\newblock Available at: https://arxiv.org/abs/1604.02730

%\bibitem{SHCC}
%{\sc S.  Hod\v zi\'c, E. Pasalic.}
%\newblock Generalized bent functions - sufficient conditions and related constructions.
%\newblock {\em Advances in Mathematics of Communications}, vol. 11, no. 3, pp. 549 -- 566, 2017.


%\bibitem{Xiang}
%{\sc X.-D. Hou, P. Langevin.}
%\newblock Results on bent functions.
%\newblock {\em Journal of Combinatorial Theory, Series A}, vol. 80, no. 2, pp. 232--246, 1997.

%\bibitem{Xiang2}
%{\sc X.-D. Hou.}
%\newblock Classification of self dual quadratic bent functions.
%\newblock {\em Designs, Codes and Cryptography}, vol. 63, no. 2, pp. 183--198, 2012.





%\bibitem{DillonSurvey}
%{\sc J. F. Dillon.}
%\newblock A survey of bent functions.
%\newblock {\em NSA Technical Journal Special Issue}, pp. 191-215, 1972.

%\bibitem{Hang}
%{\sc K. H. Kim, K.F. W. Roush.}
%\newblock  Inverses of Boolean matrices.
%\newblock {\em Linear Algebra and its Applications}, vol. 22, pp. 247--262, 1978.

%\bibitem{Distance}
%{\sc  N. Kolomeec, A. Pavlov.}
%\newblock  Bent functions on the minimal distance.
%\newblock {\em IEEE Region 8 SIBIRCON, Irkutsk Listvyanka, Russia, July 11 -- 15}, 2010.


%\bibitem{Svetla}
%{\sc O. A. Logachev, A. A. Salnikov, V. V. Yashchenko.}
%\newblock Boolean functions in coding theory and cryptography.
%\newblock {\em Translations of Mathematical Monographs, American Mathematical Society Boston, MA, USA}, 2012.

%\bibitem{Bimal}
%{\sc B. Mandal, P. Stanica, S. Gangopadhyay, E. Pasalic.}
%\newblock An analysis of the $\mathcal{C}$ class of bent functions.
%\newblock {\em  Fundamenta Informaticae}, vol. 146, no. 3, pp. 271--292, 2016. Available at: https://eprint.iacr.org/2015/588.pdf

\bibitem{Secondary}
{\sc S. Hod\v zi\'c, E. Pasalic, Y. Wei.}
\newblock A general framework for  secondary constructions  of bent and plateaued functions.
\newblock {\em IEEE Trans.  Inf. Theory}, under revision, 2017.

\bibitem{HouLang}
{\sc   X-D. Hou, P. Langevin.}
\newblock Results on bent functions.
\newblock {\em Journal of Combinatorial Theory, Series A,}  vol. 80, pp. 232--246,  1997.

\bibitem{Jong}
{\sc J. Y. Hyun, J. Lee, Y. Lee.}
\newblock Explicit criteria for construction of plateaued functions.
\newblock {\em IEEE Transactions on Information Theory}, vol. 62, no. 12, pp. 7555--7565, 2016.



\bibitem{Logatchev}
{\sc A. O. Logatchev.}
\newblock On a recursive class of plateaued Boolean functions.
\newblock {\em Discrete Mathematics and Applications}, vol. 22, No. 4, pp. 20--33, 2010.

%\bibitem{SihemBook}
%{\sc S. Mesnager.}
%\newblock Bent Functions - Fundamentals and Results.
%\newblock {\em Theoretical Computer Science - Springer International Publishing}, 2016.
%
%\bibitem{Sihem2}
%{\sc S. Mesnager.}
%\newblock Further constructions of infinite families of bent functions from new permutations and their duals.
%\newblock {\em Cryptography and Communications}, vol. 8, no. 2, pp. 229--246, 2016.

\bibitem{MMclass}
{\sc R. L. McFarland.}
\newblock A family of noncyclic difference sets.
\newblock {\em Journal of Combinatorial Theory, Series A,} vol. 15, no. 1, pp. 1--10, 1973.

\bibitem{SihemPOP2014}
{\sc S. Mesnager.}
\newblock On semi-bent functions and related plateaued functions over the Galois field $F_{2^n}$.
\newblock In Proceedings,  {\em Open Problems in Mathematics and Computational Science}, LNCS, Springer, pp. 243--273, 2014.

%\bibitem{Sihem}
%{\sc S. Mesnager.}
%\newblock Several new infinite families of bent functions and their duals.
%\newblock {\em IEEE Transactions on Information Theory}, vol. 60, no. 7, pp. 4397--4407, 2014.

\bibitem{SihemM2017IT}
{\sc S. Mesnager, C. Tang, Y.  Qi.}
\newblock Generalized plateaued functions and admissible (plateaued) functions.
\newblock{\em IEEE Transactions on Information Theory}, vol. 63, no. 10, pp. 6139--6148, 2017.


\bibitem{SihemPlate}
{\sc S. Mesnager, F. {\"O}zbudak, A. Sinak.}
\newblock Results on characterizations of plateaued functions in arbitrary characteristic.
\newblock {\em Cryptography and Information Security in the Balkans}, LNCS vol. 9450, pp. 17--30, 2016.

%\bibitem{SihemTransl}
%S.~Mesnager, P.~Ongan, F.~\"{O}zbudak,
%\newblock New bent functions from permutations and linear translators,
%\newblock C2SI 2017: Codes, Cryptology and Information Security, pp. 282-297, 2017.

\bibitem{SihemM2017AMC}
{\sc S. Mesnager, and F. Zhang.}
\newblock On constructions of bent, semi-bent and five valued spectrum functions from old bent functions.
\newblock {\em Advances in Mathematics of Communications} vol. 11, no. 2, pp. 339--345, 2017.

%\bibitem{Nyberg}
%{\sc K. Nyberg.}
%\newblock Perfect nonlinear S-boxes.
%\newblock {\em  EUROCRYPT'91 Proceedings of the 10th annual international conference on Theory and application of cryptographic techniques}, pp. 378--386, 1992.


%\bibitem{kus}
%{\sc K.U. Schmidt}.
%\newblock Quaternary constant-amplitude codes for multicode CDMA.
%\newblock {\em IEEE Transactions on Information Theory}, vol. 55, no 4, pp. 1824--1832, 2009.

%\bibitem{Rao}
%{\sc P.S.S.N.V. Prasada Rao, K.P.S. Bhaskara Rao.}
%\newblock On generalized inverses of Boolean matrices.
%\newblock {\em Linear Algebra and its Applications}, vol. 11, no. 2, pp. 135--153, 1975.

\bibitem{Houextension2017}
{\sc E. Pasalic, S. Hod\v zi\'c, F. Zhang, Y. Wei}
\newblock Bent functions from nonlinear permutations and conversely.
\newblock {\em Cryptography and communications}, pp. 1--19, 2018.

%\bibitem{Rot}
%{\sc O. S. Rothaus.}
%\newblock On "bent" functions.
%\newblock {\em Journal of Combinatorial Theory, Series A}, vol. 20, no. 3, pp. 300--305, 1976.

\bibitem{CCA}
{\sc P. Sarkar, S. Maitra.}
\newblock Cross-correlation analysis of cryptographically useful Boolean functions and S-boxes.
\newblock {\em Theory of Computing Systems}, vol. 35, no. 1, pp. 39--57, 2002.

%\bibitem{SeberryX}
%{\sc J. Seberry, X-M. Zhang.}
%\newblock Highly nonlinear $0-1$ balanced Boolean functions satisfying strict avalanche criterion.
%\newblock {\em Advances in Cryptography - Auscrypt'92., Berlin-Springer}, LNCS vol. 718, pp. 145--755, 1993.

%\bibitem{SeberryX2}
%{\sc J. Seberry, X-M. Zhang.}
%\newblock Constructions of bent functions from two known bent.
%\newblock {\em Australasian Journal of Combinatorics}, vol. 9, pp. 21--35, 1994.


%\bibitem{BentGbent}
%\newblock P. Stanica, T. Martinsen, S. Gangopadhyay, B. K. Singh.
%\newblock Bent and generalized bent Boolean functions.
%\newblock \emph{Designs, Codes and Cryptography}, vol. 69, no. 1, pp. 77--94, 2013.

%\bibitem{tang}
%{\sc C. Tang, C. Xiang, Y. Qi, K. Feng.}
%\newblock Complete characterization of generalized bent and $2^k$-bent Boolean functions.
%\newblock {\em IEEE Transactions on Information Theory}, vol. PP, no. 99, 2017.

%\bibitem{SHSC}
%{\sc S.  Hod\v zi\'c, E. Pasalic.}
%\newblock Generalized bent functions - sufficient conditions and related constructions.
%\newblock To appear. Available at: http://arxiv.org/abs/1601.08084



%\bibitem{mesna}
%{\sc S. Mesnager.}
%\newblock Several infinite classes of bent functions and their duals.
%\newblock {\em IEEE Transactions on Information Theory}, vol. 60, no. 7, pp. 4397--4407, 2014.

%\bibitem{Wang}
%{\sc W. Weiqiong, X. Guozhen.}
%\newblock Decomposition and construction of plateaued functions.
%\newblock {\em Chinese Journal of Electronics}, vol. 18, no. 4, 2009.


%\bibitem{Hyper}
%{\sc A.M. Youssef, G. Gong.}
%\newblock Hyper-bent functions.
%\newblock {\em Proceedings of EUROCRYPT 2001. Lecture Notes in Computer Science, vol. 2045, Springer, Berlin}, pp. 406?19, 2001.



\bibitem{Feng2}
{\sc F. Zhang, C. Carlet, Y. Hu, T. J. Cao.}
\newblock Secondary constructions of highly nonlinear Boolean functions and disjoint spectra plateaued functions.
\newblock {\em  Information Sciences}, vol. 283, pp. 94--106, 2014.

%\bibitem{Feng}
%{\sc F. Zhang, C. Carlet, Y. Hu, W. Zhang.}
%\newblock New secondary constructions of bent functions.
%\newblock {\em  Applicable Algebra in Engineering, Communication and Computing}, vol. 27, no. 5, pp. 413--434, 2016.

%\bibitem{FengMM}
%{\sc F. Zhang, E. Pasalic, N. Cepak,  Y. Wei.}
%\newblock Bent functions in $\mathcal{C}$ and $\mathcal{D}$ outside the completed Maiorana-McFarland class.
%\newblock {\em  International Conference on Codes, Cryptology, and Information Security,  Lecture Notes in Computer Science}, LNCS, vol. 10194, pp 298--313, 2017.
%
%\bibitem{RothEnes2016}
%{\sc F. Zhang, E. Pasalic, Y. Wei, N. Cepak.}
%\newblock Constructing bent functions outside the Maiorana-McFarland class using a general form of Rothaus.
%\newblock {\em  IEEE Transactions on Information Theory}, vol. 63, no. 8, pp. 5336 -- 5349, 2017.

%\bibitem{Feng4}
%{\sc F. Zhang, Y. Wei, E. Pasalic.}
%\newblock Constructions of bent - negabent functions and their relation to the completed Maiorana - McFarland class.
%\newblock {\em  IEEE Transactions on Information Theory}, vol. 61,  no. 3, pp. 1496--1506, 2015.

\bibitem{Fengrong2018IT}
{\sc F. Zhang, Y. Wei, E. Pasalic, S. Xia.}
\newblock Large sets of disjoint spectra plateaued functions inequivalent to partially linear functions.
\newblock {\em  IEEE Transactions on Information Theory},  vol. 64, no. 4, pp. 2987--2999, 2018.

\bibitem{Zhang2009}
{\sc W. Zhang, G. Xiao.}
\newblock Constructions of almost optimal resilient boolean functions on large even number of variables.
\newblock \emph{IEEE Transactions on Information Theory}, vol. 55, no. 12,  pp. 5822--5831, 2009.

\bibitem{WeiGuoCDMA}
{\sc W. Zhang, C.  Xie, E. Pasalic.}
\newblock Large sets of orthogonal sequences suitable for applications in CDMA systems.
\newblock {\em IEEE Transactions on Information Theory}, vol. 62, no. 6, pp. 3757--3767, 2016.

\bibitem{Zheng}
{\sc Y. Zheng, X. M. Zhang.}
\newblock On plateaued functions.
\newblock {\em IEEE Transactions on Information Theory}, vol. 47, no. 3, pp. 1215--1223, 2001.

\bibitem{Zheng2}
{\sc Y. Zheng, X. M. Zhang.}
\newblock Relationships between bent functions and complementary plateaued functions.
\newblock {\em Information Security and Cryptology - ICISC99}, LNCS vol. 1787, pp. 60--75, 2000.



\end{thebibliography}

\section{Appendix}
In the context of Remark \ref{rem:ordering}, we provide an alternative  proof of Theorem \ref{th:EA1}. For this purpose, we also need  to modify Lemma \ref{prop:Hrow1} slightly.


%Now we firstly introduce some necessary notation.
Let $(i)_2$ ($i\in[0,2^{m}-1]$) denote the binary representation of length $m$ of the integer $i$. The linear subspaces $E,E'\subset \mathbb{F}^n_2$ ($\dim(E)=\dim(E')=n-s$) and ${\Bbb F}_2^{n-s}$ we order  as
\begin{eqnarray}\label{eq:EE'}\nonumber
E&=&\{0_n,(1)_2\cdot \mathbf{B}_E,\ldots, (2^{n-s}-1)_2\cdot \mathbf{B}_E\}=\{(i)_2\cdot \mathbf{B}_E: i\in[0,2^{n-s}-1]\},\\
{\Bbb F}_2^{n-s}&=&\{(i)_2\cdot (\beta^{(1)},\beta^{(2)},\ldots,\beta^{(n-s)}): i\in[0,2^{n-s}-1]\},\\\nonumber
E'&=&\{0_n,(1)_2\cdot \mathbf{B}_{E'},\ldots, (2^{n-s}-1)_2\cdot \mathbf{B}_{E'}\}=\{(i)_2\cdot \mathbf{B}_{E'}: i\in[0,2^{n-s}-1]\},
\end{eqnarray}
where the corresponding bases are given as $\mathbf{B}_E=(\alpha^{(1)},\alpha^{(2)},\ldots,\alpha^{(n-s)})$,
$\mathbf{B}_{{\Bbb F}_2^{n-s}}=(\beta^{(1)},\beta^{(2)},$ $\ldots,\beta^{(n-s)})$ and
  $\mathbf{B}_{E'}=(\alpha'^{(1)},\alpha'^{(2)},\ldots,\alpha'^{(n-s)})$.
To clarify the notation, let $n-s=3$. Then the previous notation means that for $(3)_2=(0,1,1)$ and basis $\mathbf{B}_E=(\alpha^{(1)},\alpha^{(2)},\alpha^{(3)})$ we have that $(3)_2\cdot \mathbf{B}_E=\alpha^{(2)}+\alpha^{(3)}.$ Notice that all sets in (\ref{eq:EE'}) satisfy the recursion mentioned in Remark \ref{rem:ordering}.
%%%%%%%%%%%%%%%%%%%%%%%%%%%%%%%%%
\begin{lemma}\label{prop:Hrow2}
Let a linear subspace $E\subset {\Bbb F}_2^n$ and space ${\Bbb F}_2^{n-s}$ be ordered as in (\ref{eq:EE'}). Let the mapping $ \sigma:{\Bbb F}_2^{n-s} \rightarrow E$ be defined in terms of indices as
$$\sigma((i)_2\cdot (\beta^{(1)},\beta^{(2)},\ldots,\beta^{(n-s)}))=(i)_2\cdot (\alpha^{(1)},\alpha^{(2)},\ldots,\alpha^{(n-s)}),$$
for fixed bases $\mathbf{B}_E$, $\mathbf{B}_{\mathbb{F}^{n-s}_2}$. Then for an arbitrary matrix $B\in GL(n-s,\mathbb{F}_2)$ and vector $t=c\cdot \mathbf{B}_{\mathbb{F}^{n-s}_2} \in \mathbb{F}^n_2, c\in {\Bbb F}_2^{n-s}$, there exists a matrix $D\in GL(n,\mathbb{F}_2)$ and vector $\gamma\in E$ such that
 $$\sigma\left(\left((i)_2\cdot (\beta^{(1)},\beta^{(2)},\ldots,\beta^{(n-s)})\right)B+t\right)=\left((i)_2\cdot (\alpha^{(1)},\alpha^{(2)},\ldots,\alpha^{(n-s)})\right)D+\gamma.$$
 Furthermore, $\left((i)_2\cdot (\alpha^{(1)},\alpha^{(2)},\ldots,\alpha^{(n-s)})\right)D+\gamma \in E$ for $i\in[0,2^{n-s}-1]$.
\end{lemma}
\begin{proof}
%Let $E\subset {\Bbb F}_2^n$ be a  subspace of dimension $n-s$. For ease of description, we provide an order of $E$. Let $ (\alpha^{(1)},\alpha^{(2)},\ldots,\alpha^{(n-s)})$ be a basis of $E$. Set the ordered set $E=\{0_n,(1)_2\cdot(\alpha^{(1)},\alpha^{(2)},\ldots,\alpha^{(n-s)}),(2)_2\cdot(\alpha^{(1)},\alpha^{(2)},\ldots,\alpha^{(n-s)}),\ldots, (2^{n-s}-1)_2\cdot (\alpha^{(1)},\alpha^{(2)},\ldots,\alpha^{(n-s)})\}=\{(i)_2\cdot (\alpha^{(1)},\alpha^{(2)},\ldots,\alpha^{(n-s)}): i=0,1,2,\ldots,2^{n-s}-1\}$, where $(i)_2$ denotes the binary representation of $i$.
% Let $ {\Bbb F}_2^{n-s}=\{(i)_2\cdot (\beta^{(1)},\beta^{(2)},\ldots,\beta^{(n-s)}): i=0,1,2,\ldots,2^{n-s}-1\}$, where $(\beta^{(1)},\beta^{(2)},\ldots,\beta^{(n-s)})$ is a basis of $ {\Bbb F}_2^{n-s}$.  Thus, we can construct one to one map $ \phi$ such that
%$\phi((i)_2\cdot (\beta^{(1)},\beta^{(2)},\ldots,\beta^{(n-s)}))=(i)_2\cdot (\alpha^{(1)},\alpha^{(2)},\ldots,\alpha^{(n-s)})$.
%Let $B\in GL(n-s,\mathbb{F}_2)$, $t\in \mathbb{F}^{n-s}_2$. For $(i)_2\cdot (\beta^{(1)},\beta^{(2)},\ldots,\beta^{(n-s)})B+t$, there must exist $A\in GL(n,\mathbb{F}_2)$ such that $\phi\left(\left((i)_2\cdot (\beta^{(1)},\beta^{(2)},\ldots,\beta^{(n-s)})\right)B+t\right)=\left((i)_2\cdot (\alpha^{(1)},\alpha^{(2)},\ldots,\alpha^{(n-s)})\right)A+\gamma$, where $\gamma\in E$.
If there exist integers $i_{1},i_{2},\ldots, i_{{n-s}}$ such that the vectors
$$(i_{1})_2\cdot \mathbf{B}_{{\Bbb F}_2^{n-s}}, (i_{2})_2\cdot \mathbf{B}_{{\Bbb F}_2^{n-s}},\ldots, (i_{{n-s}})_2\cdot \mathbf{B}_{{\Bbb F}_2^{n-s}}$$
constitute a basis of ${\Bbb F}_2^{n-s}$, then $\left((i_{1})_2,(i_{2})_2, \ldots, (i_{{n-s}})_2 \right)$ is a basis of ${\Bbb F}_2^{n-s}$. Furthermore, the vectors
$$(i_{1})_2\cdot \mathbf{B}_E, (i_{2})_2\cdot  \mathbf{B}_E,\ldots, (i_{{n-s}})_2\cdot  \mathbf{B}_E$$ also constitute a basis of $E$, since the pair of basis $\mathbf{B}_{{\Bbb F}_2^{n-s}}$ and $\mathbf{B}_E$ are related via the same linear combinations of basis vectors determined by $(i)_2$, where $i\in[0,2^{n-s}-1]$.

Without loss of generality, let the vectors $\beta^{(j)}$ and $\beta^{(i_{j})}$  be related as
 $$ \beta^{(j)}B+t= \beta^{(i_{j})}\in \mathbb{F}^{n-s}_2,\;\;\;j\in[1,n-s],$$
where $B\in GL(n-s,\mathbb{F}_2)$ and $t=c\cdot \mathbf{B}_{{\Bbb F}_2^{n-s}}$ (for some vector $c\in {\Bbb F}_2^{n-s}$). Clearly, since  $(\beta^{(i_1)},\beta^{(i_2)},\ldots,\beta^{(i_{n-s)}})$ is a basis of ${\Bbb F}_2^{n-s}$, then $(\alpha^{(i_1)},\alpha^{(i_2)},\ldots,\alpha^{(i_{n-s)}})$ is a basis of $E$.

Denoting by $\widetilde{\mathbf{B}}_E=(\alpha^{(i_1)},\alpha^{(i_2)},\ldots,\alpha^{(i_{n-s)}})$, there exists a matrix $D\in GL(n,\mathbb{F}_2)$ such that
$$\widetilde{\mathbf{B}}_E=\mathbf{B}_ED+c\cdot\mathbf{B}_E.$$
Considering the basis ${\mathbf{B}}_{{\Bbb F}_2^n}=(\alpha^{(1)},\ldots,\alpha^{({n-s)}}, \alpha^{({n-s+1)}},\ldots,\alpha^{({n)}})$, and since we view $\alpha^{(i_1)},\ldots,\alpha^{(i_{n-s})}$ as linear combinations of vectors $\alpha^{(1)},\ldots,\alpha^{(n-s)}$, then $\widetilde{\mathbf{B}}_{{\Bbb F}_2^n}=(\alpha^{(i_1)},\ldots,\alpha^{(i_{n-s)}}, \alpha^{({n-s+1)}},\ldots,\alpha^{({n)}})$ is a basis of ${\Bbb F}_2^n$ as well.
Denoting by
$ D=\left(
    \begin{array}{cccc}
      d_{11} &   d_{12} & \ldots &   d_{1n} \\
      d_{21} &   d_{22} & \ldots &   d_{2n} \\
     \vdots &   \vdots & \ddots &   \vdots \\
     d_{n1} &  d_{n2} & \ldots &   d_{nn} \\
    \end{array}
  \right),
$
and $\alpha^{(j)}=(\alpha^{(j)}_1,\alpha^{(j)}_2,\ldots,\alpha^{(j)}_n)$, the system
$$\left(
    \begin{array}{c}
      \alpha^{(i_1)} \\
      \vdots \\
      \alpha^{(i_{n-s)}}\\
       \alpha^{({n-s+1)}}\\
       \vdots\\
      \alpha^{({n)}}  \\
    \end{array}
  \right)=\left(
    \begin{array}{c}
      \alpha^{(1)} \\
      \vdots \\
      \alpha^{({n-s)}}\\
       \alpha^{({n-s+1)}}\\
       \vdots\\
      \alpha^{({n)}}  \\
    \end{array}
  \right)D+c'\cdot \mathbf{B}_{{\Bbb F}_2^n}
 $$
 means that
 $$
 \left\{
   \begin{array}{l}
     \alpha^{(1)}_1d_{1k}+\alpha^{(1)}_2d_{2k}+\ldots+\alpha^{(1)}_nd_{nk}+\gamma_1=\alpha^{(i_1)}_k, \\
    \alpha^{(2)}_1d_{1k}+\alpha^{(2)}_2d_{2k}+\ldots+\alpha^{(2)}_nd_{nk}+\gamma_2=\alpha^{(i_2)}_k, \;\;\;\;\;k\in[1,n],\\
\vdots\\
 \alpha^{(n)}_1d_{1k}+\alpha^{(n)}_2d_{2k}+\ldots+\alpha^{(n)}_nd_{nk}+\gamma_n=\alpha^{(i_n)}_k,
   \end{array}
 \right.
 $$
where for $c'=(c,0,\ldots,0)\in {\Bbb F}_2^n$ the vector $\gamma$ is given as $\gamma=c'\cdot{\mathbf{B}}_{{\Bbb F}_2^n}=c\cdot\mathbf{B}_E\in E$. This means that the matrix $D$ can be computed.
%
%Thus, for $k=1,2,\ldots,n$, we can obtain $d_{1k}, d_{2k},\ldots, d_{nk}$. Hence, we can obtain a matrix $D$.
Furthermore, denoting by $\widetilde{\mathbf{B}}_{\mathbb{F}^{n-s}_2}=(\beta^{(i_1)},\beta^{(i_2)},\ldots,\beta^{(i_{n-s})})$ we  have that
$$\sigma((i)_2\cdot \widetilde{\mathbf{B}}_{\mathbb{F}^{n-s}_2})=\sigma
((i)_2\cdot \mathbf{B}_{\mathbb{F}^{n-s}_2}B+t)=(i)_2\cdot\mathbf{B}_ED+\gamma,$$
holds for all $i\in[0,2^{n-s}-1]$. Since $\mathbf{B}_ED$ is also a basis of $E$ and $\gamma \in E$, we have $(i)_2\cdot\mathbf{B}_ED+\gamma\in E$
 which completes the proof.\qed
\end{proof}
%The following result characterizes the EA-equivalence between trivial plateaued functions.
\begin{theo}%\label{th:EA1}
Two trivial $s$-plateaued functions $f,h:\mathbb{F}^n_2\rightarrow \mathbb{F}_2$ are EA-equivalent if and only if their duals $\overline{f}^*,\overline{h}^*:\mathbb{F}^{n-s}_2\rightarrow\mathbb{F}_2$ %(defined as functions on $\mathbb{F}^{n-s}_2$ via (\ref{eq:equiv22}) by setting $z_i\leftrightsquigarrow \omega_i\leftrightsquigarrow  x_i$, $x_i\in \mathbb{F}^{n-s}_2$)
are EA-equivalent bent functions. %, i.e., \textcolor[rgb]{0.98,0.00,0.00}{in terms of (\ref{eq:equiv22})} there exists a function $\ell\in \mathcal{A}_{n-s}$ such that the equality $f^*= h^*+\ell$ holds.
\end{theo}
\begin{proof} $(\Rightarrow)$ By \cite[Theorem 3.1-$(ii)$]{Secondary} and relations (\ref{eq:identif})-(\ref{eq:equiv22}) it is clear that $\overline{f}^*,\overline{h}^*:\mathbb{F}^{n-s}_2\rightarrow\mathbb{F}_2$  are EA-equivalent bent functions.

$(\Leftarrow)$ In this part, we assume that $f$ and $h$ are two fixed trivial plateaued functions and thus the values of $h^*$ (on $S_h$) and $f^*$ (on $S_f$) are fixed in advance. In this context, we need to specify their values on $\mathbb{F}^{n-s}_2$.

 Suppose that $S_f$ and $S_h$ are given as $S_f=v+E$ and $S_h=v'+E'$ respectively,   where  $E$ and $E'$ are defined as (\ref{eq:EE'}). Without loss of generality, denote  $E= \{e_0,\ldots,e_{2^{n-s}-1}\}$ and $E= \{e'_0,\ldots,e'_{2^{n-s}-1}\}$, that is $e_i= (i)_2\cdot \mathbf{B}_E $  and $e'_i= (i)_2\cdot \mathbf{B}_{E'} $.
% Note that \cite[Lemma 3.1-$(i)$]{Secondary} implies the fact that $E$ and $E'$ are ordered as in (\ref{eq:EE'}).
 With respect to $S_f=v+E$, the dual $\overline{f}^*:\mathbb{F}^{n-s}_2\rightarrow \mathbb{F}_2$ is defined as
\begin{equation}\label{equ add1}\overline{f}^*(x_i)=f^*(v+e_i)=f^*(\omega_i),\;\;\;x_i\in \mathbb{F}^{n-s}_2,\;e_i\in E.\end{equation}
 With respect to $S_h=v'+E'$, the dual $\overline{h}^*:\mathbb{F}^{n-s}_2\rightarrow \mathbb{F}_2$ is defined as
\begin{equation}\label{equ add2}\overline{h}^*(x_i)=h^*(v'+e'_i)=h^*(z_i),\;\;\;x_i\in \mathbb{F}^{n-s}_2,\;e'_i\in E'.\end{equation}
Due to Proposition \ref{prop:Hrow} and Theorem \ref{theo:plateH}-$(ii)$, it is not difficult to see that both duals $\overline{h}^*$ and $\overline{f}^*$ are bent functions.


Now, let us assume that $\overline{f}^*,\overline{h}^*$  are affine equivalent. That is, there exist a matrix $B\in GL(n-s,\mathbb{F}_2)$ ant two vectors $t,r\in \mathbb{F}^{n-s}_2$ and $\kappa\in \mathbb{F}_2$ such that  $$\overline{h}^*(x_i)=\overline{f}^*(x_iB+t)+r\cdot x_i+\kappa, \;\;\;\;\; x_i\in \mathbb{F}^{n-s}_2.$$
By Proposition \ref{prop:Hrow}, there  exists  $c\in \mathbb{F}^n_2$ such that $c\cdot \omega_i=r\cdot x_i$ holds for all $i\in[0,2^{n-s}-1]$, i.e.,
$$(c\cdot \omega_0,\ldots,c\cdot \omega_{2^{n-s}-1})=(r\cdot x_0,\ldots,r\cdot x_{2^{n-s}-1})$$
is truth table of a linear function in $n-s$ variables ($S_f$ is an affine subspace). Furthermore, by Lemma \ref{prop:Hrow2}, (\ref{equ add1}) and  (\ref{equ add2}),  we have
  \begin{eqnarray}\label{equ add3}
  {h}^*(z_i)&=& \overline{h}^*(x_i)=\overline{f}^*(x_iB+t)+r\cdot x_i+\kappa=f^*(v+(e_iD+\gamma ))+c\cdot \omega_i+\kappa,
 % &=&f^*(\omega_iD+vD+v+\gamma )+c\cdot \omega_i+\kappa,
  \end{eqnarray}
where $D$ and $ \gamma$ can be obtained by $B$ and $t$. Note that by Lemma \ref{prop:Hrow2} we have that $e_iD+\gamma\in E$ and  consequently $v+(e_iD+\gamma)\in S_f$.

Now setting ${\mathbf{B}}_{{\Bbb F}_2^n}=(\alpha^{(1)},\ldots,\alpha^{({n-s)}}, \alpha^{({n-s+1)}},\ldots,\alpha^{({n)}})$ and $\widetilde{\mathbf{B}}_{{\Bbb F}_2^n}=(\alpha'^{(1)},\ldots,\alpha'^{({n-s)}},$ $ \alpha'^{({n-s+1)}},\ldots,\alpha'^{({n)}})$, there exists a matrix $A\in GL(n,\mathbb{F}_2)$
such that
$$ \left(
     \begin{array}{c}
       \alpha^{(1)}\\
       \vdots \\
       \alpha^{({n)}} \\
     \end{array}
   \right)=  \left(
     \begin{array}{c}
       \alpha'^{(1)}\\
       \vdots \\
       \alpha'^{({n)}} \\
     \end{array}
   \right)A,
$$
holds, i.e., we have $\alpha^{(i)}=\alpha'^{(i)}A$, $i=1,2,\ldots, n$.  Furthermore, we have $e_i=e'_i A$ for $i=0,1,2\ldots,2^{n-s}-1$. Now, for $x\in \mathbb{F}^n_2$, by (\ref{WHT}), (\ref{equ add1}), (\ref{equ add2}) and  (\ref{equ add3}) we have
 \begin{equation}\label{equ IWHT}
 \begin{array}{rl}
 (-1)^{h(x)}=&2^{-n}\sum\limits_{e'_i\in E'}W_h(e'_i+v')(-1)^{e'_i\cdot x+v'\cdot x}\\
=&2^{\frac{s-n}{2}}\sum\limits_{e'_i\in E'}(-1)^{h^*(e'_i+v')+e'_i\cdot x+v'\cdot x}\\
=&2^{\frac{s-n}{2}}\sum\limits_{e_i\in E}(-1)^{f^*(v+e_iD+\gamma)+e_iA^{-1}\cdot x+v'\cdot x}\\
=&2^{\frac{s-n}{2}}\sum\limits_{e_i\in E}(-1)^{f^*(\omega_iD+vD+v+\gamma )+c\cdot \omega_i+\kappa+e_iA^{-1}\cdot x+v'\cdot x}\\
=&2^{-n}\sum\limits_{\omega_i\in S_f}W_f(\omega_iD+vD+v+\gamma)(-1)^{e_iA^{-1}\cdot x+v'\cdot x+c\cdot \omega_i+\kappa}\\
\stackrel{t_i=\omega_iD+vD+v+\gamma}{=} &2^{-n}\sum\limits_{t_i\in S_f}W_f(t_i)(-1)^{[(t_i+\gamma+v)D^{-1}A^{-1}+v']\cdot x+c\cdot (t_i+vD+\gamma+v)D^{-1}+\kappa}\\
{=} &2^{-n}\sum\limits_{t_i\in S_f}W_f(t_i)(-1)^{t_i\cdot [x(D^{-1}A^{-1})^T+cD^{-T}]+[(\gamma+v)D^{-1}A^{-1}+v']\cdot x+[c\cdot (vD+\gamma+v)D^{-1}+\kappa]}\\
{=}& (-1)^{f(x(D^{-1}A^{-1})^T+cD^{-T})+[(\gamma+v)D^{-1}A^{-1}+v']\cdot x+[c\cdot (vD+\gamma+v)D^{-1}+\kappa]}.\\
 \end{array}
 \end{equation}
Hence $f$ and $h$ are EA-equivalent, and the proof is completed.\qed
\end{proof}


\end{document}
